% Leçon 3 — Expression orale : amour de la famille et la solidarité familiale — Chinois 1ère A — Cours signé OnBuch+
% Programme officiel MINESEC. Compiler avec : tectonic ZHA03-oral-amour-de-la-famille-et-la-solidarite-familial.tex
\newcommand{\DOCMATIERE}{Chinois}
\newcommand{\DOCNIVEAU}{1ère A}
\newcommand{\DOCMODULE}{Module 1 : Vie familiale et intégration sociale}
\newcommand{\DOCLECON}{ZHA03}
\newcommand{\DOCTITRE}{Leçon 3 — Expression orale : amour de la famille et la solidarité familiale}
% OnBuch+ — préambule « cours signés » ENRICHI (compilé avec tectonic / XeLaTeX)
% Système visuel « Pop » d'OnBuch+ : fond crème (jamais blanc), bordures encre
% épaisses, ombres « dures » décalées, titres Archivo Black en capitales.
% Version MATHÉMATIQUES : graphes (pgfplots/tikz), boîtes pédagogiques
% (définition, propriété, méthode, attention, le savais-tu, exercice résolu,
% exercice, TP, bilan), page de garde et signature OnBuch+ ancrée en bas.
%
% Macros attendues dans main.tex AVANT \input{preamble} :
%   \DOCMATIERE  \DOCNIVEAU  \DOCMODULE  \DOCLECON (numéro)  \DOCTITRE
\documentclass[11pt]{article}

\usepackage{fontspec}
\usepackage[french]{babel} % français = langue principale ; le chinois via \zh{..} / textebox (xeCJK)
\usepackage{xeCJK}
\usepackage{pifont} % \ding{51} \ding{55} utilisés par les modèles
\usepackage[a4paper,margin=2cm,top=3.4cm,bottom=2.8cm,headheight=1.4cm,footskip=1.3cm]{geometry}
\usepackage{amsmath,amssymb,amsfonts}
\usepackage{mathtools} % \xmapsto, \coloneqq... souvent utilisés par les modèles
\usepackage{cancel} % simplifications barrées (bilans de forces)
% \abs{z} (module, valeur absolue) et \norm{u} (norme), utilisés par les modèles
\providecommand{\abs}{}\let\abs\relax\DeclarePairedDelimiter{\abs}{\lvert}{\rvert}
\providecommand{\norm}{}\let\norm\relax\DeclarePairedDelimiter{\norm}{\lVert}{\rVert}
\usepackage[table]{xcolor} % [table] : fournit aussi \rowcolors (alternance de lignes)
\usepackage{pagecolor}
\usepackage{tikz}
\usetikzlibrary{arrows.meta,positioning,calc,shapes.geometric,shapes.misc,shapes.arrows,decorations.pathmorphing,decorations.pathreplacing,decorations.markings,patterns,shadows,mindmap,trees,fit,backgrounds,matrix,intersections,angles,quotes}
\usepackage{pgfplots}
\usepackage[version=4]{mhchem} % équations nucléaires \ce{^{235}_{92}U}
\usepackage[european,siunitx]{circuitikz} % schémas électriques
\pgfplotsset{compat=1.18}
\usepgfplotslibrary{fillbetween,groupplots} % aires entre courbes (calcul intégral)
\usepackage{enumitem}
\usepackage{fancyhdr}
% [most] charge toutes les bibliothèques tcolorbox (listings, minted, raster,
% documentation...) et gonfle inutilement la mémoire TeX sur un document long
% (~300 boîtes) : on ne charge que ce qui est réellement utilisé.
\usepackage[skins,breakable]{tcolorbox}
\usepackage{graphicx}
\usepackage{colortbl}
\usepackage{booktabs}
\usepackage{tabularx}
\usepackage{diagbox} % case d'en-tête diagonale des tableaux à double entrée
% Colonnes extensibles centrée / gauche / droite (C, L, R), souvent utilisées par les modèles
\newcolumntype{C}{>{\centering\arraybackslash}X}
\newcolumntype{L}{>{\raggedright\arraybackslash}X}
\newcolumntype{R}{>{\raggedleft\arraybackslash}X}
\usepackage{array}
\usepackage{multicol}
\usepackage{multirow}
\usepackage{siunitx}
% parse-numbers=false : accepte aussi \SI{2,5 \times 10^{-3}}{...}
\sisetup{locale=FR,output-decimal-marker={,},group-minimum-digits=4,parse-numbers=false}
\DeclareSIUnit\ohm{\ensuremath{\symup{\Omega}}} % U+2126 absent de la police
\DeclareSIUnit\division{div} % réglages d'oscilloscope (ms/div, V/div)
\DeclareSIUnit\tr{tr} % tours (vitesse de rotation, ex. \SI{1500}{\tr\per\minute})
\DeclareSIUnit\molar{M} % concentration molaire (ex. \SI{3}{\molar}, \SI{1}{\micro\molar})
\DeclareSIUnit\an{an} % année (le modèle confond parfois avec \year, un compteur TeX) — ex. \SI{5}{\kilo\meter\per\an}
\DeclareSIUnit\FCFA{FCFA} % franc CFA (ex. \SI{500}{\FCFA\per\kilo\gram})
\DeclareSIUnit\degreeBrix{\textdegree Brix} % teneur en sucre d'un sirop/jus (réfractométrie)
\DeclareSIUnit\month{mois} % le modèle utilise parfois \month, un compteur TeX, comme unité de durée
\usepackage{qrcode}
% \degree : commande fréquemment utilisée par les modèles pour un angle ou une
% température en dehors de \SI{}{} (non fournie par les paquets chargés).
\providecommand{\degree}{^{\circ}}
% \qed (fin de démonstration), utilisé par les modèles sans amsthm
\providecommand{\qed}{\hfill$\square$}
% \barre : double barre des valeurs interdites dans un tableau de variations (tkz-tab)
\providecommand{\barre}{\ensuremath{\|}}
% environnement proof (amsthm non chargé) utilisé par les modèles
\ifdefined\proof\else\newenvironment{proof}[1][Démonstration]{\par\noindent\textit{#1.}\ }{\qed\par}\fi
\usepackage{lastpage}
\usepackage{hyperref}
\hypersetup{hidelinks}

% ============================================================
% Palette « Pop » OnBuch+ — mêmes hex que la classe P (plus_ui.dart)
% ============================================================
\definecolor{poporange}{HTML}{F59321}
\definecolor{popdark}{HTML}{E07A0C}
\definecolor{popcream}{HTML}{FFF6E8}
\definecolor{popink}{HTML}{1C1714}
\definecolor{poppurple}{HTML}{7A5AE0}
\definecolor{popgreen}{HTML}{2E9E6A}
\definecolor{poppink}{HTML}{E0457B}
\definecolor{popblue}{HTML}{2F7DE1}
\definecolor{popgold}{HTML}{C98A12}
\definecolor{popmuted}{HTML}{8A7F76}
\definecolor{popline}{HTML}{EADFD2}
\definecolor{popgray}{HTML}{8A7F76}
\colorlet{popgrey}{popgray}
% Alias tolérants (le modèle nomme parfois les couleurs différemment)
\colorlet{popwhite}{white}
\colorlet{popblack}{popink}
\colorlet{popred}{poppink}
\colorlet{popyellow}{popgold}
% Teintes claires (fonds de tableaux/encadrés) — le modèle nomme parfois
% les couleurs avec un suffixe L (ex. popblueL) sans les avoir définies.
\colorlet{poporangeL}{poporange!20}
\colorlet{popdarkL}{popdark!20}
\colorlet{poppurpleL}{poppurple!20}
\colorlet{popgreenL}{popgreen!20}
\colorlet{poppinkL}{poppink!20}
\colorlet{popblueL}{popblue!20}
\colorlet{popgoldL}{popgold!20}
\colorlet{popredL}{popred!20}
\colorlet{popgrayL}{popgray!20}
% Teintes claires pour fonds de tableaux / zones
\colorlet{popblueL}{popblue!10!white}
\colorlet{poporangeL}{poporange!14!white}
\colorlet{popgreenL}{popgreen!12!white}
\colorlet{poppurpleL}{poppurple!12!white}
\colorlet{poppinkL}{poppink!12!white}
\colorlet{popgoldL}{popgold!14!white}

\pagecolor{popcream}

% ============================================================
% Typographie
% ============================================================
\defaultfontfeatures{Ligatures=TeX}
\setmainfont{PlusJakartaSans-Medium.ttf}[Path=../../../fonts/,BoldFont=PlusJakartaSans-Bold.ttf]
% Symboles, API et emojis absents de la police principale : repli sur DejaVu Sans (caractères actifs)
\newfontface\symfont{DejaVuSans.ttf}[Path=../../../fonts/]
\makeatletter
\newcommand\symactive[1]{\catcode#1=13 \begingroup\lccode`\~=#1 \lowercase{\endgroup\def~}{{\symfont\char#1}}}
\newcommand\symblank[1]{\catcode#1=13 \begingroup\lccode`\~=#1 \lowercase{\endgroup\def~}{}}
\newcommand\symhyph[1]{\catcode#1=13 \begingroup\lccode`\~=#1 \lowercase{\endgroup\def~}{\mbox{-}}}
\newcount\symc
\newcommand\symrange[2]{\symc=#1 \@whilenum\symc<#2 \do{\expandafter\symactive\expandafter{\the\symc}\advance\symc by 1 }}
\newcommand\symblankrange[2]{\symc=#1 \@whilenum\symc<#2 \do{\expandafter\symblank\expandafter{\the\symc}\advance\symc by 1 }}
\makeatother
\symrange{"0250}{"02FF}\symactive{"2713}\symactive{"2714}\symactive{"2717}\symactive{"2718}\symactive{"2610}\symactive{"2611}\symactive{"2612}
\symrange{"2600}{"27BF}
\catcode"2705=13 \begingroup\lccode`\~="2705 \lowercase{\endgroup\def~}{{\symfont\char"2713}}\symblank{"26BD}\symblank{"26A1}
\symrange{"2460}{"257F}\symactive{"223C}\symactive{"2248}\symactive{"2260}
\symhyph{"2011}\symhyph{"2E1A}\symblankrange{"1F300}{"1FAFF}\symblank{"FE0F}
% Caractères chinois : WenQuanYi Zen Hei (sous-ensemble GB2312 + pinyin), gras/italique simulés
\setCJKmainfont{WenQuanYiZenHei-subset.ttf}[Path=../../../fonts/,AutoFakeBold=3,AutoFakeSlant=0.2]
\xeCJKsetup{CJKspace=false}
% Pinyin : voyelles à caron (ǎ ǐ ǒ ǔ ǖ ǘ ǚ ǜ) absentes de la police principale -> caractères actifs
\newfontface\pyfont{WenQuanYiZenHei-subset.ttf}[Path=../../../fonts/]
\newcommand\pyactive[1]{\catcode#1=13 \begingroup\lccode`\~=#1 \lowercase{\endgroup\def~}{{\pyfont\char#1}}}
\pyactive{"01CD}\pyactive{"01CE}\pyactive{"01CF}\pyactive{"01D0}\pyactive{"01D1}\pyactive{"01D2}\pyactive{"01D3}\pyactive{"01D4}\pyactive{"01D5}\pyactive{"01D6}\pyactive{"01D7}\pyactive{"01D8}\pyactive{"01D9}\pyactive{"01DA}\pyactive{"01DB}\pyactive{"01DC}
% Maths en sans-serif (Fira Math) avec chiffres et lettres droites de Plus
% Jakarta, pour que les formules se fondent dans le texte.
\usepackage[math-style=ISO,bold-style=ISO]{unicode-math}
\setmathfont{FiraMath-Regular.otf}[Path=../../../fonts/]
\setmathfont{PlusJakartaSans-Medium.ttf}[Path=../../../fonts/,range={up/num,up/{Latin,latin},\mathup}]
\usepackage{icomma} % virgule décimale sans espace en mode math
% Caractères mathématiques Unicode absents des polices de texte (Plus Jakarta,
% Archivo Black) : rendus par leur équivalent mathématique, en texte comme en
% formule — sinon ils s'affichent en carré vide (ex. « Arithmétique dans ℤ »).
\usepackage{newunicodechar}
\newunicodechar{ℝ}{\ensuremath{\mathbb{R}}}
\newunicodechar{ℤ}{\ensuremath{\mathbb{Z}}}
\newunicodechar{ℂ}{\ensuremath{\mathbb{C}}}
\newunicodechar{ℕ}{\ensuremath{\mathbb{N}}}
\newunicodechar{ℚ}{\ensuremath{\mathbb{Q}}}
\newunicodechar{𝔻}{\ensuremath{\mathbb{D}}}
\newunicodechar{α}{\ensuremath{\alpha}}
\newunicodechar{β}{\ensuremath{\beta}}
\newunicodechar{γ}{\ensuremath{\gamma}}
\newunicodechar{δ}{\ensuremath{\delta}}
\newunicodechar{ε}{\ensuremath{\varepsilon}}
\newunicodechar{θ}{\ensuremath{\theta}}
\newunicodechar{λ}{\ensuremath{\lambda}}
\newunicodechar{μ}{\ensuremath{\mu}}
\newunicodechar{σ}{\ensuremath{\sigma}}
\newunicodechar{τ}{\ensuremath{\tau}}
\newunicodechar{φ}{\ensuremath{\varphi}}
\newunicodechar{ψ}{\ensuremath{\psi}}
\newunicodechar{ω}{\ensuremath{\omega}}
\newunicodechar{Σ}{\ensuremath{\Sigma}}
% majuscules produites par \MakeUppercase dans les titres (α-aminés -> Α)
\newunicodechar{Α}{\ensuremath{\alpha}}
\newunicodechar{Β}{\ensuremath{\beta}}
\newunicodechar{Γ}{\ensuremath{\Gamma}}
\newunicodechar{Δ}{\ensuremath{\Delta}}
\newunicodechar{Ε}{\ensuremath{\varepsilon}}
\newunicodechar{Θ}{\ensuremath{\Theta}}
\newunicodechar{Λ}{\ensuremath{\Lambda}}
\newunicodechar{Μ}{\ensuremath{\mu}}
\newunicodechar{Τ}{\ensuremath{\tau}}
\newunicodechar{Φ}{\ensuremath{\Phi}}
\newunicodechar{Ψ}{\ensuremath{\Psi}}
\newunicodechar{Ω}{\ensuremath{\Omega}}
\newunicodechar{↦}{\ensuremath{\mapsto}}
\newunicodechar{∈}{\ensuremath{\in}}
\newunicodechar{∉}{\ensuremath{\notin}}
\newunicodechar{∧}{\ensuremath{\wedge}}
\newunicodechar{⇔}{\ensuremath{\Leftrightarrow}}
\newunicodechar{⟺}{\ensuremath{\Longleftrightarrow}}
\newunicodechar{⇒}{\ensuremath{\Rightarrow}}
\newunicodechar{⟹}{\ensuremath{\Longrightarrow}}
\newunicodechar{∘}{\ensuremath{\circ}}
\newunicodechar{∪}{\ensuremath{\cup}}
\newunicodechar{∩}{\ensuremath{\cap}}
\newunicodechar{≡}{\ensuremath{\equiv}}
\newunicodechar{⊂}{\ensuremath{\subset}}
\newunicodechar{⊥}{\ensuremath{\perp}}
\newunicodechar{∃}{\ensuremath{\exists}}
\newunicodechar{∀}{\ensuremath{\forall}}
\newunicodechar{∛}{\ensuremath{\sqrt[3]{\,}}}
\newunicodechar{∜}{\ensuremath{\sqrt[4]{\,}}}
\newunicodechar{⃗}{\ensuremath{{}^{\to}}}
\newunicodechar{ⁿ}{\ifmmode^{n}\else\textsuperscript{\ensuremath{n}}\fi}
\newunicodechar{ˣ}{\ifmmode^{x}\else\textsuperscript{\ensuremath{x}}\fi}
\newunicodechar{ᵇ}{\ifmmode^{b}\else\textsuperscript{\ensuremath{b}}\fi}
\newunicodechar{ʳ}{\ifmmode^{r}\else\textsuperscript{\ensuremath{r}}\fi}
\newunicodechar{ᵘ}{\ifmmode^{u}\else\textsuperscript{\ensuremath{u}}\fi}
\newunicodechar{ʸ}{\ifmmode^{y}\else\textsuperscript{\ensuremath{y}}\fi}
\newunicodechar{ᵃ}{\ifmmode^{a}\else\textsuperscript{\ensuremath{a}}\fi}
\newunicodechar{ᵉ}{\ifmmode^{e}\else\textsuperscript{\ensuremath{e}}\fi}
\newunicodechar{ᵏ}{\ifmmode^{k}\else\textsuperscript{\ensuremath{k}}\fi}
\newunicodechar{ᵖ}{\ifmmode^{p}\else\textsuperscript{\ensuremath{p}}\fi}
\newunicodechar{ᴺ}{\ifmmode^{N}\else\textsuperscript{\ensuremath{N}}\fi}
\newunicodechar{ᶜ}{\ifmmode^{c}\else\textsuperscript{\ensuremath{c}}\fi}
\newunicodechar{ʰ}{\ifmmode^{h}\else\textsuperscript{\ensuremath{h}}\fi}
\newunicodechar{ᵠ}{\ifmmode^{q}\else\textsuperscript{\ensuremath{q}}\fi}
\newunicodechar{ₙ}{\ifmmode_{n}\else\textsubscript{\ensuremath{n}}\fi}
\newunicodechar{ₐ}{\ifmmode_{a}\else\textsubscript{\ensuremath{a}}\fi}
\newunicodechar{ᵢ}{\ifmmode_{i}\else\textsubscript{\ensuremath{i}}\fi}
\newunicodechar{ᵣ}{\ifmmode_{r}\else\textsubscript{\ensuremath{r}}\fi}
\newunicodechar{ₖ}{\ifmmode_{k}\else\textsubscript{\ensuremath{k}}\fi}
\newunicodechar{ᵦ}{\ifmmode_{\beta}\else\textsubscript{\ensuremath{\beta}}\fi}
\newunicodechar{ₓ}{\ifmmode_{x}\else\textsubscript{\ensuremath{x}}\fi}
\newfontfamily\popdisplay{ArchivoBlack-Regular.ttf}[Path=../../../fonts/]
\newfontfamily\poplogo{SpaceGrotesk-Bold.ttf}[Path=../../../fonts/]
\newfontfamily\popsemi{PlusJakartaSans-SemiBold.ttf}[Path=../../../fonts/]
\newfontfamily\popxbold{PlusJakartaSans-ExtraBold.ttf}[Path=../../../fonts/]
\newfontfamily\popxboldls{PlusJakartaSans-ExtraBold.ttf}[Path=../../../fonts/,LetterSpace=3.0]

\setlist[enumerate]{leftmargin=1.7em,itemsep=5pt,topsep=4pt}
\setlist[itemize]{leftmargin=1.7em,itemsep=4pt,topsep=4pt,label={\color{poporange}\textbullet}}
\setlength{\parindent}{0pt}
\setlength{\parskip}{5pt}
\renewcommand{\arraystretch}{1.25}

% pgfplots : style Pop par défaut
\pgfplotsset{
  every axis/.append style={axis line style={popink,line width=1pt},tick style={popink},
    grid style={popline,line width=0.5pt},label style={font=\small},tick label style={font=\footnotesize}},
  popcurve/.style={poporange,line width=1.8pt,no markers,smooth},
}
% Même style disponible dans un \draw TikZ simple (hors pgfplots)
\tikzset{popcurve/.style={poporange,line width=1.8pt}}
\tikzset{popgrid/.style={popline,very thin}}
\colorlet{popcurve}{poporange} % le nom du style est parfois utilisé comme couleur

% ============================================================
% Pill / squiggle
% ============================================================
\newcommand{\poppill}[3]{%
  \tikz[baseline=-0.55ex]\node[draw=popink,line width=1.2pt,rounded corners=2.6pt,%
    fill=#2,inner xsep=6.5pt,inner ysep=3pt]{%
    \popxboldls\fontsize{7}{7}\selectfont\color{#3}\MakeUppercase{#1}};%
}
\newcommand{\popsquiggle}[1]{%
  \tikz[baseline=-0.4ex]{
    \draw[#1,line width=2.6pt,line cap=round]
      plot [smooth] coordinates {(0,0.09) (0.34,0.01) (0.68,0.21) (1.02,0.01) (1.36,0.10)};
  }%
}
% Mot-clé surligné (marqueur orange sous le texte)
\newcommand{\cle}[1]{{\popxbold\color{popdark}#1}}

% ============================================================
% En-tête / pied
% ============================================================
\pagestyle{fancy}
\fancyhf{}
\renewcommand{\headrulewidth}{0pt}
\renewcommand{\footrulewidth}{0pt}
\lhead{\raisebox{-0.15ex}{\poplogo\fontsize{13}{13}\selectfont\color{popink}OnBuch\color{poporange}+}}
\rhead{\poppill{\DOCMATIERE\ \textbullet\ \DOCNIVEAU\ \textbullet\ Leçon \DOCLECON}{white}{popink}}
\lfoot{\popxbold\fontsize{7.6}{7.6}\selectfont\color{popmuted}\MakeUppercase{Cours signé OnBuch+}\ \textbullet\ {\popsemi\DOCTITRE}}
\rfoot{\tikz[baseline=-0.6ex]\node[rounded corners=8.5pt,fill=popink,minimum height=17pt,minimum width=17pt,inner xsep=5pt,inner sep=0]{\popxbold\fontsize{8}{8}\selectfont\color{popcream}\thepage\,/\,\pageref*{LastPage}};}

% ============================================================
% Titres
% ============================================================
\newcommand{\chaptitre}[2]{%
  \begin{tcolorbox}[enhanced,colback=poporange,colframe=popink,boxrule=2.6pt,arc=11pt,
      drop shadow=popink,left=15pt,right=15pt,top=12pt,bottom=11pt,before skip=3pt,after skip=18pt]
    \poppill{#2}{popink}{popcream}\par\vspace{9pt}
    {\raggedright\hyphenpenalty=10000\exhyphenpenalty=10000\popdisplay\fontsize{22}{24}\selectfont\color{popcream}\MakeUppercase{#1}\par}
  \end{tcolorbox}
}
% Section numérotée : pastille orange avec le numéro + titre Archivo
\newcounter{popsec}
\newcommand{\coursec}[1]{%
  \stepcounter{popsec}\setcounter{popsub}{0}%
  \par\vspace{16pt}\noindent
  \begin{minipage}{\linewidth}
  \tikz[baseline=-0.7ex]\node[draw=popink,line width=1.6pt,fill=poporange,rounded corners=4pt,minimum size=22pt,inner sep=0]{\popdisplay\fontsize{12}{12}\selectfont\color{popcream}\thepopsec};%
  \hspace{8pt}{\popdisplay\fontsize{15}{17}\selectfont\color{popink}\MakeUppercase{#1}}\par
  \vspace{4pt}\noindent{\color{poporange}\rule{36pt}{3.6pt}}
  \end{minipage}\par\nopagebreak\vspace{8pt}
}
\newcounter{popsub}
\newcommand{\courssub}[1]{%
  \stepcounter{popsub}%
  \par\vspace{9pt}\noindent{\popxbold\fontsize{11.5}{13}\selectfont\color{popink}{\color{poporange}\thepopsec.\thepopsub}\hspace{6pt}#1}\par\nopagebreak\vspace{4pt}
}

% ============================================================
% Boîtes pédagogiques — même recette : fond blanc, bordure épaisse colorée,
% ombre dure encre, badge plein. Titre optionnel [#1].
% ============================================================
\tcbset{popbase/.style={breakable,enhanced,colback=white,boxrule=2.3pt,arc=9pt,
  shadow={3pt}{-3pt}{0pt}{popink},left=14pt,right=14pt,top=11pt,bottom=11pt,before skip=11pt,after skip=15pt,
  fontupper=\popsemi\fontsize{10.6}{14.5}\selectfont\color{popink},
  fontlower=\popsemi\fontsize{10.6}{14.5}\selectfont\color{popink}}}
% Coche dessinée (le glyphe ✓ n'existe pas dans les polices du document)
\newcommand{\popcheck}[1]{\tikz[baseline=-0.5ex]\draw[#1,line width=1.6pt,line cap=round,line join=round](-0.13,0.0)--(-0.04,-0.09)--(0.13,0.1);}
\providecommand{\coche}{\popcheck{popgreen}}
\providecommand{\resultat}[1]{\par\smallskip\noindent\fcolorbox{popgreen}{popgreenL}{\parbox{\dimexpr\linewidth-2\fboxsep-2\fboxrule\relax}{\textbf{Résultat.} #1}}\par\smallskip}
\newcommand{\popboxhead}[3]{\poppill{#1}{#2}{white}\ifx\relax#3\relax\else\hspace{6pt}{\popxbold\fontsize{10.6}{12}\selectfont\color{popink}#3}\fi\par\vspace{7pt}}

\newtcolorbox{aretenir}[1][]{popbase,colframe=popgreen,before upper={\popboxhead{À retenir}{popgreen}{#1}}}
% Texte en chinois (document de lecture, dialogue, extrait) : cadre
% d'étude. Un titre optionnel (source, auteur) s'affiche dans l'en-tête.
\newtcolorbox{textebox}[1][]{popbase,colframe=popblue,before upper={\popboxhead{文本 · Texte}{popblue}{#1}},after upper={}}
% Marques ✓ / ✗ (forme correcte / fautive) souvent utilisées par les modèles
\providecommand{\cross}{\textcolor{poppink}{\ensuremath{\times}}}
\providecommand{\xmark}{\cross}\providecommand{\crossmark}{\cross}\providecommand{\cmark}{\ensuremath{\checkmark}}
% Ligne de réponse / justification à compléter (inventée par les modèles)
\providecommand{\justification}[1]{\par\noindent\textit{Justification~:} \dotfill\par}
% \textipa (package tipa non chargé) : la police ne couvre pas l'API -> simple passage
\providecommand{\textipa}[1]{#1}
% Soulignement (ulem non chargé) : \uline, \ul
\providecommand{\uline}[1]{\underline{#1}}\providecommand{\ul}[1]{\underline{#1}}
% \glossaire{\item ...} inventée par les modèles : liste de glossaire
\providecommand{\glossaire}[1]{\begin{itemize}#1\end{itemize}}
% \alert (beamer) inventée par les modèles : texte mis en évidence
\providecommand{\alert}[1]{\textbf{\textcolor{poppink}{#1}}}
% variantes ASCII d'ordinaux écrites en macro
\providecommand{\iere}{\textsuperscript{re}}\providecommand{\ieme}{\textsuperscript{e}}\providecommand{\ere}{\textsuperscript{re}}\providecommand{\eme}{\textsuperscript{e}}\providecommand{\er}{\textsuperscript{er}}\providecommand{\ie}{\textsuperscript{e}}
% Ordinaux français écrits en macro par les modèles : 1\ière, 3\ième
\newcommand{\ière}{\textsuperscript{re}}\newcommand{\ième}{\textsuperscript{e}}
% \exemple (inventée par les modèles) : étiquette « Exemple : »
\providecommand{\exemple}{\textbf{Exemple~:}\ }
% \square utilisable aussi hors mode math (cases à cocher)
\AtBeginDocument{\let\squareMath\square\renewcommand{\square}{\ensuremath{\squareMath}}}
% Mot ou phrase en chinois à l'intérieur d'un paragraphe français
\newcommand{\zh}[1]{\ifmmode\text{#1}\else{#1}\fi}
\let\eng\zh \let\esp\zh \let\all\zh % alias tolérés
% flèches d'intonation inventées par les modèles
\providecommand{\textsearrow}{}\renewcommand{\textsearrow}{\ensuremath{\searrow}}\providecommand{\textnearrow}{}\renewcommand{\textnearrow}{\ensuremath{\nearrow}}
% \fr{..} : mot français (inventé par les modèles)
\providecommand{\fr}[1]{\textit{#1}}
% zhbox : encadré de texte chinois inventé par les modèles
\newenvironment{zhbox}{\begin{quote}}{\end{quote}}
% \vs : « versus » inventé par les modèles
\providecommand{\vs}{\textit{vs}\ }
% \blank : case à compléter inventée par les modèles
\providecommand{\blank}[1][]{\underline{\hspace{1.6em}}}
\newtcolorbox{exemplebox}[1][]{popbase,colframe=poporange,before upper={\popboxhead{Exemple}{poporange}{#1}}}
\newtcolorbox{definition}[1][]{popbase,colframe=popblue,before upper={\popboxhead{Définition}{popblue}{#1}}}
\newtcolorbox{propriete}[1][]{popbase,colframe=poppurple,before upper={\popboxhead{Propriété}{poppurple}{#1}}}
\newtcolorbox{methode}[1][]{popbase,colframe=popgold,before upper={\popboxhead{Méthode}{popgold}{#1}}}
\newtcolorbox{attention}[1][]{popbase,colframe=poppink,before upper={\popboxhead{Attention}{poppink}{#1}}}
\newtcolorbox{experience}[1][]{popbase,colframe=popblue,colback=popblueL,before upper={\popboxhead{Expérience}{popblue}{#1}}}
% « Le savais-tu ? » : carte violette pleine (contexte, culture, Cameroun)
\newtcolorbox{savaistu}[1][]{popbase,colback=poppurple,colframe=popink,
  fontupper=\popsemi\fontsize{10.4}{14.2}\selectfont\color{white},
  before upper={\poppill{Le savais-tu\,?}{white}{poppurple}\ifx\relax#1\relax\else\hspace{6pt}{\popxbold\color{white}#1}\fi\par\vspace{7pt}}}
% Exercice résolu : énoncé en haut, \tcblower puis la solution sur fond crème
\newcounter{popexo}\newcounter{popexr}
\newtcolorbox{exoresolu}[1][]{popbase,colframe=poporange,colbacklower=popcream,
  segmentation style={popink,line width=1.2pt,dashed},
  before upper={\stepcounter{popexr}\popboxhead{Exercice résolu \thepopexr}{poporange}{#1}},
  before lower={\poppill{Solution}{popink}{popcream}\par\vspace{6pt}}}
% Exercice (non résolu en place) — la correction est dans la partie Corrigés
\newtcolorbox{exercice}[1][]{popbase,colframe=popink,
  before upper={\stepcounter{popexo}\popboxhead{Exercice \thepopexo}{popink}{#1}}}
% Corrigé
\newtcolorbox{corrige}[1][]{popbase,colframe=popgreen,colback=popgreenL,
  before upper={\popboxhead{Corrigé}{popgreen}{#1}}}
% Objectifs / prérequis / bilan
\newtcolorbox{objectifs}{popbase,colframe=popink,colback=poporangeL,
  before upper={\popboxhead{Objectifs de la leçon}{popink}{}}}
\newtcolorbox{prerequis}{popbase,colframe=popink,colback=popblueL,
  before upper={\popboxhead{Prérequis}{popblue}{}}}
\newtcolorbox{bilan}{popbase,colframe=popink,colback=white,boxrule=2.8pt,
  before upper={\popboxhead{Fiche bilan}{poporange}{L'essentiel à connaître pour l'examen}}}
% Figure encadrée avec légende
\newtcolorbox{popfigure}[1][]{enhanced,breakable=false,colback=white,colframe=popline,boxrule=1.4pt,arc=8pt,
  left=10pt,right=10pt,top=10pt,bottom=8pt,before skip=10pt,after skip=12pt,halign=center,
  lower separated=false,halign lower=center,
  fontlower=\popsemi\fontsize{9}{11.5}\selectfont\color{popmuted},
  #1}
\newcommand{\legende}[1]{\par\vspace{6pt}{\popsemi\fontsize{9}{11.5}\selectfont\color{popmuted}#1}}

% ============================================================
% Signature OnBuch+
% ============================================================
\newcommand{\poplogoinline}[1]{{\poplogo\fontsize{#1}{#1}\selectfont\color{popink}OnBuch\color{poporange}+}}
\newcommand{\signatureonbuch}{%
  \begin{tcolorbox}[enhanced,colback=popink,colframe=popink,boxrule=2.2pt,arc=10pt,
      drop shadow=poporange,left=14pt,right=14pt,top=10pt,bottom=10pt,before skip=0pt,after skip=0pt]
    \tikz[baseline=-0.7ex]\node[circle,fill=popgreen,draw=popcream,line width=1.3pt,minimum size=22pt,inner sep=0]{\popcheck{white}};%
    \hspace{9pt}\begin{minipage}[c]{0.62\linewidth}
      {\popxbold\fontsize{12}{13}\selectfont\color{popcream}Signé\ \textbullet\ {\poplogo OnBuch\color{poporange}+}}\par\vspace{2pt}
      {\raggedright\popsemi\fontsize{8}{10.5}\selectfont\color{popline}\MakeUppercase{Équipe pédagogique OnBuch+}\\\MakeUppercase{Conforme au programme officiel MINESEC}\par}
    \end{minipage}\hfill
    {\poplogo\fontsize{22}{22}\selectfont\color{popcream}OnBuch\color{poporange}+}
  \end{tcolorbox}%
}

% ============================================================
% Page de garde
% #1 titre · #2 matière · #3 niveau · #4 module · #5 n° leçon · #6 durée ·
% #7 description / objectifs (texte ou itemize)
% ============================================================
\newcommand{\pagegarde}[7]{%
  \thispagestyle{empty}
  \newgeometry{a4paper,margin=1.7cm,top=1.6cm,bottom=1.4cm}%
  \begin{tikzpicture}[remember picture,overlay]
    \fill[poporange!18] ([xshift=-3.2cm,yshift=-3.6cm]current page.north east) circle (4.6cm);
    \fill[poppurple!14] ([xshift=2.4cm,yshift=7.6cm]current page.south west) circle (3.8cm);
  \end{tikzpicture}%
  {\poplogo\fontsize{30}{30}\selectfont\color{popink}OnBuch\color{poporange}+}\hfill
  \raisebox{0.6ex}{\poppill{Cours signé \textbullet\ Premium}{popink}{popcream}}\par
  \vspace{1pt}\popsquiggle{poppurple}\par
  \vspace{8pt}
  \begin{tcolorbox}[enhanced,colback=poporange,colframe=popink,boxrule=2.8pt,arc=13pt,
      drop shadow=popink,left=18pt,right=18pt,top=12pt,bottom=13pt,after skip=10pt]
    \poppill{#2 \textbullet\ #3}{popink}{popcream}\hspace{5pt}\poppill{#4}{white}{popink}\par\vspace{8pt}
    {\popdisplay\fontsize{14}{14}\selectfont\color{popink}LEÇON #5}\par\vspace{4pt}
    {\raggedright\hyphenpenalty=10000\exhyphenpenalty=10000\popdisplay\fontsize{25}{27}\selectfont\color{popcream}\MakeUppercase{#1}\par}
    \vspace{8pt}
    {\popxbold\fontsize{9.5}{9.5}\selectfont\color{popink}\MakeUppercase{Durée conseillée : #6}}
  \end{tcolorbox}
  \begin{tcolorbox}[enhanced,colback=white,colframe=popink,boxrule=2.2pt,arc=10pt,
      drop shadow=popink,left=16pt,right=16pt,top=10pt,bottom=10pt,after skip=8pt]
    \poppill{Dans cette leçon}{poppurple}{white}\par\vspace{5pt}
    \popsemi\fontsize{9.3}{12.6}\selectfont\color{popink}#7
  \end{tcolorbox}
  \noindent\begin{minipage}[t]{0.487\linewidth}
    \begin{tcolorbox}[enhanced,colback=popgreen,colframe=popink,boxrule=2.2pt,arc=10pt,
        drop shadow=popink,left=11pt,right=11pt,top=10pt,bottom=10pt,height=3.1cm,valign=center]
      \tcbox[colback=white,colframe=popink,boxrule=1.3pt,arc=5pt,boxsep=0pt,nobeforeafter,
        left=3pt,right=3pt,top=3pt,bottom=3pt,valign=center]{\qrcode[height=1.9cm]{https://play.google.com/store/apps/details?id=com.luvvix.onbuch&hl=fr}}%
      \hspace{8pt}\begin{minipage}[c]{0.5\linewidth}
        {\popdisplay\fontsize{11}{12.5}\selectfont\color{white}\MakeUppercase{Télécharge OnBuch}}\par\vspace{4pt}
        {\raggedright\popsemi\fontsize{8.4}{11}\selectfont\color{white}Gratuit \textbullet\ Google Play\\Tuteur IA, annales, concours, résultats\par}
      \end{minipage}
    \end{tcolorbox}
  \end{minipage}\hfill
  \begin{minipage}[t]{0.487\linewidth}
    \begin{tcolorbox}[enhanced,colback=poppurple,colframe=popink,boxrule=2.2pt,arc=10pt,
        drop shadow=popink,left=11pt,right=11pt,top=10pt,bottom=10pt,height=3.1cm,valign=center]
      \tikz[baseline=-0.7ex]\node[circle,fill=white,draw=popink,line width=1.3pt,minimum size=1.6cm,inner sep=0]{\popdisplay\fontsize{24}{24}\selectfont\color{poppurple}+};%
      \hspace{8pt}\begin{minipage}[c]{0.6\linewidth}
        {\popdisplay\fontsize{11}{12.5}\selectfont\color{white}\MakeUppercase{Passe à OnBuch+}}\par\vspace{4pt}
        {\raggedright\popsemi\fontsize{8.4}{11}\selectfont\color{white}Tous les cours signés, Tuteur IA illimité, fiches PDF — onglet OnBuch+ de l'app\par}
      \end{minipage}
    \end{tcolorbox}
  \end{minipage}
  \vspace{8pt}
  \popcontenu
  \vfill
  \signatureonbuch
  \restoregeometry
  \newpage
}

% Rangée « Ce cours contient » (page de garde)
\newcommand{\poptile}[3]{% #1 couleur #2 gros texte #3 légende
  \begin{tcolorbox}[enhanced,colback=#1,colframe=popink,boxrule=2pt,arc=9pt,shadow={2.5pt}{-2.5pt}{0pt}{popink},
      left=6pt,right=6pt,top=9pt,bottom=9pt,halign=center,valign=center,height=1.75cm,width=\linewidth,nobeforeafter]
    {\popdisplay\fontsize{12.5}{13}\selectfont\color{white}\MakeUppercase{#2}}\par\vspace{3pt}
    {\centering\popsemi\fontsize{7.8}{9.5}\selectfont\color{white}#3\par}
  \end{tcolorbox}}
\newcommand{\popcontenu}{%
  \par\noindent{\popxboldls\fontsize{7.5}{7.5}\selectfont\color{popmuted}CE COURS SIGNÉ CONTIENT}\par\vspace{7pt}
  \noindent\begin{minipage}[t]{0.235\linewidth}\poptile{popblue}{Cours}{complet, illustré, pas à pas}\end{minipage}\hfill
  \begin{minipage}[t]{0.235\linewidth}\poptile{poporange}{Méthodes}{et exercices résolus}\end{minipage}\hfill
  \begin{minipage}[t]{0.235\linewidth}\poptile{poppink}{Exercices}{type examen + corrigés}\end{minipage}\hfill
  \begin{minipage}[t]{0.235\linewidth}\poptile{popgreen}{Fiche bilan}{l'essentiel à retenir}\end{minipage}\par}

% Sommaire
\newtcolorbox{sommaire}{enhanced,colback=white,colframe=popink,boxrule=2.2pt,arc=10pt,
  drop shadow=popink,left=18pt,right=18pt,top=13pt,bottom=13pt,before skip=6pt,after skip=20pt,
  before upper={\poppill{Sommaire}{poppurple}{white}\par\vspace{9pt}\popsemi\fontsize{10.6}{14.5}\selectfont\color{popink}}}

% Signature de fin de cours — poussée tout en bas de la dernière page
\newcommand{\findecours}{%
  \par\vfill\nopagebreak
  \begin{center}\popsquiggle{poporange}\end{center}\vspace{4pt}
  \signatureonbuch
}
\AtBeginDocument{\def\llbracket{\mathopen{[\mkern-3.5mu[}}\def\rrbracket{\mathclose{]\mkern-3.5mu]}}} % intervalles d'entiers (glyphe ⟦ absent de la police)
% Glyphes absents de Fira Math : redéfinis à partir de glyphes disponibles
\AtBeginDocument{%
  \def\setminus{\mathbin{\backslash}}\let\smallsetminus\setminus
  \def\vdots{\mathord{\vbox{\baselineskip3.2pt\lineskiplimit0pt\kern4pt\hbox{.}\hbox{.}\hbox{.}}}}%
  \def\ddots{\mathinner{\mkern1mu\raise6.5pt\hbox{.}\mkern2mu\raise3.6pt\hbox{.}\mkern2mu\raise0.7pt\hbox{.}\mkern1mu}}%
  \def\triangle{\mathord{\scalebox{0.9}{$\symup{\Delta}$}}}%
  \def\nabla{\mathord{\raisebox{\depth}{\scalebox{1}[-1]{$\symup{\Delta}$}}}}%
  \def\checkmark{\popcheck{popdark}}%
  \def\star{\mathbin{\ast}}\def\bigstar{\mathord{\ast}}%
  \def\diamond{\mathbin{\scalebox{0.6}{$\Diamond$}}}%
  \def\bigcup{\mathop{\mathchoice{\raisebox{-0.3em}{\scalebox{1.7}{$\cup$}}}{\scalebox{1.25}{$\cup$}}{\cup}{\cup}}\displaylimits}%
  \def\bigcap{\mathop{\mathchoice{\raisebox{-0.3em}{\scalebox{1.7}{$\cap$}}}{\scalebox{1.25}{$\cap$}}{\cap}{\cap}}\displaylimits}%
  \def\lesseqqgtr{\mathrel{\lessgtr}}\def\bigoplus{\mathop{\oplus}\displaylimits}%
}
\providecommand{\ssi}{si et seulement si} % abréviation fréquente
\AtBeginDocument{\providecommand{\wideparen}{\overparen}} % arc de cercle

%% FIGFIX-BEGIN v5 — correctifs globaux des figures (docs/onbuch-generation/tools/figfix)
\usepackage{adjustbox}\usepackage{etoolbox}\tcbuselibrary{raster}
% toute figure TikZ/pgfplots plus large que la ligne est réduite à la largeur disponible
\BeforeBeginEnvironment{tikzpicture}{\begin{adjustbox}{max width=\linewidth}}
\AfterEndEnvironment{tikzpicture}{\end{adjustbox}}
% légendes pgfplots lisibles par-dessus les courbes
\pgfplotsset{every axis/.append style={legend style={fill=white,fill opacity=0.92,text opacity=1,draw=black!15,inner sep=3pt}}}
% étiquettes d'arêtes (syntaxe edge["..."]) avec fond blanc
\tikzset{every edge quotes/.append style={fill=white,inner sep=1.2pt}}
%% FIGFIX-END
\begin{document}

\pagegarde{Leçon 3 — Expression orale : amour de la famille et la solidarité familiale}{Chinois}{1ère A}{Module 1 : Vie familiale et intégration sociale}{ZHA03}{2 h}{Cette leçon d'expression orale entraîne les élèves de 1ère A à présenter un exposé simple et à participer à un débat guidé en chinois sur le thème de l'amour de la famille et de la solidarité familiale. Elle fournit les dialogues modèles, les formules fonctionnelles et le travail des tons nécessaires à une prise de parole claire et correcte.
\vspace{4pt}
\begin{multicols}{2}\small
\begin{itemize}
\item À la fin de cette leçon, je sais saluer et me présenter avant un exposé oral en chinois
\item À la fin de cette leçon, je sais présenter ma famille et exprimer l'amour familial avec des phrases simples
\item À la fin de cette leçon, je sais utiliser les formules fonctionnelles pour donner mon avis (我觉得..., 我认为...)
\item À la fin de cette leçon, je sais exprimer l'accord et le désaccord poliment dans un débat
\item À la fin de cette leçon, je sais relancer la parole avec des questions simples (你呢？你觉得怎么样？)
\item À la fin de cette leçon, je sais prononcer correctement les quatre tons et le ton neutre dans les mots-clés de la leçon
\end{itemize}\end{multicols}}

\begin{sommaire}
\begin{multicols}{2}
\begin{enumerate}
\item Mise en route : parler de sa famille
\item Dialogues modèles
\item Formules fonctionnelles pour exposer et débattre
\item Les tons à travailler
\item Jeux de rôle guidés
\item Débat guidé : la famille d'abord ?
\item Bilan et entraînement à l'exposé final
\item Activité d'intégration
\item Méthodes et exercices résolus
\item Exercices
\item Corrigés des exercices
\item Fiche bilan
\end{enumerate}
\end{multicols}
\end{sommaire}

\begin{objectifs}
\begin{itemize}
\item À la fin de cette leçon, je sais saluer et me présenter avant un exposé oral en chinois
\item À la fin de cette leçon, je sais présenter ma famille et exprimer l'amour familial avec des phrases simples
\item À la fin de cette leçon, je sais utiliser les formules fonctionnelles pour donner mon avis (我觉得..., 我认为...)
\item À la fin de cette leçon, je sais exprimer l'accord et le désaccord poliment dans un débat
\item À la fin de cette leçon, je sais relancer la parole avec des questions simples (你呢？你觉得怎么样？)
\item À la fin de cette leçon, je sais prononcer correctement les quatre tons et le ton neutre dans les mots-clés de la leçon
\item À la fin de cette leçon, je sais jouer un dialogue à deux voix sur la solidarité familiale
\item À la fin de cette leçon, je sais conclure un exposé avec une formule de politesse (谢谢大家！)
\end{itemize}
\end{objectifs}

\begin{prerequis}
\begin{itemize}
\item Vocabulaire de la famille : 爸爸, 妈妈, 哥哥, 姐姐, 弟弟, 妹妹, 爷爷, 奶奶
\item Structures de base : 我是..., 我家有...口人, 我爱...
\item Les quatre tons du mandarin et le pinyin
\item Verbes fréquents : 爱, 帮助, 有, 是, 做
\item Nombres et classificateur 口 pour compter les membres de la famille
\end{itemize}
\end{prerequis}

\newpage

\coursec{Mise en route : parler de sa famille}

\courssub{Situation d'accroche}

Au Cameroun, quand un oncle tombe malade à Douala, toute la famille se mobilise : les cousins de Yaoundé envoient de l'argent par Mobile Money, les tantes préparent le repas, les enfants rendent visite. En Chine aussi, la famille est le cœur de la société : le respect des aînés et la solidarité entre générations sont des valeurs fondamentales. Mais saurais-tu dire en chinois \zh{我爱我的家人} « j'aime ma famille » ou \zh{我们互相帮助} « nous nous aidons mutuellement » ?

\textbf{Problématique :} comment présenter sa famille en chinois et exprimer l'amour et la solidarité familiale, afin de préparer un exposé oral et un débat en classe ?

Aujourd'hui, tu vas apprendre à présenter ta famille, à donner ton avis et à débattre sur l'amour familial et la solidarité, comme lors d'un exposé devant la classe.

\courssub{Observation : un premier texte}

Lis d'abord ce petit texte de présentation, écrit par une élève chinoise de ton âge. Observe les phrases soulignées : que remarques-tu sur leur construction ?

\begin{textebox}[Ma famille — 我的家]
我家有六口人：爸爸、妈妈、哥哥、姐姐、妹妹和我。我很爱我的家人。\\
Wǒ jiā yǒu liù kǒu rén : bàba, māma, gēge, jiějie, mèimei hé wǒ. Wǒ hěn ài wǒ de jiārén.\\
« Il y a six personnes dans ma famille : papa, maman, grand frère, grande sœur, petite sœur et moi. J'aime beaucoup ma famille. »
\end{textebox}

\textbf{Observations guidées :}
\begin{itemize}
\item La phrase commence par \zh{我家} « ma famille » : le sujet vient en premier, comme en français.
\item \zh{有} yǒu exprime « il y a » : pas besoin de verbe « être ».
\item \zh{六口人} : le nombre \zh{六} liù « six » est suivi du classificateur \zh{口} kǒu, le mot de mesure pour compter les membres d'une famille.
\item \zh{很爱} hěn ài : l'adverbe \zh{很} hěn « très » se place \emph{avant} le verbe \zh{爱} ài « aimer ».
\end{itemize}

\courssub{Réactivation du vocabulaire : les membres de la famille}

Tu as déjà rencontré ces mots dans les leçons précédentes. Réactivons-les : ils seront tes outils de base pour tout l'exposé.

\begin{center}
\begin{tabularx}{\linewidth}{|X|X|X|X|}
\hline
\rowcolor{popblueL} Caractères & Pinyin & Nature & Traduction \\
\hline
家 & jiā & nom & famille, foyer \\
\hline
家人 & jiārén & nom & membres de la famille (collectif) \\
\hline
爸爸 & bàba & nom & papa \\
\hline
妈妈 & māma & nom & maman \\
\hline
哥哥 & gēge & nom & grand frère \\
\hline
姐姐 & jiějie & nom & grande sœur \\
\hline
弟弟 & dìdi & nom & petit frère \\
\hline
妹妹 & mèimei & nom & petite sœur \\
\hline
爷爷 & yéye & nom & grand-père (paternel) \\
\hline
奶奶 & nǎinai & nom & grand-mère (paternelle) \\
\hline
爱 & ài & verbe & aimer \\
\hline
帮助 & bāngzhù & verbe & aider \\
\hline
团结 & tuánjié & adjectif/verbe & uni ; s'unir \\
\hline
互相 & hùxiāng & adverbe & mutuellement, l'un l'autre \\
\hline
\end{tabularx}
\end{center}

\begin{attention}[Le redoublement affectueux]
Les mots de famille chinois sont souvent \cle{redoublés} : \zh{爸爸}, \zh{妈妈}, \zh{哥哥}. Attention à la prononciation : la \emph{deuxième} syllabe se prononce au \textbf{ton neutre} (léger, sans accent) : \emph{bàba}, \emph{māma}, \emph{gēge}. Ne dis pas \emph{bàbà} avec deux tons marqués : cela sonnerait artificiel. C'est le même phénomène qu'en français familier « papa », « maman ».
\end{attention}

\begin{savaistu}[Le caractère 家 jiā]
En Chine, le caractère \zh{家} jiā « famille » représente un cochon (\zh{豕}) sous un toit (\zh{宀}) : autrefois, avoir un cochon sous son toit signifiait avoir un foyer prospère. La famille et la maison ne font qu'un dans la culture chinoise : on dit \zh{回家} huí jiā « rentrer à la maison » aussi bien que « retrouver sa famille ».
\end{savaistu}

\courssub{Carte mentale : le lexique de la famille}

\begin{popfigure}
\begin{center}\tcbox[colback=popblue,colframe=popblue,coltext=white,arc=9pt,boxsep=4pt,left=6pt,right=6pt]{\bfseries\small 家 jiā famille}\end{center}
\vspace{-2pt}
\begin{tcbraster}[raster columns=2,raster equal height=rows,raster column skip=6pt,raster row skip=6pt,enhanced,blankest]
\begin{tcolorbox}[enhanced,colback=popblue,colframe=popblue,coltext=white,boxrule=0.9pt,arc=3pt,left=4pt,right=4pt,top=3pt,bottom=3pt,boxsep=1pt,halign=left]\bfseries\scriptsize 家人 jiārén membres 爸爸 妈妈 哥哥 姐姐\end{tcolorbox}
\begin{tcolorbox}[enhanced,colback=poporange,colframe=poporange,coltext=white,boxrule=0.9pt,arc=3pt,left=4pt,right=4pt,top=3pt,bottom=3pt,boxsep=1pt,halign=left]\bfseries\scriptsize 爱 ài amour 我爱我家\end{tcolorbox}
\begin{tcolorbox}[enhanced,colback=popgreen,colframe=popgreen,coltext=white,boxrule=0.9pt,arc=3pt,left=4pt,right=4pt,top=3pt,bottom=3pt,boxsep=1pt,halign=left]\bfseries\scriptsize 帮助 bāngzhù aide 互相帮助\end{tcolorbox}
\begin{tcolorbox}[enhanced,colback=poppurple,colframe=poppurple,coltext=white,boxrule=0.9pt,arc=3pt,left=4pt,right=4pt,top=3pt,bottom=3pt,boxsep=1pt,halign=left]\bfseries\scriptsize 团结 tuánjié union 很团结\end{tcolorbox}
\end{tcbraster}

\legende{Carte mentale du lexique de la famille : quatre branches autour du mot-clé 家.}
\end{popfigure}

\courssub{La phrase de présentation : 我家有五口人}

\begin{propriete}[Présenter le nombre de personnes de sa famille]
Structure : \textbf{Sujet (我家) + 有 + nombre + 口 + 人}.\\
\zh{我家有五口人。} \emph{Wǒ jiā yǒu wǔ kǒu rén.} « Il y a cinq personnes dans ma famille. »
\begin{itemize}
\item \zh{有} yǒu = « il y a / avoir » : c'est le verbe d'existence et de possession.
\item \zh{口} kǒu est le \cle{classificateur} (mot de mesure) des membres d'une famille. En chinois, on ne peut pas juxtaposer directement un nombre et un nom : il faut un classificateur entre les deux.
\item Le mot \zh{人} rén « personne » reste invariable : pas de pluriel en chinois.
\end{itemize}
\end{propriete}

\begin{exemplebox}[Déclinaisons]
\zh{我家有四口人。} \emph{Wǒ jiā yǒu sì kǒu rén.} « Il y a quatre personnes dans ma famille. »\\
\zh{他家有八口人。} \emph{Tā jiā yǒu bā kǒu rén.} « Il y a huit personnes dans sa famille. »\\
\zh{你家有几口人？} \emph{Nǐ jiā yǒu jǐ kǒu rén?} « Combien y a-t-il de personnes dans ta famille ? »
\end{exemplebox}

\begin{attention}[Erreurs fréquentes des francophones]
\begin{itemize}
\item Ne dis pas \zh{我家是五口人} : le français « ma famille \emph{est} de cinq personnes » pousse à utiliser \zh{是} shì « être », mais le chinois exige \zh{有} yǒu.
\item N'oublie pas le classificateur : \zh{我家有五人} est incorrect ; il faut \zh{我家有五口人}.
\item Ordre figé : nombre + 口 + 人, jamais \zh{五人口}.
\end{itemize}
\end{attention}

\courssub{Exprimer l'amour et la solidarité}

\begin{propriete}[Le verbe 爱 ài et l'adverbe 很 hěn]
Structure : \textbf{Sujet + (很) + 爱 + complément}.\\
\zh{我爱我的家人。} \emph{Wǒ ài wǒ de jiārén.} « J'aime ma famille. »\\
\zh{我很爱我的妈妈。} \emph{Wǒ hěn ài wǒ de māma.} « J'aime beaucoup ma maman. »
\begin{itemize}
\item \zh{很} hěn « très » se place toujours \emph{avant} le verbe ou l'adjectif qu'il nuance, jamais après (contrairement au français « beaucoup » qui suit le verbe).
\item \zh{的} de marque la possession : \zh{我的家人} « ma famille » (litt. « les membres-de-famille de moi »).
\end{itemize}
\end{propriete}

\begin{propriete}[Décrire les relations : 互相帮助]
Structure : \textbf{Sujet pluriel + 互相 + verbe}.\\
\zh{我们互相帮助。} \emph{Wǒmen hùxiāng bāngzhù.} « Nous nous aidons mutuellement. »\\
L'adverbe \zh{互相} hùxiāng « mutuellement » remplace le pronom réciproque français « se / l'un l'autre ». Il se place juste avant le verbe. Le chinois n'a pas besoin de pronom réfléchi : \zh{互相} suffit.
\end{propriete}

\begin{exemplebox}[Exemples traduits]
\zh{我们一家人很团结。} \emph{Wǒmen yì jiā rén hěn tuánjié.} « Notre famille est très unie. »\\
\zh{弟弟常常帮助我。} \emph{Dìdi chángcháng bāngzhù wǒ.} « Mon petit frère m'aide souvent. »\\
\zh{爷爷奶奶很爱我们。} \emph{Yéye nǎinai hěn ài wǒmen.} « Papi et mamie nous aiment beaucoup. »\\
\zh{姐姐和我互相学习。} \emph{Jiějie hé wǒ hùxiāng xuéxí.} « Ma grande sœur et moi apprenons l'un de l'autre. »
\end{exemplebox}

\begin{attention}[家人 est déjà collectif]
Ne dis jamais \zh{我爱我的家人们} : \zh{家人} jiārén désigne déjà l'ensemble des membres de la famille ; on n'ajoute pas le suffixe de pluriel \zh{们} men. En revanche, \zh{们} est correct avec les pronoms et les noms de personnes : \zh{我们} « nous », \zh{同学们} « les camarades ».
\end{attention}

\courssub{Comparaison chinois / français}

\begin{center}
\begin{tabularx}{\linewidth}{|X|X|X|}
\hline
\rowcolor{popblueL} Notion & Chinois & Français \\
\hline
Nombre de membres & 我家有五口人。(有 + nombre + 口 + 人) & « Il y a cinq personnes dans ma famille. » \\
\hline
Possession & 我的家人 (的 obligatoire) & « ma famille » (adjectif possessif) \\
\hline
Nuance & 很爱 (很 avant le verbe) & « aimer beaucoup » (adverbe après) \\
\hline
Réciprocité & 互相帮助 (adverbe 互相) & « s'aider mutuellement » (pronom réfléchi) \\
\hline
Pluriel & 家人 (collectif, sans 们) & « les membres de la famille » (pluriel marqué) \\
\hline
\end{tabularx}
\end{center}

\courssub{Méthode : présenter sa famille à l'oral}

\begin{methode}[Les trois étapes de la présentation]
Pour présenter ta famille à l'oral, suis toujours cet ordre :
\begin{enumerate}
\item \textbf{Le nombre de personnes} : \zh{我家有……口人。} « Il y a ... personnes dans ma famille. »
\item \textbf{La liste des membres} : \zh{爸爸、妈妈、哥哥和我。} « papa, maman, mon grand frère et moi. » (utilise \zh{、} entre les éléments et \zh{和} hé « et » avant le dernier).
\item \textbf{Le sentiment} : \zh{我很爱我的家人。} « J'aime beaucoup ma famille. » Tu peux ajouter une phrase sur la solidarité : \zh{我们互相帮助。}
\end{enumerate}
Parle lentement, soigne les tons, et regarde ton auditoire.
\end{methode}

\courssub{Exercice résolu}

\begin{exoresolu}[Remets la phrase dans l'ordre]
Énoncé : remets les mots suivants dans le bon ordre pour former une phrase correcte : \zh{口人 / 我家 / 有 / 五}. Puis traduis-la en français.
\tcblower
Solution : le sujet vient d'abord (\zh{我家}), puis le verbe d'existence (\zh{有}), puis le nombre suivi de son classificateur et du nom (\zh{五口人}).\\
\zh{我家有五口人。} \emph{Wǒ jiā yǒu wǔ kǒu rén.} « Il y a cinq personnes dans ma famille. »\\
Vérification : Sujet + 有 + nombre + 口 + 人. La structure est respectée, le classificateur \zh{口} n'est pas oublié.
\end{exoresolu}

\begin{exoresolu}[Traduis en chinois]
Énoncé : traduis en chinois : « Nous nous aidons mutuellement. J'aime beaucoup ma famille. »
\tcblower
Solution : pour la réciprocité, on utilise l'adverbe \zh{互相} avant le verbe ; pour la nuance, \zh{很} avant \zh{爱}.\\
\zh{我们互相帮助。我很爱我的家人。}\\
\emph{Wǒmen hùxiāng bāngzhù. Wǒ hěn ài wǒ de jiārén.}\\
Piège évité : pas de \zh{们} après \zh{家人}, et \zh{很} bien placé avant le verbe.
\end{exoresolu}

\courssub{Échauffement oral : à toi de parler !}

\begin{experience}[Présente ta famille à ton voisin]
Par groupes de deux, chaque élève présente sa famille en \textbf{deux phrases minimum} à son voisin, en suivant la méthode des trois étapes :
\begin{enumerate}
\item \zh{我家有……口人。}
\item \zh{有爸爸、妈妈、……和我。}
\item \zh{我很爱我的家人。我们互相帮助。}
\end{enumerate}
Le voisin écoute, puis pose une question : \zh{你家有几口人？} « Combien de personnes dans ta famille ? » ou \zh{你爱你的家人吗？} « Aimes-tu ta famille ? ». Ensuite, on échange les rôles. Le professeur circule et corrige les tons.
\end{experience}

\begin{aretenir}[Bilan de la mise en route]
\begin{itemize}
\item Présenter sa famille : \zh{我家有 + nombre + 口人。}
\item Exprimer l'amour : \zh{我很爱我的家人。} (\zh{很} avant le verbe ; pas de \zh{们} après \zh{家人}).
\item Exprimer la solidarité : \zh{我们互相帮助。} (\zh{互相} = « mutuellement »).
\item Classificateur obligatoire : \zh{口} kǒu pour les membres de la famille.
\item Ton neutre sur la deuxième syllabe des mots redoublés : \emph{bàba}, \emph{māma}, \emph{gēge}.
\end{itemize}
\end{aretenir}

\coursec{Dialogues modèles}

\courssub{Transition depuis la mise en route}
Après avoir évoqué en classe les membres de votre famille et les mots-clés du vocabulaire familial (\zh{爸爸、妈妈、哥哥、姐姐、弟弟、妹妹}), nous passons maintenant à l'\textbf{oral interactif}. Deux dialogues types vous permettront d'observer les structures de base pour \textbf{présenter sa famille} (Dialogue 1) et pour \textbf{exprimer la solidarité et l'entraide} (Dialogue 2). Écoutez, répétez, analysez, puis appropriez-vous ces modèles pour vos propres exposés et débats.

\courssub{Dialogue 1 : Présentation de la famille entre deux camarades}
\begin{textebox}[Dialogue 1 — Présentation de la famille]
A : 你家有几口人？\\
Nǐ jiā yǒu jǐ kǒu rén ?\\
(Combien de personnes y a-t-il dans ta famille ?)

B : 我家有四口人：爸爸、妈妈、妹妹和我。\\
Wǒ jiā yǒu sì kǒu rén : bàba, māma, mèimei hé wǒ.\\
(Quatre : papa, maman, ma petite sœur et moi.)

A : 你爱你的家人吗？\\
Nǐ ài nǐ de jiārén ma ?\\
(Tu aimes ta famille ?)

B : 当然！我很爱他们。\\
Dāngrán ! Wǒ hěn ài tāmen.\\
(Bien sûr ! Je les aime beaucoup.)
\end{textebox}

\begin{propriete}[Analyse linguistique du Dialogue 1]
\textbf{1. Structure interrogative : \cle{有} + \cle{几} + classifieur + nom}\\
Schéma : \textbf{Sujet + 有 + 几 + 口/个 + Nom ?}\\
\zh{几} (jǐ) s'utilise pour demander un nombre inférieur à 10. Le classifieur \zh{口} (kǒu) est spécifique pour compter les membres d'une famille (bouches à nourrir). Pour les personnes en général, on utilise \zh{个} (gè).

\textbf{2. Énumération avec \cle{和} (hé)}\\
\zh{爸爸、妈妈、妹妹和我} : la virgule chinoise \zh{、} sépare les éléments, \zh{和} relie le dernier. En français, on met « et » avant le dernier ; en chinois, \zh{和} ne relie que deux noms (ou groupes nominaux), pas des phrases entières.

\textbf{3. Question fermée avec \cle{吗} (ma)}\\
\zh{你爱你的家人吗？} : ajout de la particule \zh{吗} en fin de phrase affirmative pour former une question oui/non. L'ordre des mots ne change pas.

\textbf{4. Adverbe de degré \cle{很} (hěn) + Verbe d'état}\\
\zh{我很爱他们} : \zh{很} (hěn) lie le sujet au verbe d'état/adjectif \zh{爱} (àimer). Attention : \zh{很} ici a souvent une valeur grammaticale (liaison) plus qu'intensive (« très »). Pour dire « vraiment beaucoup », on utilisera \zh{非常} (fēicháng) ou \zh{特别} (tèbié).
\end{propriete}

\begin{center}
\begin{tabularx}{\linewidth}{|X|X|X|X|}
\hline
\rowcolor{popblueL} \textbf{Caractères} & \textbf{Pinyin} & \textbf{Nature} & \textbf{Traduction} \\
\hline
有 & yǒu & verbe & avoir, il y a \\
\hline
几 & jǐ & pronom interrogatif & combien (moins de 10) \\
\hline
口 & kǒu & classifieur & (pour membres famille) \\
\hline
和 & hé & conjonction & et (entre noms) \\
\hline
爱 & ài & verbe & aimer \\
\hline
家人 & jiārén & nom & membres de la famille, proches \\
\hline
当然 & dāngrán & adverbe & bien sûr, naturellement \\
\hline
他们 & tāmen & pronom & ils, elles (personnes) \\
\hline
\end{tabularx}
\end{center}

\courssub{Dialogue 2 : Entraide familiale — Aider un frère malade, partager les tâches}
\begin{textebox}[Dialogue 2 — Solidarité et entraide]
A : 你哥哥生病了，你怎么办？\\
Nǐ gēge shēngbìng le, nǐ zěnme bàn ?\\
(Ton grand frère est malade, que fais-tu ?)

B : 我照顾他，给他做饭。\\
Wǒ zhàogù tā, gěi tā zuòfàn.\\
(Je prends soin de lui, je lui prépare à manger.)

A : 你真好！\\
Nǐ zhēn hǎo !\\
(Tu es vraiment gentil !)

B : 家人应该互相帮助。\\
Jiārén yīnggāi hùxiāng bāngzhù.\\
(Les membres de la famille doivent s'entraider.)
\end{textebox}

\begin{propriete}[Analyse linguistique du Dialogue 2]
\textbf{1. Aspect accompli / changement d'état : \cle{了} (le) après le verbe}\\
\zh{生病了} (shēngbìng le) : \zh{了} marque ici la survenue d'un nouvel état (« est tombé malade » / « est malade maintenant »). Ce n'est pas seulement le passé, c'est l'actualisation de la situation.

\textbf{2. Verbe composé \cle{照顾} (zhàogù) — \textbf{Point de prononciation crucial}}\\
Voir la boîte \textbf{Attention} ci-dessous. Structure : \textbf{Sujet + 照顾 + Objet (personne)}.

\textbf{3. Construction \cle{给} (gěi) + Bénéficiaire + Verbe d'action}\\
\zh{给他做饭} (gěi tā zuòfàn) : « faire à manger \textbf{pour} lui ». \zh{给} introduit le bénéficiaire de l'action. Ordre : \zh{给} + Personne + Verbe. Comparer : \zh{我做饭给他吃} (Je cuisine pour qu'il mange / Je lui cuisine à manger).

\textbf{4. Modal \cle{应该} (yīnggāi) « devoir, doit »}\\
\zh{家人应该互相帮助} : exprime une obligation morale, un devoir. Négation : \zh{不应该} (bù yīnggāi).

\textbf{5. Réciprocité : \cle{互相} (hùxiāng) + Verbe}\\
\zh{互相帮助} (s'entraider), \zh{互相理解} (se comprendre mutuellement), \zh{互相尊敬} (se respecter mutuellement). Sujet obligatoirement pluriel ou collectif.
\end{propriete}

\begin{center}
\begin{tabularx}{\linewidth}{|X|X|X|X|}
\hline
\rowcolor{poporangeL} \textbf{Caractères} & \textbf{Pinyin} & \textbf{Nature} & \textbf{Traduction} \\
\hline
哥哥 & gēge & nom & grand frère \\
\hline
生病 & shēngbìng & verbe/VO & tomber malade, être malade \\
\hline
怎么办 & zěnme bàn & locution & que faire ? comment faire ? \\
\hline
照顾 & zhàogù & verbe & prendre soin de, soigner \\
\hline
给 & gěi & préposition/verbe & pour (bénéficiaire), donner \\
\hline
做饭 & zuòfàn & VO & cuisiner, faire à manger \\
\hline
真 & zhēn & adverbe & vraiment, véritablement \\
\hline
应该 & yīnggāi & verbe modal & devoir (obligation morale) \\
\hline
互相 & hùxiāng & adverbe & mutuellement, l'un l'autre \\
\hline
帮助 & bāngzhù & verbe/nom & aider, aide \\
\hline
\end{tabularx}
\end{center}

\courssub{Travail de la prononciation : les tons à surveiller}
\begin{attention}[Prononciation de \zh{照顾} (zhàogù) — Deux 4e tons]
Le mot \zh{照顾} combine deux quatrièmes tons (tons descendants : \textbf{ˋ ˋ}).
\begin{itemize}
    \item \zh{照} (zhào) : départ haut, chute brutale. Attention au \zh{zh-} (rétroflexe, langue recourbée vers le palais dur) + \zh{-ao} (diphongue a-o ouvert).
    \item \zh{顾} (gù) : même contour descendant. \zh{g-} (vélaire non aspiré) + \zh{-u} (u arrondi, lèvres très rondes).
\end{itemize}
\textbf{Erreur fréquente des francophones :} prononcer le deuxième ton comme un 1er ton (haut plat) ou un 2e ton (montant) par fatigue, ou « avaler » la chute du premier. \textbf{Exercice :} tapez dans les mains sur chaque syllabe en exagérant la chute : \zh{ZHÀO! GÙ!} Répétez 5 fois de suite : \zh{照顾、照顾、照顾、照顾、照顾}.
\end{attention}

\begin{attention}[Autres pièges de tons dans les dialogues]
\begin{itemize}
    \item \zh{几} (jǐ, 3e ton) vs \zh{鸡} (jī, 1er ton, poulet). Contexte : \zh{几口人}.
    \item \zh{和} (hé, 2e ton) vs \zh{河} (hé, 2e ton, fleuve) vs \zh{喝} (hē, 1er ton, boire). Ici conjonction \zh{hé}.
    \item \zh{爱} (ài, 4e ton) : chute nette. Ne pas confondre avec \zh{艾} (ài, herbe) ou \zh{矮} (ǎi, petit de taille, 3e ton).
    \item \zh{病} (bìng, 4e ton) : finale \zh{-ing} (ng nasal vélaire). Ne pas dire « bine ».
\end{itemize}
\end{attention}

\courssub{Méthode de lecture et mémorisation progressive}
\begin{methode}[Lecture à voix haute et appropriation en 3 étapes]
\begin{enumerate}
    \item \textbf{Lecture pinyin seule :} Cachez les caractères. Lisez le pinyin à voix haute 3 fois. Concentrez-vous \textbf{uniquement} sur la mélodie des tons (contours) et l'enchaînement des syllabes (liaisons). Vérifiez les 4e tons de \zh{zhàogù}, \zh{shēngbìng}, \zh{ài}.
    \item \textbf{Lecture caractères + pinyin :} Découvrez les caractères. Lisez en pointant du doigt chaque caractère. Associez forme, son et sens. Repérez les mots-outils : \zh{了、和、给、应该、互相}.
    \item \textbf{Lecture caractères seuls (sans pinyin) :} Cachez le pinyin. Lisez le texte chinois. Si vous bloquez, ne regardez que le mot difficile.
    \item \textbf{Jeu de rôle sans texte :} En binôme, jouez la scène. A pose la question, B répond. Inversez les rôles. Puis changez les informations (voir exercice de substitution ci-dessous).
\end{enumerate}
\end{methode}

\courssub{Exercices de substitution en binôme (Variantes contextuelles)}
L'objectif est de rendre le dialogue vivant et personnel. Travaillez en binôme : l'un joue A, l'autre B. Changez les éléments soulignés.

\begin{exemplebox}[Substitution Dialogue 1 — Taille et composition de la famille]
\begin{tabular}{ll}
\textbf{Modèle :} & \zh{我家有四口人：爸爸、妈妈、妹妹和我。} \\
\textbf{Variante 1 (Famille nombreuse, Cameroun) :} & \zh{我家有七口人：爸爸、妈妈、两个哥哥、一个姐姐、妹妹和我。} \\
 & (Wǒ jiā yǒu qī kǒu rén : bàba, māma, liǎng gè gēge, yí gè jiějie, mèimei hé wǒ.) \\
\textbf{Variante 2 (Famille monoparentale) :} & \zh{我家有两口人：妈妈和我。} \\
 & (Wǒ jiā yǒu liǎng kǒu rén : māma hé wǒ.) \\
\textbf{Variante 3 (Avec grands-parents) :} & \zh{我家有六口人：爷爷、奶奶、爸爸、妈妈、弟弟和我。} \\
 & (Wǒ jiā yǒu liù kǒu rén : yéye, nǎinai, bàba, māma, dìdi hé wǒ.)
\end{tabular}
\end{exemplebox}

\begin{exemplebox}[Substitution Dialogue 2 — Situations d'entraide]
\begin{tabular}{ll}
\textbf{Modèle :} & \zh{我照顾他，给他做饭。} \\
\textbf{Variante 1 (Sœur malade) :} & \zh{我陪着她，给她倒水、喂药。} \\
 & (Wǒ péizhe tā, gěi tā dào shuǐ, wèi yào. — Je la tiens compagnie, lui verse de l'eau, lui donne ses médicaments.) \\
\textbf{Variante 2 (Maman fatiguée, tâches ménagères) :} & \zh{我帮妈妈洗衣服、扫地。} \\
 & (Wǒ bāng māma xǐ yīfu, sǎo dì. — J'aide maman à laver le linge, balayer.) \\
\textbf{Variante 3 (Petit frère devoirs) :} & \zh{我教弟弟写汉字、练习拼音。} \\
 & (Wǒ jiāo dìdi xiě hànzì, liànxí pīnyīn. — J'apprends à mon petit frère à écrire les caractères, réviser le pinyin.)
\end{tabular}
\end{exemplebox}

\courssub{Identification des formules fonctionnelles (Boîte à outils oral)}
Repérez dans les deux dialogues les « briques » réutilisables pour votre exposé et le débat.

\begin{center}
\begin{tabularx}{\linewidth}{|X|X|X|}
\hline
\rowcolor{popgreenL} \textbf{Fonction communicative} & \textbf{Formule chinoise (Caractères + Pinyin)} & \textbf{Équivalent français} \\
\hline
Demander la composition familiale & \zh{你家有几口人？} (Nǐ jiā yǒu jǐ kǒu rén ?) & Combien de personnes chez toi ? \\
\hline
Présenter sa famille (énumération) & \zh{我家有[数字]口人：[Liste] 和 我。} & Nous sommes [nombre] : [Liste] et moi. \\
\hline
Exprimer l'affection & \zh{我(很/非常/特别)爱[人]。} & J'aime [personne] (beaucoup/énormément). \\
\hline
Demander la réaction face à un problème & \zh{……，你怎么办？} (…, nǐ zěnme bàn ?) & …, que fais-tu ? / Comment réagis-tu ? \\
\hline
Décrire une action de soin & \zh{我照顾[人]，给[人][Action]。} & Je soigne [personne], je fais [action] pour [personne]. \\
\hline
Exprimer l'admiration / validation & \zh{你真[Adj]！} (Nǐ zhēn … !) & Tu es vraiment [Adj] ! (gentil, courageux, filial) \\
\hline
Énoncer un principe moral (débat) & \zh{[Sujet] 应该 [Verbe/Verbe+Obj]。} & [Sujet] doit / devrait [Action]. \\
\hline
Exprimer la réciprocité & \zh{[Sujet pluriel] 互相 [Verbe]。} & [Sujet] se [Verbe] mutuellement. \\
\hline
\end{tabularx}
\end{center}

\courssub{Exercice résolu : Transformer le dialogue en exposé court}
\begin{exoresolu}[Transformer le Dialogue 1 en mini-exposé de 30 secondes]
\textbf{Consigne :} À partir du Dialogue 1, produisez un monologue cohérent (5-6 phrases) pour vous présenter vous et votre famille, en utilisant les connecteurs logiques \zh{然后} (ensuite), \zh{而且} (de plus), \zh{因为……所以……} (parce que… donc).

\textbf{Solution détaillée :}
Voici une proposition de mini-exposé structuré :

\bigskip
\noindent
\textbf{Mini-exposé : Ma famille}
\begin{itemize}
    \item \zh{大家好。} \emph{Dàjiā hǎo.} (Bonjour à tous.)
    \item \zh{我叫[Votre Prénom]，我家有五口人。} \emph{Wǒ jiào [Votre Prénom], wǒ jiā yǒu wǔ kǒu rén.} (Je m'appelle…, ma famille compte 5 personnes.)
    \item \zh{有爸爸、妈妈、哥哥、姐姐和我。} \emph{Yǒu bàba, māma, gēge, jiějie hé wǒ.} (Il y a papa, maman, grand frère, grande sœur et moi.)
    \item \zh{我哥哥在上大学，我姐姐在工作。} \emph{Wǒ gēge zài shàng dàxué, wǒ jiějie zài gōngzuò.} (Mon grand frère est à l'université, ma grande sœur travaille.)
    \item \zh{我们很亲密，因为我们互相帮助，互相爱护。} \emph{Wǒmen hěn qīnmì, yīnwèi wǒmen hùxiāng bāngzhù, hùxiāng àihù.} (Nous sommes très proches, parce que nous nous entraidons, nous nous chérissons mutuellement.)
    \item \zh{我爱我的家。} \emph{Wǒ ài wǒ de jiā.} (J'aime ma famille.)
\end{itemize}

\textbf{Analyse des structures utilisées :}
\begin{itemize}
    \item \zh{我叫…} (Présentation nom) — HSK 1.
    \item \zh{有…} (Existence) sans sujet (sujet « ma famille » sous-entendu ou topic-comment).
    \item \zh{在 + Verbe} (Action en cours / situation actuelle) : \zh{在上大学、在工作}.
    \item \zh{因为……所以……} (Cause-conséquence) : structure clé pour l'argumentation (débat).
    \item \zh{互相 + Verbe} (Réciprocité) réemployée du Dialogue 2.
\end{itemize}
\end{exoresolu}

\courssub{Mémorisation progressive : Dialogue troncé (Technique du « trou »)}
Appliquez cette technique chez vous ou en binôme pour ancrer la fluidité. Lisez la version complète, puis la version à trous, essayez de la dire de mémoire, puis refermez le livre.

\begin{exemplebox}[Dialogue 1 — Version à trous (niveau 1 : 1 mot sur 2)]
A : 你家 \_\_\_\_ 几 \_\_\_\_ 人？ \\
B : 我家 \_\_\_\_ 四 \_\_\_\_ 人：爸爸、\_\_\_\_、妹妹 \_\_\_\_ 我。 \\
A : 你 \_\_\_\_ 你的 \_\_\_\_ 吗？ \\
B : \_\_\_\_！我 \_\_\_\_ 爱 \_\_\_\_.
\end{exemplebox}

\begin{exemplebox}[Dialogue 2 — Version à trous (niveau 2 : mots-outils retirés)]
A : 你哥哥 生病 \_\_\_\_，你 怎么办？ \\
B : 我 \_\_\_\_ 他，\_\_\_\_ 他 做饭。 \\
A : 你 真 好！ \\
B : 家人 \_\_\_\_ 互相 \_\_\_\_.
\end{exemplebox}

\courssub{Schéma visuel : Structure interactionnelle A $\leftrightarrow$ B}
\begin{popfigure}
\begin{tikzpicture}[
    node distance=1.2cm and 3cm,
    boxA/.style={draw=popblue, thick, fill=popblueL, rounded corners, minimum width=4cm, minimum height=1.2cm, align=center, font=\small},
    boxB/.style={draw=poporange, thick, fill=poporangeL, rounded corners, minimum width=4cm, minimum height=1.2cm, align=center, font=\small},
    arrow/.style={-Stealth, thick, popdark},
    labelArrow/.style={midway, fill=white, inner sep=2pt, font=\tiny, text=popdark}
]
% Dialogue 1
\node[boxA, align=center] (A1) at (0, 3){A : Question ouverte\\ \zh{你家有几口人？}};
\node[boxB, align=center] (B1) at (6, 3){B : Info factuelle\\ \zh{我家有四口人…}};
\node[boxA, align=center] (A2) at (0, 1.5){A : Question sentiment\\ \zh{你爱你的家人吗？}};
\node[boxB, align=center] (B2) at (6, 1.5){B : Affirmation + Intensité\\ \zh{当然！我很爱他们。}};

% Dialogue 2
\node[boxA, align=center] (A3) at (0, -1){A : Situation problème\\ \zh{哥哥生病了，怎么办？}};
\node[boxB, align=center] (B3) at (6, -1){B : Action concrète\\ \zh{照顾、给做饭}};
\node[boxA, align=center] (A4) at (0, -2.5){A : Validation sociale\\ \zh{你真好！}};
\node[boxB, align=center] (B4) at (6, -2.5){B : Principe moral (Débat)\\ \zh{家人应该互相帮助}};

% Arrows Dialogue 1
\draw[arrow] (A1.east) -- node[labelArrow] {Demande info} (B1.west);
\draw[arrow, popblue] (B1.west) -- node[labelArrow] {Réponse} (A1.east); % retour implicite
\draw[arrow] (A2.east) -- node[labelArrow] {Sonde affect} (B2.west);
\draw[arrow, popblue] (B2.west) -- node[labelArrow] {Engagement} (A2.east);

% Arrows Dialogue 2
\draw[arrow] (A3.east) -- node[labelArrow] {Provoque action} (B3.west);
\draw[arrow] (A4.east) -- node[labelArrow] {Renforce} (B4.west);
\draw[arrow, poporange] (B4.west) -- node[labelArrow] {Justifie} (A4.east);

% Brackets labels
\node[align=center, font=\small\bfseries, text=popblue] at (-3.5, 2.25) {Dialogue 1\\\textbf{Présentation}\\\textbf{\& Affect}};
\node[align=center, font=\small\bfseries, text=poporange] at (-3.5, -1.75) {Dialogue 2\\\textbf{Action}\\\textbf{\& Valeur}};
\end{tikzpicture}
\legende{Schéma des échanges : Locuteur A (bleu) initie, relance, valide ; Locuteur B (orange) informe, agit, justifie. Structure typique Exposé $\rightarrow$ Débat.}
\end{popfigure}

\courssub{Exemples supplémentaires pour l'entraînement personnel}
\begin{itemize}
    \item \zh{妈妈每天给我们做饭。} \emph{Māma měitiān gěi wǒmen zuòfàn.} (Maman nous prépare à manger tous les jours.) — Structure \zh{给} + bénéficiaire.
    \item \zh{我们应该尊敬老人。} \emph{Wǒmen yīnggāi zūnjìng lǎorén.} (Nous devons respecter les personnes âgées.) — \zh{尊敬} (zūnjìng, respecter) vocabulaire HSK 3/4, valeur culturelle forte.
    \item \zh{弟弟生病了，姐姐照顾他。} \emph{Dìdi shēngbìng le, jiějie zhàogù tā.} (Le petit frère est malade, la grande sœur s'occupe de lui.) — Remplacement pronominal \zh{他} (lui).
    \item \zh{在喀麦隆，大家都说“家是港湾”。} \emph{Zài Kāmàilóng, dàjiā dōu shuō “Jiā shì gǎngwān”.} (Au Cameroun, tout le monde dit « La famille est un port/havre ».) — Expression idiomatique \zh{港湾} (gǎngwān), lien socio-culturel.
\end{itemize}

\courssub{Bilan de la section : Check-list avant de passer aux jeux de rôle}
\begin{itemize}
    \item[$\square$] Je sais compter ma famille avec \zh{口} (kǒu) et énumérer avec \zh{和} (hé).
    \item[$\square$] Je prononce correctement les deux 4e tons de \zh{照顾} (zhàogù).
    \item[$\square$] Je sais utiliser \zh{给} + personne + verbe pour dire « faire qqch pour qqn ».
    \item[$\square$] Je peux exprimer un devoir moral avec \zh{应该} (yīnggāi) et la réciprocité avec \zh{互相} (hùxiāng).
    \item[$\square$] J'ai joué les deux dialogues en binôme sans lire le texte (étape 3 de la méthode).
\end{itemize}

\emph{Prêts ? Passons maintenant à la \textbf{Section 3 : Formules fonctionnelles pour exposer et débattre} pour outiller votre prise de parole longue.}

\coursec{Leçon 3 — Expression orale : Amour de la famille et solidarité familiale}

\courssub{Mise en route : parler de sa famille}

\begin{methode}[Échauffement oral (5 min)]
\begin{enumerate}
\item \textbf{Tour de table rapide} : Chaque élève dit une phrase sur sa famille en utilisant \zh{我家有...人} (Wǒ jiā yǒu ... rén) et \zh{我最爱...} (Wǒ zuì ài...).
\item \textbf{Mots-clés au tableau} : L'enseignant note les mots revenus : \zh{父母} (fùmǔ), \zh{兄弟姐妹} (xiōngdì jiěmèi), \zh{爷爷奶奶} (yéye nǎinai), \zh{外公外婆} (wàigōng wàipó), \zh{团结} (tuánjié), \zh{孝顺} (xiàoshùn), \zh{分担} (fēndān), \zh{家务} (jiāwù).
\item \textbf{Question déclencheur} : \zh{“家里谁做家务最多？为什么？”} (Jiālǐ shéi zuò jiāwù zuì duō? Wèishéme?) — Passage naturel vers le thème de la solidarité.
\end{enumerate}
\end{methode}

\begin{center}
\begin{tabularx}{\linewidth}{|X|X|X|X|}
\hline
\rowcolor{popblueL} \textbf{Caractères} & \textbf{Pinyin} & \textbf{Nature} & \textbf{Traduction} \\
\hline
家庭 & jiātíng & n. & famille, foyer \\
亲人 & qīnrén & n. & proches, parents (au sens large) \\
团结 & tuánjié & adj./v. & uni, solidaire / s'unir, se serrer les coudes \\
孝顺 & xiàoshùn & adj./v. & pieux, respectueux (envers les parents) / respecter, honorer \\
分担 & fēndān & v. & partager (une charge), répartir \\
家务 & jiāwù & n. & tâches ménagères \\
责任 & zérèn & n. & responsabilité \\
关爱 & guānai & n./v. & soin, affection / chérir, prendre soin de \\
理解 & lǐjiě & v./n. & comprendre / compréhension \\
支持 & zhīchí & v./n. & soutenir / soutien \\
\hline
\end{tabularx}
\end{center}

\courssub{Dialogues modèles}

\begin{textebox}[Dialogue 1 : Le week-end en famille (Registre courant)]
\textbf{A :} \zh{周末你去哪儿了？} (Zhōumò nǐ qù nǎr le?) \\
\textbf{B :} \zh{我陪父母去市场买菜，顺便帮他们拿东西。} (Wǒ péi fùmǔ qù shìchǎng mǎicài, shùnbiàn bāng tāmen ná dōngxi.) \\
\textbf{A :} \zh{真孝顺！我周末通常都在家看手机。} (Zhēn xiàoshùn! Wǒ zhōumò tōngcháng dōu zài jiā kàn shǒujī.) \\
\textbf{B :} \zh{少看点手机，多陪陪家人吧。家庭团结最重要。} (Shǎo kàn diǎn shǒujī, duō péipei jiārén ba. Jiātíng tuánjié zuì zhòngyào.) \\
\textbf{Traduction :} \\
A : Tu as fait quoi le week-end ? \\
B : J'ai accompagné mes parents au marché acheter des légumes, et j'ai porté les affaires pour eux. \\
A : Vraiment filial ! Moi le week-end je reste d'habitude à la maison à regarder mon téléphone. \\
B : Regarde moins le téléphone, accompagne plus ta famille. L'unité familiale est le plus important.
\end{textebox}

\begin{textebox}[Dialogue 2 : Discussion sur le partage des tâches (Registre semi-formel)]
\textbf{A :} \zh{我觉得做家务不只是大人的事。} (Wǒ juéde zuò jiāwù bù zhǐ shì dàrén de shì.) \\
\textbf{B :} \zh{我同意。孩子分担家务能培养责任感。} (Wǒ tóngyì. Háizi fēndān jiāwù néng péiyǎng zérèngǎn.) \\
\textbf{A :} \zh{而且能体谅父母的辛苦。比如洗碗、扫地、倒垃圾。} (Érqiě néng tǐliàng fùmǔ de xīnkǔ. Bǐrú xǐwǎn, sǎodì, dào lājī.) \\
\textbf{B :} \zh{对。只要大家动手，家里就干净又温馨。} (Duì. Zhǐyào dàjiā dòngshǒu, jiālǐ jiù gānjìng yòu wēnxīn.) \\
\textbf{Traduction :} \\
A : Je pense que faire les tâches ménagères n'est pas seulement l'affaire des adultes. \\
B : Je suis d'accord. Les enfants qui partagent les tâches développent le sens des responsabilités. \\
A : Et puis ça permet de comprendre la peine des parents. Par exemple laver la vaisselle, balayer, sortir les poubelles. \\
B : Exact. Tant que tout le monde met la main à la pâte, la maison est propre et chaleureuse.
\end{textebox}

\courssub{Formules fonctionnelles pour exposer et débattre}

Après avoir analysé les \emph{dialogues modèles} ci-dessus, où les interlocuteurs échangeaient de manière naturelle sur leur famille, nous passons maintenant à la prise de parole en continu : l'\textbf{exposé oral structuré} et le \textbf{débat argumenté}. En classe de Première (niveau HSK 3 visé), vous devez être capable de tenir un discours cohérent de 1 à 2 minutes, d'interagir avec l'auditoire et de gérer l'accord ou le désaccord avec politesse.

\courssub{Les cinq étapes d'un exposé oral réussi}

Un bon exposé oral en chinois suit une progression logique, immuable, que l'on peut schématiser ainsi :

\begin{methode}[Structure type d'un exposé oral (1--2 minutes)]
\emph{Étape 1 -- Saluer et se présenter} : \zh{大家好！我叫...} (Dàjiā hǎo! Wǒ jiào...) \\
\emph{Étape 2 -- Annoncer le sujet clairement} : \zh{今天我要谈一谈...} (Jīntiān wǒ yào tántan...) \\
\emph{Étape 3 -- Développer : Opinion + Exemples/Arguments} : \zh{我觉得... 因为... 而且...} (Wǒ juéde... yīnwèi... érqiě...) \\
\emph{Étape 4 -- Conclure brièvement} : \zh{总之...} (Zǒng'ér...) ou \zh{所以...} (Suǒyǐ...) \\
\emph{Étape 5 -- Remercier l'auditoire} : \zh{我的话说完了，谢谢大家！} (Wǒ de huà shuō wán le, xièxie dàjiā!)
\end{methode}

\begin{popfigure}
\begin{tikzpicture}[
    node distance=1.2cm and 1.5cm,
    every node/.style={font=\small, align=center},
    box/.style={rectangle, rounded corners, draw=popblue, thick, fill=popblueL, minimum width=3.2cm, minimum height=1.4cm},
    arrow/.style={-Stealth, thick, popdark}
]
\node[box, align=center] (s1){\textbf{1. Salutation}\\\zh{大家好！我叫...}};
\node[box, right=of s1, align=center] (s2){\textbf{2. Sujet}\\\zh{今天我要谈一谈...}};
\node[box, right=of s2, align=center] (s3){\textbf{3. Développement}\\\zh{我觉得... 因为...}\\\zh{而且... 例如...}};
\node[box, right=of s3, align=center] (s4){\textbf{4. Conclusion}\\\zh{总之... / 所以...}};
\node[box, right=of s4, align=center] (s5){\textbf{5. Remerciement}\\\zh{谢谢大家！}};

\draw[arrow] (s1) -- (s2);
\draw[arrow] (s2) -- (s3);
\draw[arrow] (s3) -- (s4);
\draw[arrow] (s4) -- (s5);
\end{tikzpicture}
\legende{Organigramme du déroulé type d'un exposé oral structuré}
\end{popfigure}

\courssub{Se présenter et annoncer le sujet}

Ces formules sont le « sésame » de toute prise de parole publique. Elles fixent le cadre formel.

\begin{center}
\begin{tabularx}{\linewidth}{|X|X|X|X|}
\hline
\rowcolor{popblueL} \textbf{Caractères} & \textbf{Pinyin} & \textbf{Nature} & \textbf{Traduction} \\
\hline
大家好 & dàjiā hǎo & Loc. & Bonjour à tous \\
老师好 & lǎoshī hǎo & Loc. & Bonjour professeur \\
我叫... & wǒ jiào... & Phrase & Je m'appelle... \\
我是... & wǒ shì... & Phrase & Je suis... (nom / classe) \\
今天我要谈一谈... & jīntiān wǒ yào tántan... & Phrase & Aujourd'hui je vais parler de... \\
我的题目是... & wǒ de tímù shì... & Phrase & Mon sujet est... \\
\hline
\end{tabularx}
\end{center}

\begin{propriete}[Annoncer le sujet : \zh{谈一谈} vs \zh{介绍一下}]
\zh{谈一谈} (tántan) « parler un peu de, discuter de » $\rightarrow$ marque un ton \textbf{conversationnel}, idéal pour un exposé en classe.
\zh{介绍一下} (jièshào yíxià) « présenter brièvement » $\rightarrow$ plus \textbf{formel}, utilisé pour une présentation de projet ou un exposé officiel.
Structure : \zh{今天} (Sujet temps) + \zh{我要/想} (Verbe modal) + \zh{谈一谈/介绍一下} (Verbe redupliqué) + \zh{题目/话题} (Objet).
\end{propriete}

\begin{exemplebox}[Exemples de début d'exposé]
\zh{大家好！我叫李明，是高一A班的学生。今天我要谈一谈“家庭的爱与团结”。} \\
Dàjiā hǎo! Wǒ jiào Lǐ Míng, shì Gāoyī A bān de xuéshēng. Jīntiān wǒ yào tántan “Jiātíng de ài yǐ tuánjié”. \\
« Bonjour à tous ! Je m'appelle Li Ming, élève de 1ère A. Aujourd'hui je vais parler de "L'amour et l'unité dans la famille". »

\zh{老师好，同学们好。我的题目是“孝顺与分担家务”。我来介绍一下我的看法。} \\
Lǎoshī hǎo, tóngxuémen hǎo. Wǒ de tímù shì “Xiàoshùn yǔ fēndān jiāwù”. Wǒ lái jièshào yíxià wǒ de kànfǎ. \\
« Bonjour professeur, bonjour camarades. Mon sujet est "La piété filiale et le partage des tâches ménagères". Je vais présenter mon point de vue. »
\end{exemplebox}

\courssub{Exprimer son opinion : \zh{我觉得} vs \zh{我认为}}

C'est le cœur de l'exposé (Étape 3). Le choix du verbe marque le registre de langue.

\begin{propriete}[Nuance entre \zh{我觉得} et \zh{我认为}]
\begin{itemize}
\item \zh{我觉得} (wǒ juéde) : « Je trouve que / J'ai l'impression que / Je pense que (subjectif) ». \textbf{Registre courant, oral, sincère}. Utilisé pour exprimer un ressenti, une préférence, une opinion personnelle basée sur l'expérience.
\item \zh{我认为} (wǒ rènwéi) : « J'estime que / Je considère que / Mon avis est que ». \textbf{Registre plus formel, écrit, argumenté}. Utilisé dans un débat structuré, une dissertation orale, pour avancer une thèse.
\end{itemize}
\textbf{Structure commune :} \zh{Sujet + 觉得/认为 + [Proposition complète (Sujet + Verbe + Complément)]}.
\end{propriete}

\begin{attention}[Piège majeur des francophones : l'ellipse du sujet dans la proposition subordonnée]
En français, on dit : « Je pense \textbf{que c'est} important » (sujet « ce » obligatoire).
En chinois, \zh{我觉得} / \zh{我认为} introduit une proposition \textbf{qui doit garder son propre sujet} si ce n'est pas le même que le principal, ou répéter le sujet logique.
\begin{itemize}
\item \textbf{Faux} : \zh{我觉得很重要。} (Wǒ juéde hěn zhòngyào.) $\leftarrow$ Il manque le sujet de « important » : important \textbf{pour qui / quoi ?}
\item \textbf{Vrai} : \zh{我觉得\textbf{家庭}很重要。} (Wǒ juéde \textbf{jiātíng} hěn zhòngyào.) « Je pense que \textbf{la famille} est importante. »
\item \textbf{Vrai} : \zh{我觉得\textbf{这样做}很好。} (Wǒ juéde \textbf{zhèyàng zuò} hěn hǎo.) « Je trouve que \textbf{faire ainsi} est bien. »
\end{itemize}
\cle{Règle d'or} : Après \zh{我觉得/我认为}, posez-vous la question : « Qui/Quoi fait l'action ou a la qualité ? » $\rightarrow$ Écrivez-le en chinois.
\end{attention}

\begin{exemplebox}[Opinions sur la famille et la solidarité]
\zh{我觉得家人比什么都重要。} \\
Wǒ juéde jiārén bǐ shénme dōu zhòngyào. \\
« Je pense que la famille est plus importante que tout. » (Subjectif, affectif)

\zh{我认为团结是家庭幸福的基础。} \\
Wǒ rènwéi tuánjié shì jiātíng xìngfú de jīchǔ. \\
« J'estime que la solidarité est le fondement du bonheur familial. » (Argumenté, thèse)

\zh{我觉得做家务不是麻烦，而是一种关爱。} \\
Wǒ juéde zuò jiāwù bù shì máfan, ér shì yī zhǒng guānai. \\
« Je trouve que faire le ménage n'est pas une corvée, mais une forme de soin/attention. »

\zh{我认为父母应该以身作则，教导孩子分担责任。} \\
Wǒ rènwéi fùmǔ yīnggāi yǐshēnzuòzé, jiàodǎo háizi fēndān zérèn. \\
« J'estime que les parents doivent donner l'exemple et éduquer les enfants à partager les responsabilités. »
\end{exemplebox}

\courssub{Interagir : Exprimer l'accord, le désaccord poli et relancer}

Un débat ne se résume pas à juxtaposer des monologues. Il faut des \textbf{connecteurs interactionnels}.

\begin{center}
\begin{tabularx}{\linewidth}{|X|X|X|X|}
\hline
\rowcolor{popblueL} \textbf{Fonction} & \textbf{Caractères / Pinyin} & \textbf{Registre} & \textbf{Traduction / Usage} \\
\hline
Accord fort & \zh{我完全同意。 / 我赞成。} (Wǒ wánquán tóngyì. / Wǒ zànchéng.) & Formel / Standard & Je suis tout à fait d'accord. / J'approuve. \\
Accord naturel & \zh{对！/ 是啊！/ 正是如此。} (Duì! / Shì a! / Zhèngshì rúcǐ.) & Oral / Courant & Exact ! / Oui ! / C'est bien cela. \\
Accord partiel & \zh{你说得有道理，但是...} (Nǐ shuō de yǒu dàoli, dànshì...) & Poli & Tu as raison sur un point, mais... \\
\hline
Désaccord poli & \zh{对不起，我不同意。} (Duìbuqǐ, wǒ bù tóngyì.) & Standard & Pardon, je ne suis pas d'accord. \\
Désaccord nuancé & \zh{我不太赞同这个看法。} (Wǒ bù tài zàntóng zhège kànfǎ.) & Formel & Je ne souscris pas tout à fait à cet avis. \\
Transition & \zh{可是... / 不过... / 但是...} (Kěshì... / Bùguò... / Dànshì...) & Tous & Mais / Cependant / Toutefois. \\
\hline
Relance (ouvert) & \zh{你呢？} (Nǐ ne?) & Oral & Et toi ? \\
Relance (avis) & \zh{你觉得怎么样？} (Nǐ juéde zěnmeyàng?) & Standard & Qu'en penses-tu ? / Quel est ton avis ? \\
Relance (précision) & \zh{能不能具体说说？} (Néng bù néng jùtǐ shuōshuo?) & Formel & Pourrais-tu être plus précis ? \\
\hline
\end{tabularx}
\end{center}

\begin{propriete}[La politesse du désaccord : \zh{对不起} + \zh{不同意}]
En chinois, le désaccord direct (\zh{你错了} « Tu as tort ») est perçu comme agressif (\zh{没面子} « perdre la face »).
La structure standard \textbf{excuse + désaccord} (\zh{对不起，我不同意}) ou \textbf{concession + mais} (\zh{你说得对，可是...} « Tu as raison, mais... ») préserve l'harmonie relationnelle (\zh{和气} héqì), valeur centrale tant au Cameroun qu'en Chine.
\end{propriete}

\begin{exemplebox}[Échanges types en débat]
\textbf{A :} \zh{我觉得手机让家人疏远了。你觉得怎么样？} (Wǒ juéde shǒujī ràng jiārén shūyuǎn le. Nǐ juéde zěnmeyàng?) \\
« Je trouve que le téléphone éloigne les membres de la famille. Qu'en penses-tu ? »

\textbf{B (Accord) :} \zh{我完全同意。大家都看屏幕，不说话。} (Wǒ wánquán tóngyì. Dàjiā dōu kàn píngmǔ, bù shuōhuà.) \\
« Je suis tout à fait d'accord. Tout le monde regarde l'écran, ne parle pas. »

\textbf{C (Désaccord poli) :} \zh{对不起，我不太赞同。视频通话让我们联系远方的亲戚。} (Duìbuqǐ, wǒ bù tài zàntóng. Shìpìn tōnghuà ràng wǒmen liánxì yuǎnfāng de qīnqi.) \\
« Pardon, je ne suis pas tout à fait d'accord. Les appels vidéo nous permettent de contacter la famille lointaine. »

\textbf{A (Relance) :} \zh{有道理。不过，面对面交流更重要。} (Yǒu dàoli. Bùguò, miànduìmiàn jiāoliú gèng zhòngyào.) \\
« Tu as raison. Cependant, l'échange en face à face est plus important. »
\end{exemplebox}

\courssub{Conclure et remercier : clore l'échange avec élégance}

La conclusion (Étape 4) résume l'idée force ; le remerciement (Étape 5) marque la fin officielle de votre temps de parole.

\begin{exemplebox}[Formules de conclusion et de clôture]
\zh{总之，家庭的爱需要我们用行动来维护。} \\
Zǒng'ér, jiātíng de ài xūyào wǒmen yòng xíngdòng lái wéihù. \\
« En résumé, l'amour familial a besoin que nous le maintenions par des actions. »

\zh{所以，请大家多陪伴家人，多说“我爱你们”。} \\
Suǒyǐ, qǐng dàjiā duō péibàn jiārén, duō shuō “Wǒ ài nǐmen”. \\
« Donc, je vous invite tous à accompagner davantage votre famille, à dire plus souvent "Je vous aime". »

\zh{我的发言完毕，谢谢大家聆听！} \\
Wǒ de fāyán wánbì, xièxie dàjiā língtīng! \\
« Mon intervention est terminée, merci à tous pour votre écoute ! » (Plus formel que \zh{我的话说完了}).
\end{exemplebox}

\begin{savaistu}[Le proverbe clé : \zh{团结就是力量}]
\zh{团结就是力量} (Tuánjié jiùshì lìliang) — « L'union fait la force ».
C'est l'un des proverbes (\zh{成语} chéngyǔ) les plus connus en Chine, inscrit sur le fronton de la Grande Salle du Peuple à Pékin. Il résume parfaitement la valeur de \zh{团结} (tuánjié, solidarité/unité) au sein de la famille comme dans la nation. L'utiliser en conclusion d'un exposé sur la solidarité familiale fait forte impression : \zh{正如那句话所说：团结就是力量，家庭团结，万事兴。} (Jiùrú nà jù huà suǒ shuō: Tuánjié jiùshì lìliang, jiātíng tuánjié, wànshì xīng. — « Comme le dit ce proverbe : l'union fait la force, famille unie, tout prospère. »).
\end{savaistu}

\courssub{Les tons à travailler : points de vigilance}

La prononciation correcte des tons est cruciale pour l'intelligibilité de votre exposé. Voici les pièges fréquents dans le vocabulaire de cette leçon.

\begin{propriete}[Tons à surveiller (HSK 3)]
\begin{itemize}
\item \textbf{Ton 3 + Ton 3 (Sandhi) :} \zh{你好} $\rightarrow$ \zh{ní hǎo} ; \zh{很好} $\rightarrow$ \zh{hěn hǎo} ; \zh{我也} $\rightarrow$ \zh{wó yě}. Règle : le premier ton 3 devient ton 2 (montant).
\item \textbf{Ton neutre (轻声 qīngshēng) :} Particules de structure \zh{了} (le), \zh{吗} (ma), \zh{呢} (ne), \zh{吧} (ba) ; suffixes \zh{一下} (yíxià), \zh{一谈} (yítán \rightarrow tántan), \zh{说说} (shuōshuo). Ne pas les "chanter".
\item \textbf{Homophones / Tons proches :}
      \begin{itemize}
      \item \zh{赚} (zhuàn, 4e ton) « gagner de l'argent » vs \zh{转} (zhuǎn, 3e ton) « tourner / transférer ».
      \item \zh{重} (zhòng, 4e ton) « important / lourd » vs \zh{种} (zhǒng, 3e ton) « espèce / planter ».
      \item \zh{累} (lèi, 4e ton) « fatigué » vs \zh{雷} (léi, 2e ton) « tonnerre ».
      \end{itemize}
\item \textbf{Enchaînement \zh{因为... 所以...} :} \zh{yīnwèi} (4e-4e) \zh{suǒyǐ} (3e-3e $\rightarrow$ 2e-3e). Marquez la pause après la cause.
\end{itemize}
\end{propriete}

\begin{experience}[Atelier prononciation : "Le marché aux tons"]
\begin{enumerate}
\item \textbf{Écoute discriminative} : L'enseignant lit une liste de mots (zhuàn/zhuǎn, zhòng/zhǒng, lèi/léi, tántan/tǎntan). Les élèves lèvent la carte "1/2/3/4/0" correspondante.
\item \textbf{Lecture en chaîne} : Chaque élève lit une phrase des \emph{Dialogues modèles} ou des \emph{Exemples d'opinions}. La classe note les erreurs de tons sur une fiche de pair-évaluation.
\item \textbf{Correction collective} : Relecture chorale des phrases problématiques (ex: \zh{我觉得家庭很重要} — attention au ton de \zh{重} et au sandhi \zh{很重} $\rightarrow$ \zh{hén zhòng}).
\end{enumerate}
\end{experience}

\courssub{Jeux de rôle guidés}

\begin{experience}[Jeu de rôle 1 : "La réunion de famille" (Binôme, 3 min)]
\textbf{Rôles :} A = Grand frère / Grande sœur ; B = Cadet / Cadette.
\textbf{Situation :} Les parents sont fatigués. Il faut décider qui fait quoi ce soir (cuisiner, laver, balayer, sortir poubelles).
\textbf{Contraintes linguistiques :} Utiliser \zh{我觉得...} / \zh{我认为...} ; \zh{因为... 所以...} ; \zh{分担} ; \zh{责任} ; \zh{公平} (gōngpíng).
\textbf{Modèle de départ :}
A: \zh{弟弟/妹妹，今天爸妈很累。我觉得我们应该分担家务。} \\
B: \zh{我同意。但是我有很多作业。} \\
A: \zh{没事，我来洗碗，你扫地行吗？} \\
B: \zh{好，公平。那我倒垃圾。}
\end{experience}

\begin{experience}[Jeu de rôle 2 : "L'argent de poche" (Groupe de 3, 4 min)]
\textbf{Rôles :} Parent / Enfant A (veut gagner l'argent) / Enfant B (veut l'argent gratuit).
\textbf{Sujet :} \zh{“零花钱：给还是赚？”}
\textbf{Objectif :} Utiliser les formules d'accord/désaccord poli (\zh{对不起，我不同意}, \zh{你有道理，但是...}) et la relance \zh{你觉得怎么样？}.
\end{experience}

\courssub{Débat guidé : La famille d'abord ?}

\begin{methode}[Déroulement du débat en classe (15 min)]
\begin{enumerate}
\item \textbf{Préparation individuelle (3 min)} : Remplir la \emph{Fiche de préparation orale} (voir Exercice ci-dessous) pour l'un des deux camps : \zh{赞成} (Pour) ou \zh{反对} (Contre).
\item \textbf{Mise en commun par camp (3 min)} : Les "Pour" se regroupent, les "Contre" se regroupent. Ils affinent leurs 2 meilleurs arguments.
\item \textbf{Débat public (7 min)} : Un représentant de chaque camp fait son exposé (1 min). Ensuite, débat libre modéré par l'enseignant : relances, concessions, contre-arguments.
\item \textbf{Synthèse (2 min)} : L'enseignant note au tableau les structures bien utilisées et les erreurs de tons/grammaire à corriger collectivement.
\end{enumerate}
\end{methode}

\begin{aretenir}[Mots-outils pour le débat "Famille vs Individu"]
\begin{center}
\begin{tabularx}{\linewidth}{|X|X|X|}
\hline
\rowcolor{popblueL} \textbf{Connecteur} \& \textbf{Pinyin} \& \textbf{Usage} \\
\hline
虽然... 但是... \& suīrán... dànshì... \& Concession : "Bien que... pourtant..." \\
不仅... 而且... \& bùjǐn... érqiě... \& Accumulation : "Non seulement... mais encore..." \\
如果... 就... \& rúguǒ... jiù... \& Hypothèse : "Si... alors..." \\
首先... 其次... 最后... \& shǒuxiān... qícì... zuìhòu... \& Structuration : "D'abord... ensuite... enfin..." \\
\hline
\end{tabularx}
\end{center}
\end{aretenir}

\courssub{Bilan et entraînement à l'exposé final}

\courssub{Entraînement guidé : Reconstituer un mini-exposé}

\begin{exoresolu}[Remettre les phrases dans l'ordre logique (Étape 1 à 5)]
\textbf{Énoncé :} Voici les phrases d'un exposé de 1 minute sur « Partager les tâches ménagères » (\zh{分担家务}). Classez-les de 1 à 5 selon la structure apprise.
\tcblower
\textbf{Solution détaillée :}
\begin{enumerate}
\item \textbf{Étape 1 (Salutation) :} \zh{大家好！我叫王芳，是高一学生。} (Dàjiā hǎo! Wǒ jiào Wáng Fāng, shì Gāoyī xuéshēng.)
\item \textbf{Étape 2 (Sujet) :} \zh{今天我要谈一谈“分担家务的重要性”。} (Jīntiān wǒ yào tántan “Fēndān jiāwù de zhòngyào xìng”.)
\item \textbf{Étape 3 (Argument 1) :} \zh{我觉得做家务能锻炼我们的责任心。} (Wǒ juéde zuò jiāwù néng duànliàn wǒmen de zérènxīn.) $\rightarrow$ \zh{我觉得} + sujet \zh{做家务} + verbe \zh{能锻炼}.
\item \textbf{Étape 3 (Argument 2 - \zh{而且}) :} \zh{而且，帮父母干活能减轻他们的疲劳。} (Érqiě, bāng fùmǔ gànhuó néng jiǎnqīng tāmen de píláo.) $\rightarrow$ Connecteur \zh{而且} « de plus / en outre » pour ajouter un argument.
\item \textbf{Étape 5 (Clôture) :} \zh{我的话说完了，谢谢大家！} (Wǒ de huà shuō wán le, xièxie dàjiā!)
\end{enumerate}
\emph{Note :} L'étape 4 (Conclusion synthétique \zh{总之/所以}) est implicite ici car l'exposé est très court, mais en examen, ajoutez une phrase : \zh{所以，我们应该主动分担家务。} (Suǒyǐ, wǒmen yīnggāi zhǔdòng fēndān jiāwù.)
\end{exoresolu}

\courssub{Exercice d'application : Préparez votre fiche de débat}

\begin{exercice}[Fiche de préparation orale (10 minutes)]
\textbf{Sujet de débat :} \zh{“零花钱应该由孩子自己赚吗？”} (L'argent de poche doit-il être gagné par les enfants eux-mêmes ?)

Remplissez la grille ci-dessous en \textbf{caractères chinois} + \textbf{pinyin}. Utilisez le vocabulaire de la leçon (\zh{我觉得/我认为}, \zh{因为/所以}, \zh{而且}, \zh{比如}).

\begin{center}
\begin{tabularx}{\linewidth}{|X|X|}
\hline
\rowcolor{popblueL} \textbf{Étape} & \textbf{Votre brouillon (Caractères + Pinyin)} \\
\hline
1. Salutation + Identité & \\
\hline
2. Annonce du sujet & \zh{今天我要谈一谈...} \\
\hline
3. Mon avis principal (\zh{我觉得/我认为}) & \\
\hline
3. Argument 1 (\zh{因为...}) & \\
\hline
3. Argument 2 (\zh{而且... / 比如...}) & \\
\hline
4. Conclusion (\zh{总之 / 所以}) & \\
\hline
5. Remerciement & \\
\hline
\end{tabularx}
\end{center}

\textbf{Consigne orale :} Entraînez-vous à lire votre fiche à voix haute 3 fois en respectant les tons. Chronométrez-vous : visez 1 minute 30.
\end{exercice}

\begin{corrige}[Exercice : Fiche de préparation orale (Exemple de réponse type)]
Voici une proposition de réponse pour le camp \zh{赞成} (Pour / D'accord).

\begin{center}
\begin{tabularx}{\linewidth}{|X|X|}
\hline
\rowcolor{popblueL} \textbf{Étape} & \textbf{Exemple de réponse} \\
\hline
1. Salutation & \zh{大家好！我叫陈伟，来自高一A班。} (Dàjiā hǎo! Wǒ jiào Chén Wěi, lái zì Gāoyī A bān.) \\
\hline
2. Sujet & \zh{今天我要谈一谈“零花钱应该由孩子自己赚”。} (Jīntiān wǒ yào tántan “Línghuāqián yīnggāi yóu háizi zìjǐ zhuàn”.) \\
\hline
3. Avis principal & \zh{\textbf{我认为}孩子自己赚零花钱\textbf{很有必要}。} (Wǒ \textbf{rènwéi} háizi zìjǐ zhuàn línghuāqián \textbf{hěn yǒu bìyào}.) \\
\hline
3. Argument 1 & \zh{\textbf{因为}这样能让我们懂得金钱来之不易。} (\textbf{Yīnwèi} zhèyàng néng ràng wǒmen dǒngde jīnqián láizhībùyì.) \\
\hline
3. Argument 2 & \zh{\textbf{而且}，通过劳动赚钱能培养独立能力。\textbf{比如}洗车、辅导弟妹都可以赚钱。} (\textbf{Érqiě}, tōngguò láodòng zhuànqián néng péiyǎng dúlì nénglì. \textbf{Bǐrú} xǐchē, fǔdǎo dìmèi dōu kěyǐ zhuànqián.) \\
\hline
4. Conclusion & \zh{\textbf{总之}，自己赚钱花，心安理得，也减轻父母负担。} (\textbf{Zǒng'ér}, zìjǐ zhuàn qián huā, xīn'ānlǐdé, yě jiǎnqīng fùmǔ fùdān.) \\
\hline
5. Remerciement & \zh{我的发言完毕，谢谢大家聆听！} (Wǒ de fāyán wánbì, xièxie dàjiā língtīng!) \\
\hline
\end{tabularx}
\end{center}

\emph{Points de vigilance pour l'oral :}
\begin{itemize}
\item Prononcez distinctement \zh{赚} (zhuàn, 4e ton) vs \zh{转} (zhuǎn, 3e ton).
\item \zh{来之不易} (lái zhī bù yì) : expression idiomatique « venir sans facilité = fruit d'un dur labeur ».
\item \zh{心安理得} (xīn'ānlǐdé) : « l'esprit en paix = conscience tranquille ».
\item Liaison \zh{因为... 所以...} (yīnwèi... suǒyǐ...) : ne pas oublier \zh{所以} dans la conclusion si on a commencé par \zh{因为}.
\end{itemize}
\end{corrige}

\coursec{Les tons à travailler}

Après avoir acquis les formules fonctionnelles pour structurer un exposé et animer un débat (\emph{cf.} section précédente), il est indispensable de soigner la prononciation. En chinois, un changement de ton modifie le sens du mot : dire \zh{妈} (mā, mère) au lieu de \zh{马} (mǎ, cheval) peut prêter à sourire, mais dire \zh{吻} (wěn, baiser) au lieu de \zh{问} (wèn, demander) peut créer un malentendu gênant ! Cette section consolide les quatre tons, le ton neutre, la règle de sandhi du 3e ton et les particularités des mots de la famille, à travers l'observation, la gestuelle et des exercices ciblés.

\courssub{Les quatre tons et le ton neutre : rappel et gestes}

\begin{methode}[Associer un geste à chaque ton]
Pour ancrer la mélodie de chaque ton dans la mémoire corporelle, répétez chaque syllabe en effectuant le geste correspondant de la main droite :
\begin{enumerate}
    \item \textbf{1er ton (阴平 yīnpíng) — haut et plat} : main paume vers le bas, horizontale à hauteur des yeux, elle avance droit devant sans bouger. \emph{Exemple :} \zh{妈 mā} (mère).
    \item \textbf{2e ton (阳平 yángpíng) — montant} : main partie basse, paume vers le haut, elle monte vers le visage comme pour poser une question. \emph{Exemple :} \zh{麻 má} (chanvre).
    \item \textbf{3e ton (上声 shǎngshēng) — descendant-remontant (en V)} : main partie haute, elle descend bas (niveau ventre) puis remonte. \emph{Exemple :} \zh{马 mǎ} (cheval).
    \item \textbf{4e ton (去声 qùshēng) — chute brutale} : main partie haute, elle coupe vers le bas énergiquement comme un couperet. \emph{Exemple :} \zh{骂 mà} (gronder).
    \item \textbf{Ton neutre (轻声 qīngshēng) — léger, court, sans contour fixe} : poing fermé qui s'ouvre légèrement au niveau du menton, son faible et bref. \emph{Exemple :} 2e syllabe de \zh{妈妈 māma}.
\end{enumerate}
\emph{Exercice de classe :} L'énoncé donne le ton (1 à 5), la classe fait le geste et chante la voyelle \zh{a} sur ce ton.
\end{methode}

\begin{popfigure}
\begin{tikzpicture}[scale=1.2, >=Stealth]
% Axes
\draw[popline, thick] (0,0) -- (5.5,0) node[right] {Temps};
\draw[popline, thick] (0,0) -- (0,4.5) node[above] {Hauteur};
% Repères
\foreach \y/\lbl in {1/bas, 2/moyen, 3/haut, 4/très haut} {
    \draw[popmuted, dashed] (0,\y) -- (5,\y);
    \node[left, popmuted, font=\tiny] at (0,\y) {\lbl};
}
% 1er ton
\draw[popblue, very thick] (0.2,3.5) -- (1.3,3.5);
\node[popblue, above, font=\small\bfseries] at (0.75,3.7) {1er ton (55)};
% 2e ton
\draw[poporange, very thick] (1.7,1.5) .. controls (2.2,2.5) and (3.3,3.5) .. (4.3,4);
\node[poporange, above, font=\small\bfseries] at (3,4.2) {2e ton (35)};
% 3e ton
\draw[popgreen, very thick] (0.2,2.5) .. controls (0.7,1) and (1.8,1) .. (2.8,2.5);
\node[popgreen, below, font=\small\bfseries] at (1.5,0.5) {3e ton (214)};
% 4e ton
\draw[poppurple, very thick] (3.2,4) -- (4.2,1);
\node[poppurple, right, font=\small\bfseries] at (4.3,2.5) {4e ton (51)};
% Ton neutre (point)
\fill[poppink] (4.8,2) circle (3pt);
\node[poppink, above, font=\small\bfseries] at (4.8,2.4) {Ton neutre};
\end{tikzpicture}
\legende{Contours tonals des quatre tons pleins et du ton neutre (schéma simplifié). Les chiffres (55, 35, 214, 51) indiquent la hauteur relative de départ et d'arrivée sur une échelle de 1 (bas) à 5 (haut).}
\end{popfigure}

\begin{center}
\begin{tabularx}{\linewidth}{|X|X|X|X|}
\hline
\rowcolor{popblueL} \textbf{Ton} & \textbf{Contour (valeur Chao)} & \textbf{Geste de la main} & \textbf{Exemple canonique} \\
\hline
1er ton (阴平) & 55 (haut plat) & Main horizontale, avance droit & \zh{妈 mā} (mère) \\
\hline
2e ton (阳平) & 35 (montant) & Main monte de bas en haut & \zh{麻 má} (chanvre) \\
\hline
3e ton (上声) & 214 (descend-remonte) & Main dessine un V & \zh{马 mǎ} (cheval) \\
\hline
4e ton (去声) & 51 (chute) & Main coupe vers le bas & \zh{骂 mà} (gronder) \\
\hline
Ton neutre (轻声) & variable, court & Poing s'ouvre au menton & 2e syll. \zh{妈 ma} dans \zh{妈妈 māma} \\
\hline
\end{tabularx}
\end{center}

\courssub{Paires minimales à contraster}

L'opposition de tons sur un même segment phonétique est l'exercice le plus efficace pour affiner l'oreille et l'articulation. Travaillez ces paires en binôme : l'un dit l'un des deux mots, l'autre montre la carte du ton (1, 2, 3, 4, 0).

\begin{exemplebox}[Paire 1 : \zh{爱 ài} (4e ton) vs \zh{挨 āi} (1er ton)]
\begin{itemize}
    \item \zh{爱 ài} [4e ton] — \emph{aimer, affection}. \zh{我爱我的家人。} Wǒ ài wǒ de jiārén. (J'aime ma famille.)
    \item \zh{挨 āi} [1er ton] — \emph{être tout contre, subir (aussi prononcé ái en 2e ton pour « endurer »)}. \zh{挨家挨户 āi jiā āi hù} (de porte en porte / chaque famille).
\end{itemize}
\emph{Astuce :} \zh{爱} tombe comme un couperet (4e ton = sentiment fort qui s'affirme) ; \zh{挨} reste haut et plat (1er ton = contact prolongé, ça dure).
\end{exemplebox}

\begin{exemplebox}[Paire 2 : \zh{家 jiā} (1er ton) vs \zh{假 jiǎ} (3e ton)]
\begin{itemize}
    \item \zh{家 jiā} [1er ton] — \emph{famille, maison, foyer}. \zh{我的家在喀麦隆。} Wǒ de jiā zài Kāmàilóng. (Ma famille est au Cameroun.)
    \item \zh{假 jiǎ} [3e ton] — \emph{faux, feint ; vacances}. \zh{这是假的。} Zhè shì jiǎ de. (C'est un faux.) / \zh{放假 fàngjià} (être en vacances).
\end{itemize}
\emph{Astuce :} \zh{家} (toit + cochon) = le toit reste haut, stable (1er ton) ; \zh{假} (homme + non) = l'homme hésite, descend et remonte (3e ton).
\end{exemplebox}

\begin{exercice}[Discrimination : paires minimales]
L'enseignant prononce l'un des deux mots. Les élèves lèvent la carte « 1 », « 3 » ou « 4 » correspondante.
\begin{enumerate}
    \item \zh{爱} / \zh{挨}
    \item \zh{家} / \zh{假}
    \item \zh{妈 mā} / \zh{马 mǎ} / \zh{骂 mà}
    \item \zh{好 hǎo} / \zh{号 hào}
    \item \zh{买 mǎi} / \zh{卖 mài}
\end{enumerate}
\end{exercice}

\courssub{Mots-clés de la leçon : travail ciblé}

Le vocabulaire central de cette leçon (« amour de la famille et solidarité ») comporte des combinaisons tonales fréquentes. Maîtrisez-les en priorité.

\begin{center}
\begin{tabularx}{\linewidth}{|X|X|X|X|}
\hline
\rowcolor{popblueL} \textbf{Caractères} & \textbf{Pinyin} & \textbf{Nature} & \textbf{Traduction} \\
\hline
家人 & jiā rén (1-2) & nom & membres de la famille, proches \\
\hline
帮助 & bāng zhù (1-4) & verbe & aider, assistance \\
\hline
团结 & tuán jié (2-2) & verbe/adj & s'unir, solidarité, uni \\
\hline
照顾 & zhào gù (4-4) & verbe & prendre soin de, veiller sur \\
\hline
重要 & zhòng yào (4-4) & adjectif & important, capital \\
\hline
孝顺 & xiào shùn (4-4) & adjectif & filial, respectueux envers les parents \\
\hline
和睦 & hé mù (2-4) & adjectif & harmonieux (relations familiales) \\
\hline
支持 & zhī chí (1-2) & verbe/nom & soutenir, soutien \\
\hline
\end{tabularx}
\end{center}

\begin{exemplebox}[Analyse tonale détaillée]
\begin{itemize}
    \item \zh{帮助 bāngzhù} : 1er ton haut plat (\zh{bāng}) + 4e ton chute nette (\zh{zhù}). \textbf{Attention :} ne pas faire monter \zh{zhù} (erreur de 2e ton).
    \item \zh{团结 tuánjié} : deux 2e tons montants. Visualisez deux mains qui montent ensemble = solidarité qui grandit.
    \item \zh{照顾 zhàogù} : deux 4e tons. Deux coups de couperet secs : \zh{zhào} (regarder) + \zh{gù} (soin). Rythme martelé.
    \item \zh{重要 zhòngyào} : deux 4e tons aussi. \zh{zhòng} (lourd) + \zh{yào} (vouloir/important). Insistez sur la chute.
    \item \zh{家人 jiārén} : 1er ton + 2e ton. Montée douce après le plateau haut.
    \item \zh{孝顺 xiàoshùn} : deux 4e tons. \zh{xiào} (piété filiale) + \zh{shùn} (obéir, docile). Chutes successives, solennelles.
    \item \zh{和睦 hémù} : 2e ton montant (\zh{hé}) + 4e ton chute (\zh{mù}). Harmonie qui monte puis s'ancre.
    \item \zh{支持 zhīchí} : 1er ton plat (\zh{zhī}) + 2e ton montant (\zh{chí}). Soutien stable qui s'élève.
\end{itemize}
\end{exemplebox}

\courssub{La règle de sandhi du troisième ton}

\begin{propriete}[Sandhi du 3e ton (三声变调 sānshēng biàndiào)]
Lorsque \textbf{deux 3e tons se suivent}, le \textbf{premier se prononce comme un 2e ton} (montant 35), tandis que le second reste un 3e ton canonique (214). On écrit les tons d'origine en pinyin, mais on prononce le sandhi.
\begin{center}
\begin{tabularx}{\linewidth}{|X|X|X|}
\hline
\rowcolor{popgreenL} \textbf{Mot écrit (pinyin)} & \textbf{Prononciation réelle} & \textbf{Traduction} \\
\hline
很好 & hén hǎo (2 + 3) & très bien \\
\hline
我也 & wó yě (2 + 3) & moi aussi \\
\hline
你好 & ní hǎo (2 + 3) & bonjour \\
\hline
姐姐 & jié jie (2 + 0) & grande sœur (2e + neutre) \\
\hline
\end{tabularx}
\end{center}
\emph{Exception :} Si une pause ou une frontière syntaxique forte sépare les deux syllabes, le premier garde son 3e ton. Ex. : \zh{好 | 好} hǎo | hǎo (très très bien, avec insistance).
\end{propriete}

\begin{exemplebox}[Application aux mots de la leçon]
\begin{itemize}
    \item \zh{很好 hěn hǎo} → \emph{prononcé} \zh{hén hǎo}. \zh{这道菜很好吃。} Zhè dào cài hén hǎo chī. (Ce plat est très bon.)
    \item \zh{我也 wǒ yě} → \emph{prononcé} \zh{wó yě}. \zh{我也爱喀麦隆。} Wó yě ài Kāmàilóng. (Moi aussi j'aime le Cameroun.)
    \item \zh{姐姐 jiějie} : le premier 3e ton devient 2e ton (\zh{jié}), le second est ton neutre (\zh{jie}). Ne pas dire \zh{jiě jie} (deux tons pleins).
    \item \zh{我也想 wǒ yě xiǎng} (trois 3e tons) → \emph{prononcé} \zh{wó yé xiǎng} (2-2-3). Règle : en chaîne, tous les 3e tons non finaux deviennent 2e tons.
\end{itemize}
\end{exemplebox}

\courssub{Le ton neutre dans les redoublements familiaux}

\begin{attention}[Piège classique : le ton neutre des termes de parenté]
En chinois, les termes d'appel familiers pour les parents et frères/sœurs sont formés par \textbf{redoublement} : la première syllabe garde son ton d'origine, la \textbf{deuxième passe au ton neutre} (légère, courte, sans contour tonal marqué).
\begin{center}
\begin{tabularx}{\linewidth}{|X|X|X|}
\hline
\rowcolor{poporangeL} \textbf{Écrit (pinyin)} & \textbf{Prononciation réelle} & \textbf{Erreur fréquente (francophones)} \\
\hline
爸爸 & bà ba (4 + 0) & \zh{bà bà} (deux 4e tons) ✗ \\
\hline
妈妈 & mā ma (1 + 0) & \zh{mā mā} (deux 1er tons) ✗ \\
\hline
哥哥 & gē ge (1 + 0) & \zh{gē gē} (deux 1er tons) ✗ \\
\hline
姐姐 & jiě jie (3→2 + 0) & \zh{jiě jie} (3e + 1er/3e) ✗ \\
\hline
弟弟 & dì di (4 + 0) & \zh{dì dì} (deux 4e tons) ✗ \\
\hline
妹妹 & mèi mei (4 + 0) & \zh{mèi mei} (deux 4e tons) ✗ \\
\hline
\end{tabularx}
\end{center}
\emph{Pourquoi ?} Le redoublement exprime l'affection, la familiarité ; le ton neutre « allège » la seconde syllabe. En français, on ne change pas l'intonation pour « papa/maman » vs « père/mère », d'où la tendance à garder deux tons pleins.
\end{attention}

\begin{exemplebox}[Phrases modèles avec termes de parenté]
\begin{itemize}
    \item \zh{爸爸，我回家了。} Bà ba, wǒ huí jiā le. (Papa, je suis rentré à la maison.)
    \item \zh{妈妈在做饭。} Mā ma zài zuò fàn. (Maman cuisine.)
    \item \zh{哥哥帮助弟弟。} Gē ge bāngzhù dì di. (Le grand frère aide le petit frère.)
    \item \zh{姐姐很照顾妹妹。} Jiě jie hěn zhàogù mèi mei. (La grande sœur prend bien soin de la petite sœur.)
\end{itemize}
\emph{Exercice de chœur :} La classe répète chaque terme en exagérant le contraste : 1re syllabe forte et tonique, 2e syllabe murmurée, très courte.
\end{exemplebox}

\courssub{Exercice de discrimination auditive : « Levez la carte »}

\begin{methode}[Protocole « Levez la carte » (活动 : 举卡片)]
\begin{enumerate}
    \item \textbf{Préparation :} Chaque élève découpe 5 petites cartes (format carte de visite) et écrit en gros : \textbf{1}, \textbf{2}, \textbf{3}, \textbf{4}, \textbf{0} (pour ton neutre).
    \item \textbf{Consigne :} L'enseignant prononce \textbf{une seule syllabe ou un mot disyllabique} (vocabulaire de la leçon ou paires minimales). Pas de phrase, pas de contexte.
    \item \textbf{Réaction :} L'élève lève \textbf{immédiatement} la carte correspondant au ton de la \textbf{dernière syllabe entendue} (ou de l'unique syllabe).
    \item \textbf{Feedback :} L'enseignant scanne la salle, corrige collectivement (« Beaucoup de 3 au lieu de 4 pour \zh{zhù} »), puis refait passer le mot difficile.
    \item \textbf{Variante :} L'enseignant dit une phrase courte (\zh{我帮助妈妈}), les élèves lèvent la carte du ton du mot-cible annoncé à l'avance (ex. : « ton de \zh{帮} » → carte 1).
\end{enumerate}
\end{methode}

\begin{exoresolu}[Discrimination auditive — Mots de la leçon]
\textbf{Énoncé :} L'enseignant prononce les mots suivants dans un ordre aléatoire. L'élève note le ton de chaque syllabe (chiffres 1-4, 0 pour neutre).
\begin{enumerate}
    \item \zh{家人}
    \item \zh{帮助}
    \item \zh{团结}
    \item \zh{照顾}
    \item \zh{重要}
    \item \zh{爸爸}
    \item \zh{妈妈}
    \item \zh{哥哥}
    \item \zh{姐姐}
    \item \zh{弟弟}
\end{enumerate}

\textbf{Solution détaillée :}
\begin{enumerate}
    \item \zh{jiā rén} → \textbf{1 - 2}
    \item \zh{bāng zhù} → \textbf{1 - 4}
    \item \zh{tuán jié} → \textbf{2 - 2}
    \item \zh{zhào gù} → \textbf{4 - 4}
    \item \zh{zhòng yào} → \textbf{4 - 4}
    \item \zh{bà ba} → \textbf{4 - 0}
    \item \zh{mā ma} → \textbf{1 - 0}
    \item \zh{gē ge} → \textbf{1 - 0}
    \item \zh{jiě jie} → \textbf{2 (sandhi) - 0}
    \item \zh{dì di} → \textbf{4 - 0}
\end{enumerate}
\emph{Point de vigilance :} Items 9 (\zh{姐姐}) : le premier ton écrit est 3e, mais prononcé 2e (sandhi) + neutre. Items 6-10 : systématiser le ton neutre sur la 2e syllabe.
\end{exoresolu}

\courssub{Entraînement en contexte : texte à lire à voix haute}

Le texte suivant met en scène une famille camerounaise à Yaoundé. Il reprend le vocabulaire-cible et les phénomènes tonaux vus ci-dessus (sandhi du 3e ton, tons neutres, successions de 4e tons). Lisez-le à haute voix en respectant les contours tonaux, la liaison naturelle et l'intonation de phrase.

\begin{textebox}[Lecture expressive : « Une soirée à Mvog-Ada »]
我家住在雅温得的姆沃格-阿达区。\\ 
Wǒ jiā zhù zài Yǎwēndé de Mǔwògé-Ādá qū.\\
Ma famille habite le quartier Mvog-Ada à Yaoundé.\\

今晚，爸爸下班早，妈妈做了 ndolé 和 米饭。\\ 
Jīnwǎn, bà ba xiàbān zǎo, mā ma zuò le ndolé hé mǐfàn.\\
Ce soir, papa a fini le travail tôt, maman a préparé du ndolé et du riz.\\

哥哥 和 姐姐 帮助 弟弟 洗手。\\ 
Gē ge hé jiě jie bāngzhù dì di xǐshǒu.\\
Le grand frère et la grande sœur aident le petit frère à se laver les mains.\\

我们 围坐 在 一起，气氛 很 和睦。\\ 
Wǒmen wéizuò zài yīqì, qìfēn hěn hé mù.\\
Nous nous asseyons ensemble, l'atmosphère est très harmonieuse.\\

爸爸 说：「家人 团结 最 重要。」\\ 
Bà ba shuō: « Jiārén tuánjié zuì zhòngyào. »\\
Papa dit : « La solidarité familiale est ce qu'il y a de plus important. »\\

妈妈 照顾 大家，夹 菜 给 妹妹。\\ 
Mā ma zhàogù dàjiā, jiā cài gěi mèi mei.\\
Maman prend soin de tout le monde, elle met des morceaux dans le bol de la petite sœur.\\

我也 想 像 哥哥 一样，孝顺 父母，支持 家庭。\\ 
Wó yé xiǎng xiàng gē ge yīyàng, xiàoshun fùmǔ, zhīchí jiātíng.\\
Moi aussi je veux être comme grand frère, filial envers mes parents, soutenir la famille.\\

这 晚饭 真 好吃！\\ 
Zhè wǎnfàn zhēn hén hǎo chī!\\
Ce dîner est vraiment délicieux !
\end{textebox}

\begin{center}
\begin{tabularx}{\linewidth}{|X|X|X|}
\hline
\rowcolor{popblueL} \textbf{Mot / Expression} & \textbf{Pinyin (tons réels)} & \textbf{Point tonal à surveiller} \\
\hline
姆沃格-阿达 & Mǔwògé-Ādá & 3-4-2-1-2 (noms propres, tons canoniques) \\
\hline
下班早 & xiàbān zǎo & 4-1-3 (zǎo = 3e ton isolé = 214 complet) \\
\hline
ndolé & (emprunt) & prononcé \emph{ndo-lé} (tons camerounais) \\
\hline
帮助弟弟 & bāngzhù dìdi & 1-4 | 4-0 (deux 4e tons suivis + neutre) \\
\hline
很和睦 & hěn hé mù → \textbf{hén hé mù} & sandhi 3e+2e : hén (2) + hé (2) + mù (4) \\
\hline
最重要 & zuì zhòngyào & 4 | 4-4 (trois 4e tons d'affilée : rythme martelé) \\
\hline
照顾大家 & zhàogù dàjiā & 4-4 | 4-1 (chutes successives) \\
\hline
我也想 & wó yé xiǎng & sandhi triple 3e tons : wó (2) + yé (2) + xiǎng (3) \\
\hline
真好吃 & zhēn hén hǎo chī & sandhi hěn→hén (2) + hǎo (3) + chī (1) \\
\hline
\end{tabularx}
\end{center}

\textbf{Questions de compréhension orale (pour l'enseignant) :}
\begin{enumerate}
    \item Où habite la famille ? (Quartier, ville)
    \item Qu'a préparé maman pour le dîner ?
    \item Qui aide le petit frère à se laver les mains ?
    \item Quelle phrase dit le père sur la famille ?
    \item Quel engagement prend le narrateur à la fin ?
\end{enumerate}

\begin{savaistu}[Un ton, un sens : \zh{问} vs \zh{吻}]
En chinois, \textbf{seul le ton distingue} certains mots aux segments identiques.
\begin{itemize}
    \item \zh{问 wèn} [4e ton] = \emph{demander, interroger}. \zh{我问老师。} Wǒ wèn lǎoshī. (Je demande au prof.)
    \item \zh{吻 wěn} [3e ton] = \emph{embrasser (sur la bouche)}. \zh{他吻了她。} Tā wěn le tā. (Il l'a embrassée.)
\end{itemize}
\emph{Anecdote :} Un étudiant français a voulu dire « Je vais demander au professeur » (\zh{我去问老师}) mais a prononcé \zh{wěn} (3e ton) : « Je vais embrasser le professeur » ! La classe a bien ri. \textbf{Moralité :} le ton n'est pas une option décorative, c'est le sens même du mot.
\end{savaistu}

\coursec{Jeux de rôle guidés}

Dans la section précédente, nous avons travaillé les \cle{tons} : les distinguer, les enchaîner, les prononcer sans « aplatir » la phrase. Mais un ton correct prononcé tout seul ne suffit pas : il faut maintenant le placer dans une \emph{vraie situation de communication}. C'est exactement le rôle des \cle{jeux de rôle} : parler pour agir, pour aider, pour convaincre — comme dans une vraie famille, à Yaoundé comme à Pékin. Dans cette section, tu vas jouer trois scènes de solidarité familiale, avec des cartes de rôle, des répliques de départ et une grille d'observation remplie par tes camarades.

\courssub{Principe général : comment se déroule un jeu de rôle}

Un jeu de rôle n'est pas une récitation. Chaque élève reçoit une \cle{carte de rôle} (identité, situation, objectif) et doit improviser à partir des répliques et des formules apprises. La règle d'or : on parle \emph{uniquement en chinois}, même pour hésiter (\zh{那个……} nàge… « euh… »).

\begin{methode}[Le protocole « 2 + 3 + 2 »]
Pour chaque jeu de rôle, la séquence est toujours la même :
\begin{enumerate}
\item \textbf{Préparer (2 minutes)} : lis ta fiche de rôle, choisis tes répliques de départ, répète mentalement les tons des mots-clés.
\item \textbf{Jouer (3 minutes)} : joue la scène avec ton groupe, sans lire tes notes mot à mot. Regarde ton partenaire.
\item \textbf{Recevoir un retour (2 minutes)} : les camarades observateurs remplissent la grille d'observation et te donnent un point fort et un point à améliorer.
\end{enumerate}
Ensuite, on \cle{tourne les rôles} : chacun doit parler au moins deux fois dans la séance.
\end{methode}

\courssub{Les répliques de départ : notre boîte à outils}

Avant de jouer, observons les trois répliques-clés qui vont revenir dans toutes les scènes. Lis-les à voix haute en respectant les tons.

\begin{textebox}[Répliques de départ à mémoriser]
你怎么了？\\
Nǐ zěnme le ?\\
« Qu'est-ce que tu as ? »\\[4pt]
别担心，我来帮你。\\
Bié dānxīn, wǒ lái bāng nǐ.\\
« Ne t'inquiète pas, je viens t'aider. »\\[4pt]
我们是一家人！\\
Wǒmen shì yì jiā rén !\\
« Nous sommes une famille ! »
\end{textebox}

Observons la deuxième réplique. Elle contient deux structures essentielles de cette leçon :

\begin{itemize}
\item \zh{别} bié + verbe : l'\cle{impératif négatif} (« ne … pas »), pour rassurer ou interdire doucement ;
\item \zh{我来} wǒ lái + verbe : « c'est moi qui (va)… », la formule par excellence pour \emph{proposer son aide} spontanément.
\end{itemize}

\begin{propriete}[L'impératif négatif avec 别 bié]
\textbf{Structure} : 别 + verbe (+ complément) \\
\textbf{Emploi} : pour dire à quelqu'un de \emph{ne pas} faire quelque chose (conseil, interdiction, réconfort). C'est l'équivalent du français « ne … pas » à l'impératif. \\
\textbf{Exemples} :
\begin{itemize}
\item \zh{别担心。} Bié dānxīn. « Ne t'inquiète pas. »
\item \zh{别哭！} Bié kū ! « Ne pleure pas ! »
\item \zh{别太累了。} Bié tài lèi le. « Ne te fatigue pas trop. »
\end{itemize}
\textbf{Exception / comparaison} : pour les ordres plus secs ou écrits (panneaux, règlements), on utilise \zh{不要} bù yào + verbe : \zh{不要说话。} Bù yào shuōhuà. « Ne parlez pas. » À l'oral familial, 别 est plus naturel et plus doux.
\end{propriete}

\begin{attention}[别 bié, pas 不 bù !]
Erreur n°1 des francophones : traduire « ne t'inquiète pas » par \zh{不担心}. C'est \textbf{faux} : \zh{不} nie les phrases à l'indicatif (\zh{我不担心。} Wǒ bù dānxīn. « Je ne m'inquiète pas. »), jamais l'impératif. Retiens la correspondance :
\begin{itemize}
\item indicatif négatif → \zh{不} : \zh{他不来。} Tā bù lái. « Il ne vient pas. »
\item impératif négatif → \zh{别} : \zh{别来！} Bié lái ! « Ne viens pas ! »
\end{itemize}
De même, « ne t'inquiète pas » se dit \zh{别担心}, jamais \zh{不担心} dans ce sens.
\end{attention}

\begin{propriete}[Proposer son aide : 我来 + verbe]
\textbf{Structure} : (Sujet) + 我来 + verbe \\
\textbf{Sens} : « laisse-moi… », « c'est moi qui m'en charge ». Le mot \zh{来} n'exprime pas ici le mouvement « venir » mais l'\emph{engagement volontaire} du locuteur. \\
\textbf{Nuance} : c'est la formule de la solidarité active — on ne demande pas la permission, on se propose directement, ce qui est très valorisé dans la culture familiale chinoise (et camerounaise !).
\end{propriete}

\begin{exemplebox}[La solidarité en action]
\zh{我来照顾奶奶。} \emph{Wǒ lái zhàogù nǎinai.} « C'est moi qui prends soin de grand-mère. »\\[3pt]
\zh{你先休息，我来做饭。} \emph{Nǐ xiān xiūxi, wǒ lái zuòfàn.} « Repose-toi d'abord, je prépare le repas. »\\[3pt]
\zh{哥哥病了，我来送他去医院。} \emph{Gēge bìng le, wǒ lái sòng tā qù yīyuàn.} « Mon grand frère est malade, c'est moi qui l'accompagne à l'hôpital. »\\[3pt]
\zh{我来洗碗，你去学习吧。} \emph{Wǒ lái xǐ wǎn, nǐ qù xuéxí ba.} « Je fais la vaisselle, va étudier. »
\end{exemplebox}

Note l'emploi de \zh{先} xiān (« d'abord ») placé \emph{avant} le verbe : \zh{你先休息} « toi d'abord reposer » = « repose-toi d'abord ». En français l'adverbe se met après, en chinois toujours avant le verbe.

\courssub{Jeu de rôle 1 : « À l'hôpital »}

\textbf{Situation} : un membre de la famille est malade. La famille s'organise pour l'accompagner à l'hôpital (par exemple l'hôpital central de Yaoundé) et pour répartir les tâches.

\textbf{Cartes de rôle à distribuer} (consignes en français, réplique de départ en chinois) :

\begin{center}
\begin{tabularx}{\linewidth}{|X|X|X|}
\hline
\rowcolor{popblueL} Rôle & Consigne & Réplique de départ \\
\hline
Le père & Tu organises : qui conduit, qui reste à la maison ? & \zh{我们怎么办？} Wǒmen zěnme bàn ? \\
\hline
La mère & Tu rassures le malade et proposes de préparer à manger. & \zh{别担心，我来做饭。} \\
\hline
L'enfant malade & Tu expliques ce que tu as (tête, ventre, fièvre). & \zh{我头疼，还发烧。} Wǒ tóu téng, hái fāshāo. \\
\hline
Le grand frère / la grande sœur & Tu proposes d'accompagner le malade à l'hôpital. & \zh{我来送他去医院。} \\
\hline
La grand-mère & Tu donnes un conseil et bénis la famille. & \zh{我们是一家人！} \\
\hline
\end{tabularx}
\end{center}

\textbf{Vocabulaire utile pour la scène} :

\begin{center}
\begin{tabularx}{\linewidth}{|X|X|X|X|}
\hline
\rowcolor{popblueL} Caractères & Pinyin & Nature & Traduction \\
\hline
医院 & yīyuàn & nom & hôpital \\
\hline
生病 & shēngbìng & verbe & tomber malade \\
\hline
发烧 & fāshāo & verbe & avoir de la fièvre \\
\hline
头疼 & tóu téng & locution & avoir mal à la tête \\
\hline
照顾 & zhàogù & verbe & prendre soin de \\
\hline
送 & sòng & verbe & accompagner, conduire \\
\hline
休息 & xiūxi & verbe & se reposer \\
\hline
药 & yào & nom & médicament \\
\hline
\end{tabularx}
\end{center}

\courssub{Jeu de rôle 2 : « Le repas de famille »}

\textbf{Situation} : dimanche midi, toute la famille se retrouve. Chacun propose d'aider à sa manière. Objectif : utiliser au maximum la structure \zh{我来} + verbe et l'impératif négatif \zh{别}.

\textbf{Réplique vedette} : \zh{我来做饭！} \emph{Wǒ lái zuòfàn !} « Je cuisine ! »

\textbf{Répartition suggérée} : le père achète les provisions (\zh{我来买菜。} Wǒ lái mǎi cài. « Je fais les courses. »), la mère cuisine, le fils met la table (\zh{我来摆桌子。} Wǒ lái bǎi zhuōzi.), la fille fait la vaisselle, le grand-père raconte une histoire pendant le repas. Le malade du premier jeu de rôle peut être invité : \zh{你先休息，我们来做饭。} Nǐ xiān xiūxi, wǒmen lái zuòfàn. « Repose-toi d'abord, nous préparons le repas. »

\begin{savaistu}[La table ronde, symbole de l'unité]
Au Cameroun comme en Chine, les repas de famille sont des moments d'union sacrés. En Chine, on ne sert pas des assiettes individuelles : les plats sont posés \emph{au centre d'une table ronde} et chacun se sert avec ses baguettes. La table ronde (\zh{圆桌} yuánzhuō) symbolise l'unité et la plénitude de la famille, car le mot \zh{圆} yuán (« rond ») évoque la réunion, comme dans \zh{团圆} tuányuán, « réunion familiale ». Au Cameroun, partager un même plat de ndolé ou d'eru traduit exactement la même valeur : manger ensemble, c'est être une famille.
\end{savaistu}

\courssub{Jeu de rôle 3 : « L'interview »}

\textbf{Situation} : un journaliste de la radio scolaire interviewe un élève sur sa famille et la solidarité familiale. Ce jeu de rôle prépare directement l'exposé final de la section 7.

\textbf{Questions du journaliste (à adapter)} :
\begin{itemize}
\item \zh{你家有几口人？} Nǐ jiā yǒu jǐ kǒu rén ? « Combien de personnes compte ta famille ? »
\item \zh{谁照顾奶奶？} Shéi zhàogù nǎinai ? « Qui prend soin de grand-mère ? »
\item \zh{你爱你的家人吗？为什么？} Nǐ ài nǐ de jiārén ma ? Wèishénme ? « Aimes-tu ta famille ? Pourquoi ? »
\item \zh{家人有困难的时候，你怎么办？} Jiārén yǒu kùnnan de shíhou, nǐ zěnme bàn ? « Quand ta famille a des difficultés, que fais-tu ? »
\end{itemize}

\textbf{Amorces de réponse pour l'interviewé} : \zh{我家有五口人……} Wǒ jiā yǒu wǔ kǒu rén… « Ma famille compte cinq personnes… » ; \zh{我很爱我的家人，因为……} Wǒ hěn ài wǒ de jiārén, yīnwèi… « J'aime beaucoup ma famille, parce que… »

\begin{attention}[Le classificateur 口 kǒu pour la famille]
Pour compter les membres d'une famille, on n'utilise pas le classificateur général \zh{个} mais \zh{口} kǒu (littéralement « bouche » : les bouches à nourrir !). On dit \zh{五口人} wǔ kǒu rén « cinq personnes (dans la famille) », jamais \zh{五个人} dans ce contexte précis. En revanche, \zh{五个人} est correct pour cinq personnes quelconques (dans la rue, en classe).
\end{attention}

\courssub{Le cercle des rôles : qui parle à qui ?}

La figure suivante visualise les interactions possibles dans les trois scènes. Chaque flèche représente un échange possible : une question, une proposition d'aide, un conseil.

\begin{popfigure}
\begin{center}
\begin{tikzpicture}[every node/.style={font=\small}]
\node[draw=popblue, fill=popblueL, circle, minimum size=1.7cm, align=center, thick] (j) at (90:3.1) {记者\\ \emph{jìzhě}\\ journaliste};
\node[draw=poporange, fill=poporangeL, circle, minimum size=1.7cm, align=center, thick] (f) at (162:3.1) {爸爸\\ \emph{bàba}\\ père};
\node[draw=poppink, fill=poppinkL, circle, minimum size=1.7cm, align=center, thick] (m) at (234:3.1) {妈妈\\ \emph{māma}\\ mère};
\node[draw=popgreen, fill=popgreenL, circle, minimum size=1.7cm, align=center, thick] (e) at (306:3.1) {孩子\\ \emph{háizi}\\ enfant};
\node[draw=poppurple, fill=poppurpleL, circle, minimum size=1.7cm, align=center, thick] (g) at (18:3.1) {奶奶\\ \emph{nǎinai}\\ grand-mère};
\draw[-{Stealth}, thick, popdark] (j) to[bend left=12] (f);
\draw[-{Stealth}, thick, popdark] (j) to[bend right=12] (e);
\draw[-{Stealth}, thick, popdark] (f) to[bend left=12] (m);
\draw[-{Stealth}, thick, popdark] (m) to[bend left=12] (e);
\draw[-{Stealth}, thick, popdark] (e) to[bend left=12] (g);
\draw[-{Stealth}, thick, popdark] (g) to[bend left=12] (f);
\draw[-{Stealth}, thick, popdark] (m) to[bend left=12] (g);
\node[align=center, font=\small\itshape, text=popmuted] at (0,0) {我来帮你！\\ Wǒ lái bāng nǐ !};
\end{tikzpicture}
\end{center}
\legende{Le cercle des rôles : chaque flèche est un échange possible (question, aide proposée, conseil). Au centre, la formule-clé de la solidarité.}
\end{popfigure}

\courssub{La grille d'observation par les pairs}

Pendant que deux ou trois élèves jouent, les autres observent et remplissent cette grille. Elle rend l'écoute active et donne un retour précis et bienveillant.

\begin{center}
\begin{tabularx}{\linewidth}{|X|X|X|X|}
\hline
\rowcolor{popblueL} Critère & Excellent & Correct & À retravailler \\
\hline
Prononciation des tons & Tons nets, pas de « phrase plate » & Quelques tons confondus (2\textsuperscript{e}/3\textsuperscript{e}) & Tons souvent aplatis \\
\hline
Formules utilisées & 别 + verbe et 我来 + verbe bien employés & Une seule formule utilisée & Formules absentes ou fausses \\
\hline
Fluidité & Parle sans lire, regarde le partenaire & Quelques longues hésitations & Lit ses notes en permanence \\
\hline
Interaction & Relance, réagit, écoute & Répond mais ne relance pas & Attend seulement son tour \\
\hline
\end{tabularx}
\end{center}

\begin{methode}[Donner un bon retour à un camarade]
\begin{enumerate}
\item Commence toujours par un \textbf{point fort} précis : « Ton \zh{别担心} était parfait, les tons étaient nets. »
\item Donne ensuite \textbf{un seul} point à améliorer, formulé positivement : « Essaie de monter davantage le deuxième ton de \zh{来} lái. »
\item Termine par un encouragement en chinois : \zh{很好！加油！} Hěn hǎo ! Jiāyóu ! « Très bien ! Courage ! »
\end{enumerate}
\end{methode}

\begin{exoresolu}[Construire une réplique de solidarité]
\textbf{Énoncé.} Ta petite sœur est malade et s'inquiète de rater ses cours. Construis une réplique complète en deux parties : (a) rassure-la avec l'impératif négatif, (b) propose ton aide avec \zh{我来}.
\tcblower
\textbf{Solution.} Étape 1 : l'impératif négatif exige \zh{别} + verbe : \zh{别担心} Bié dānxīn « ne t'inquiète pas ». Étape 2 : proposition d'aide avec \zh{我来} + verbe ; « t'aider à faire tes devoirs » se dit \zh{帮你做作业} bāng nǐ zuò zuòyè. Réplique complète : \zh{别担心，我来帮你做作业。} \emph{Bié dānxīn, wǒ lái bāng nǐ zuò zuòyè.} « Ne t'inquiète pas, je vais t'aider à faire tes devoirs. » Vérification : \zh{别} et non \zh{不} (impératif) ; \zh{帮} + personne + verbe, structure correcte.
\end{exoresolu}

\begin{exercice}[À toi de jouer]
\begin{enumerate}
\item Traduis en chinois : « Ne pleure pas, je t'accompagne à l'hôpital. »
\item Complète avec \zh{别} ou \zh{不} : (a) \zh{他今天＿来学校。} (b) \zh{＿太累了，先休息吧。}
\item Prépare une carte de rôle « grand-père » pour le jeu de rôle 2 : écris deux répliques avec \zh{我来}.
\item En binôme, jouez la scène 1 pendant 3 minutes ; la classe remplit la grille d'observation, puis \textbf{tournez les rôles} : chacun doit avoir parlé au moins deux fois à la fin de la séance.
\end{enumerate}
\end{exercice}

\begin{corrige}[Exercice « À toi de jouer »]
1. \zh{别哭，我来送你去医院。} Bié kū, wǒ lái sòng nǐ qù yīyuàn. (Impératif négatif avec \zh{别} ; proposition d'aide avec \zh{我来}.)\\
2. (a) \zh{不} : phrase à l'indicatif — \zh{他今天不来学校。} « Il ne vient pas à l'école aujourd'hui. » (b) \zh{别} : impératif — \zh{别太累了，先休息吧。} « Ne te fatigue pas trop, repose-toi d'abord. »\\
3. Exemple : \zh{我来买菜。} Wǒ lái mǎi cài. « Je fais les courses. » / \zh{我来给大家讲故事。} Wǒ lái gěi dàjiā jiǎng gùshi. « Je raconte une histoire à tout le monde. »\\
4. Évaluation par les pairs avec la grille ; l'enseignant veille à la rotation et à ce que chaque élève ait tenu au moins deux rôles différents.
\end{corrige}

\begin{aretenir}[L'essentiel des jeux de rôle]
\begin{itemize}
\item Protocole : 2 min de préparation, 3 min de jeu, 2 min de retour ; rotation obligatoire des rôles.
\item Impératif négatif : \zh{别} + verbe (\zh{别担心}), jamais \zh{不}.
\item Proposer son aide : \zh{我来} + verbe (\zh{我来做饭！}).
\item Compter la famille : \zh{口} kǒu (\zh{五口人}).
\item Trois répliques à connaître par cœur : \zh{你怎么了？} / \zh{别担心，我来帮你。} / \zh{我们是一家人！}
\end{itemize}
\end{aretenir}

\coursec{Débat guidé : la famille d'abord ?}

Après les jeux de rôle qui ont permis de réactiver le lexique de la famille et les structures d'opinion simples (\emph{jeux de rôle guidés}), nous passons maintenant à une activité plus ouverte : un débat structuré. L'objectif n'est pas seulement de « parler », mais de \textbf{construire un argumentaire cohérent} en chinois, en utilisant le connecteur de cause \cle{因为} (yīnwèi) et les formules de désaccord poli. Ce débat prépare directement l'exposé final de la leçon.

\begin{aretenir}[Vocabulaire clé du débat (à maîtriser)]
\begin{center}
\begin{tabularx}{\linewidth}{|X|X|X|X|}
\hline
\rowcolor{popblueL} \textbf{Caractères} & \textbf{Pinyin} & \textbf{Nature} & \textbf{Traduction} \\
\hline
家人 & jiārén & nom & famille, membres de la famille \\
重要 & zhòngyào & adj & important \\
因为 & yīnwèi & conj & parce que, car \\
所以 & suǒyǐ & adv & donc, par conséquent \\
爱 & ài & nom/verbe & amour, aimer \\
帮助 & bāngzhù & verbe/nom & aider, aide \\
团结 & tuánjié & verbe/adj & unir, solidarité, uni \\
朋友 & péngyou & nom & ami(s) \\
支持 & zhīchí & verbe & soutenir, appuyer \\
理解 & lǐjiě & verbe & comprendre (profondément) \\
选择 & xuǎnzé & verbe & choisir \\
比 & bǐ & prép & (comparatif) que \\
更 & gèng & adv & encore plus \\
比如 & bǐrú & conj & par exemple \\
例如 & lìrú & conj & par exemple \\
同意 & tóngyì & verbe & être d'accord \\
不同意 & bù tóngyì & verbe & ne pas être d'accord \\
为什么 & wèishénme & pron. inter. & pourquoi \\
总之 & zǒngzhī & adv & en résumé, bref \\
孝顺 & xiàoshùn & adj/verbe & filial, respecter ses parents \\
根本 & gēnběn & nom & fondement, base \\
友谊 & yǐyì & nom & amitié \\
社会生活 & shèhuì shēnghuó & nom & vie sociale \\
主持人 & zhǔchírén & nom & modérateur, animateur \\
讨论 & tǎolùn & verbe/nom & discuter, discussion \\
根基 & gēnjī & nom & base, fondation (figuré) \\
陪伴 & péibàn & verbe/nom & accompagner, compagnie \\
缺一不可 & quē yī bù kě & chengyu & indispensable, aucun ne peut manquer \\
异乡 & yìxiāng & nom & lieu étranger, loin de chez soi \\
人脉关系 & rénmài guānxì & nom & réseau de relations \\
职业机会 & zhíyè jīhuì & nom & opportunité professionnelle \\
结合 & jiéhé & verbe & combiner, unir \\
成功 & chénggōng & verbe/adj & réussir, succès \\
互助 & hùzhù & verbe/nom & s'entraider, entraide \\
交学费 & jiāo xuéfèi & v+obj & payer les frais de scolarité \\
力量 & lìliàng & nom & force, puissance \\
懂 & dǒng & verbe & comprendre, savoir \\
骂 & mà & verbe & gronder, insulter \\
反而 & fǎn'ér & adv & au contraire \\
鼓励 & gǔlì & verbe & encourager \\
热爱 & rè'ài & verbe & aimer passionnément \\
足球 & zúqiú & nom & football \\
父母 & fùmǔ & nom & parents \\
同学 & tóngxué & nom & camarade de classe \\
练习口语 & liànxí kǒuyǔ & v+obj & pratiquer l'oral \\
永远 & yǒngyuǎn & adv & toujours, à jamais \\
\hline
\end{tabularx}
\end{center}
\end{aretenir}

\courssub{Présentation du sujet et mise en place}

\zh{家人最重要吗？} (Jiārén zuì zhòngyào ma ?) — « La famille est-elle le plus important ? » C'est une question fermée qui appelle un choix, mais la richesse du débat réside dans la justification. La classe est divisée en deux groupes : le groupe \textbf{POUR} (la famille avant tout) et le groupe \textbf{CONTRE} (les amis comptent aussi / la famille n'est pas le \emph{seul} élément important).

\begin{popfigure}
\begin{tikzpicture}[
    node distance=0.5cm and 1cm,
    box/.style={draw=popblue, thick, rounded corners, fill=popblueL, text width=6.5cm, align=center, minimum height=1.2cm},
    title/.style={draw=popdark, thick, rounded corners, fill=popgold, text width=6.5cm, align=center, font=\bfseries},
    arrow/.style={-Stealth, thick, poporange}
]
% Titres
\node[title] (pour_t) {POUR : 家人第一};
\node[title, right=of pour_t] (contre_t) {CONTRE : 朋友也重要};

% Arguments POUR
\node[box, below=of pour_t, align=center] (p1){爱 \textbf{Ài} (Amour)\\\emph{La famille donne un amour inconditionnel.}};
\node[box, below=of p1, align=center] (p2){帮助 \textbf{Bāngzhù} (Aide)\\\emph{Soutien matériel et moral garanti.}};
\node[box, below=of p2, align=center] (p3){团结 \textbf{Tuánjié} (Union)\\\emph{Force collective face aux épreuves.}};

% Arguments CONTRE
\node[box, below=of contre_t, align=center] (c1){朋友 \textbf{Péngyou} (Amis)\\\emph{Choisis, ils partagent nos valeurs.}};
\node[box, below=of c1, align=center] (c2){支持 \textbf{Zhīchí} (Soutien)\\\emph{Les amis soutiennent aussi (études, travail).}};
\node[box, below=of c2, align=center] (c3){平衡 \textbf{Pínghéng} (Équilibre)\\\emph{Vie sociale = famille + amis.}};

% Flèches décoratives
\draw[arrow] (pour_t.south) -- (p1.north);
\draw[arrow] (contre_t.south) -- (c1.north);
\end{tikzpicture}
\legende{Tableau des arguments imposés pour le débat : colonne POUR (famille) vs colonne CONTRE (amis/équilibre).}
\end{popfigure}

\courssub{Observation : le connecteur de cause \texorpdfstring{\zh{因为}}{yīnwèi}}

Avant de lancer le chronomètre, observons comment justifier son avis en chinois. En français, on dit « Je pense que... \textbf{parce que}... ». En chinois, la structure standard place la cause \textbf{après} l'opinion, introduite par \zh{因为} (yīnwèi).

\begin{textebox}[Modèles d'arguments pour le débat]
\textbf{Groupe POUR :} \\
\zh{我觉得家人最重要，因为他们给我们爱、帮助和团结。} \\
\emph{Wǒ juéde jiārén zuì zhòngyào, yīnwèi tāmen gěi wǒmen ài, bāngzhù hé tuánjié.} \\
« Je pense que la famille est la plus importante, \textbf{parce qu'}elle nous donne amour, aide et union. » \\[0.5em]

\textbf{Groupe CONTRE :} \\
\zh{我不同意。因为朋友也很重要，他们也帮助我们。} \\
\emph{Wǒ bù tóngyì. Yīnwèi péngyou yě hěn zhòngyào, tāmen yě bāngzhù wǒmen.} \\
« Je ne suis pas d'accord. \textbf{Parce que} les amis sont aussi importants, ils nous aident aussi. »
\end{textebox}

\begin{propriete}[Structure \texorpdfstring{\zh{因为}}{yīnwèi} + Cause]
\zh{因为} (yīnwèi) signifie « parce que / car ». Il introduit la \textbf{cause} ou la \textbf{raison}.
\medskip

\textbf{Schéma :} \quad \cle{Sujet + 觉得 / 认为 + Opinion ， 因为 + Raison 。}

\begin{center}
\begin{tabularx}{\linewidth}{|X|X|X|}
\hline
\rowcolor{popblueL} \textbf{Structure} & \textbf{Exemple (Caractères + Pinyin)} & \textbf{Traduction} \\
\hline
Avis + \zh{因为} + Cause & \zh{我爱家人，因为他们帮助我。} \newline \emph{Wǒ ài jiārén, yīnwèi tāmen bāngzhù wǒ.} & J'aime ma famille \textbf{parce qu'}ils m'aident. \\
\hline
Avis négatif + \zh{因为} + Cause & \zh{我不觉得朋友次要，因为他们支持我。} \newline \emph{Wǒ bù juéde péngyou cìyào, yīnwèi tāmen zhīchí wǒ.} & Je ne trouve pas les amis secondaires \textbf{car} ils me soutiennent. \\
\hline
Réponse adverse + \zh{因为} & \zh{我不同意，因为朋友也给爱。} \newline \emph{Wǒ bù tóngyì, yīnwèi péngyou yě gěi ài.} & Je ne suis pas d'accord, \textbf{parce que} les amis donnent aussi de l'amour. \\
\hline
\end{tabularx}
\end{center}

\textbf{Point de grammaire clé :} Contrairement au français où « \emph{Parce que}..., \emph{donc}... » est une faute de syntaxe (pléonasme), le chinois standard \textbf{accepte et utilise souvent la corrélation} \zh{因为... 所以...} (yīnwèi... suǒyǐ... : « Parce que... donc... »).
\emph{Exemple :} \zh{因为下雨了，所以我不去。} (Yīnwèi xiàyǔ le, suǒyǐ wǒ bù qù.) — « Parce qu'il pleut, donc je n'y vais pas. » $\rightarrow$ Correct en chinois.
\emph{À l'oral rapide (débat) :} \zh{因为} seul suffit amplement. \zh{所以} est souvent omis pour la fluidité.
\end{propriete}

\begin{attention}[Piège francophone : \texorpdfstring{\zh{因为... 所以...}}{yīnwèi... suǒyǐ...} vs « Parce que... donc... »]
\begin{itemize}
    \item \textbf{Erreur fréquente :} Ne pas traduire mot à mot « Parce que... donc... » en pensant que c'est une erreur en chinois. \textbf{C'est correct en chinois.} La structure \zh{因为... 所以...} est la norme écrite et formelle.
    \item \textbf{À l'oral (débat) :} Privilégiez \zh{因为} seul + phrase principale. Ex: \zh{我觉得家人重要，因为他们爱我。} (Pas besoin de \zh{所以}).
    \item \textbf{Ordre des mots :} En français, la cause peut être en début de phrase (« Parce qu'il pleut, je reste »). En chinois, la proposition en \zh{因为} suit généralement l'opinion principale (Sujet + Verbe d'opinion + Objet, \zh{因为}...). Mettre \zh{因为} en tout début (\zh{因为他们爱我，所以我觉得家人重要}) est possible mais plus formel / écrit.
\end{itemize}
\end{attention}

\courssub{Méthode : Structurer son argument oral en 3 temps}

Pour réussir l'exercice de « 1 minute d'avis personnel », il ne suffit pas de dire « Je pense que... ». Il faut \textbf{argumenter}. Voici la méthode standard pour un exposé oral court (niveau HSK 3/4).

\begin{methode}[Construire un argument oral solide : Avis $\rightarrow$ Raison $\rightarrow$ Exemple]
\underline{Étape 1 : L'Avis (开门见山 -- Kāiménjiànshān : aller droit au but)}
Utilisez \zh{我觉得} (Wǒ juéde - Je sens/pense que) ou \zh{我认为} (Wǒ rènwéi - Je considère que / plus formel).
\emph{Modèle :} \zh{我觉得家人最重要。} / \zh{我认为朋友也很重要。}

\underline{Étape 2 : La Raison (因为 -- Yīnwèi)}
Enchaînez immédiatement avec \zh{因为} + cause. Utilisez le vocabulaire imposé : \zh{爱, 帮助, 团结, 朋友, 支持}.
\emph{Modèle :} \zh{因为他们给我们爱和帮助。} / \zh{因为朋友也支持我们。}

\underline{Étape 3 : L'Exemple concret (例如 / 比如 -- Lìrú / Bǐrú)}
C'est ce qui fait la différence au Bac. Donnez un exemple personnel, court et vrai.
\emph{Mots-outils :} \zh{例如} (lìrú - par exemple), \zh{比如} (bǐrú - par exemple / comme), \zh{有一次} (yǒu yī cì - une fois).
\emph{Modèle :} \zh{比如，上次我生病了，妈妈照顾我。} (Bǐrú, shàngcì wǒ shēngbìng le, māma zhàogù wǒ. -- Par exemple, la dernière fois que j'ai été malade, maman s'est occupée de moi.)

\underline{Étape 4 (Optionnel mais recommandé) : Conclusion / Transition}
\zh{所以，家人最重要。} (Suǒyǐ, jiārén zuì zhòngyào.) ou \zh{因此，我们要孝顺父母。} (Yīncǐ, wǒmen yào xiàoshùn fùmǔ.)
\end{methode}

\begin{exemplebox}[Application de la méthode : Deux mini-exposés types (1 minute)]
\textbf{Élève A (Camp POUR) :} \\
\zh{我觉得家人最重要。} (Avis) \\
\zh{因为家人给我们无条件的爱、帮助和团结。} (Raison : vocabulaire imposé) \\
\zh{比如，我考试不及格的时候，爸爸妈妈没有骂我，反而鼓励我。} (Exemple concret : \zh{骂} mà = gronder, \zh{反而} fǎn'ér = au contraire, \zh{鼓励} gǔlì = encourager) \\
\zh{所以，我永远爱我的家人。} (Conclusion) \\[0.5em]

\textbf{Élève B (Camp CONTRE) :} \\
\zh{我认为朋友也很重要，不只是家人。} (Avis nuancé) \\
\zh{因为朋友是在学校和工作中帮助我们、支持我们的人。} (Raison : \zh{支持} zhīchí) \\
\zh{例如，我学中文很难，是我的同学天天跟我练习口语。} (Exemple : \zh{同学} tóngxué = camarade de classe, \zh{练习口语} liànxí kǒuyǔ = pratiquer l'oral) \\
\zh{所以，家人和朋友都很重要。} (Conclusion commune anticipée)
\end{exemplebox}

\courssub{Conduite du débat : Règles et chronométrage}

L'enseignant joue le rôle de modérateur (\zh{主持人} zhǔchírén). Le débat se déroule en trois phases strictement chronométrées pour forcer la concision et la clarté (compétence Bac).

\begin{center}
\begin{tabularx}{\linewidth}{|X|X|X|}
\hline
\rowcolor{popblueL} \textbf{Phase} & \textbf{Durée} & \textbf{Consigne linguistique} \\
\hline
1. Exposé initial (Tour de table) & 1 min / élève & \zh{我觉得 / 我认为... 因为... 比如...} \\
\hline
2. Réplique / Contre-argument & 30 sec / élève & \zh{我不同意，因为... / 我不太同意，可是...} \\
\hline
3. Synthèse & 1 min (groupe) & \zh{总之... 家人和朋友都很重要。} \\
\hline
\end{tabularx}
\end{center}

\begin{aretenir}[Formules fonctionnelles pour le débat (Fiche de révision orale)]
\begin{itemize}
    \item \textbf{Donner son avis :} \zh{我觉得...} (Wǒ juéde...) / \zh{我认为...} (Wǒ rènwéi...) / \zh{依我看来...} (Yī wǒ kànlái... - À mon avis).
    \item \textbf{Justifier :} \zh{因为...} (Yīnwèi...) / \zh{原因是...} (Yuányīn shì... - La raison est que...).
    \item \textbf{Donner un exemple :} \zh{比如...} (Bǐrú...) / \zh{例如...} (Lìrú...).
    \item \textbf{Exprimer un accord :} \zh{我同意。} (Wǒ tóngyì.) / \zh{对，我也觉得...} (Duì, wǒ yě juéde...).
    \item \textbf{Exprimer un désaccord poli :} \zh{我不同意，因为...} (Wǒ bù tóngyì, yīnwèi...) / \zh{我不太同意，我觉得...} (Wǒ bù tài tóngyì, wǒ juéde... - Je ne suis pas tout à fait d'accord, je pense que...).
    \item \textbf{Relancer / Demander précision :} \zh{为什么？} (Wèishénme ?) / \zh{能不能举个例子？} (Néng bù néng jǔ gè lìzi ? - Peux-tu donner un exemple ?).
    \item \textbf{Conclure :} \zh{总之...} (Zǒngzhī... - En résumé / Bref...) / \zh{所以...} (Suǒyǐ... - Donc...).
\end{itemize}
\end{aretenir}

\courssub{Exercice résolu : Transformer une opinion française en argument chinois structuré}

\begin{exoresolu}[De l'idée à l'oral chinois]
\textbf{Consigne :} Un élève francophone veut dire : « Moi, je pense que les amis sont plus importants que la famille parce qu'on choisit ses amis, et ils nous comprennent mieux. Par exemple, mes parents ne comprennent pas ma passion pour le foot, mais mes potes, si. Propose un script oral chinois (caractères + pinyin) utilisable en 1 minute. »

\tcblower

\textbf{Solution détaillée (étape par étape) :}

\textbf{1. Analyse lexicale nécessaire :}
\begin{itemize}
    \item « Choisir » : \zh{选择} (xuǎnzé).
    \item « Comprendre » : \zh{理解} (lǐjiě).
    \item « Passion » : \zh{热爱} (rè'ài) ou \zh{喜欢} (xǐhuan).
    \item « Foot » : \zh{足球} (zúqiú).
    \item « Potes / Copains » : \zh{朋友} (péngyou) / \zh{哥们儿} (gēmenr - familier nord, éviter au Bac) $\rightarrow$ gardons \zh{朋友}.
    \item « Plus important que » : \zh{比...重要} (bǐ... zhòngyào).
\end{itemize}

\textbf{2. Construction selon la méthode (Avis $\rightarrow$ Raison $\rightarrow$ Exemple) :}

\begin{center}
\begin{tabularx}{\linewidth}{|X|X|X|}
\hline
\rowcolor{popgreenL} \textbf{Étape} & \textbf{Chinois (Caractères)} & \textbf{Pinyin} \\
\hline
Avis & \zh{我觉得朋友比家人更重要。} & Wǒ juéde péngyou bǐ jiārén gèng zhòngyào. \\
\hline
Raison & \zh{因为朋友是我们自己选择的，他们更理解我们。} & Yīnwèi péngyou shì wǒmen zìjǐ xuǎnzé de, tāmen gèng lǐjiě wǒmen. \\
\hline
Exemple & \zh{比如，我热爱足球，父母不太理解，但我的朋友懂我，常跟我踢球。} & Bǐrú, wǒ rè'ài zúqiú, fùmǔ bù tài lǐjiě, dàn wǒ de péngyou dǒng wǒ, cháng gēn wǒ tīqiú. \\
\hline
Conclusion & \zh{所以，朋友在我生命里很重要。} & Suǒyǐ, péngyou zài wǒ shēngmìng lǐ hěn zhòngyào. \\
\hline
\end{tabularx}
\end{center}

\textbf{3. Version fluide pour l'oral (sans pauses tableaux) :}
\zh{我觉得朋友比家人更重要，因为朋友是我们自己选择的，他们更理解我们。比如，我热爱足球，父母不太理解，但我的朋友懂我，常跟我踢球。所以，朋友在我生命里很重要。}

\textbf{4. Points de vigilance pour l'élève :}
\begin{itemize}
    \item \zh{比} (bǐ) structure : \zh{A 比 B Adj} (A est plus Adj que B). Ne pas dire \zh{更重要比家人}.
    \item \zh{理解} (lǐjiě) vs \zh{懂} (dǒng) : \zh{理解} = comprendre profondément (sentiments, raisons) ; \zh{懂} = savoir, connaître, comprendre (langue, fait). Ici \zh{懂我} (me comprendre/me cerner) est très naturel à l'oral.
    \item \zh{常} (cháng) = souvent (plus court que \zh{经常} jīngcháng, idéal pour le débit oral).
\end{itemize}
\end{exoresolu}

\courssub{Le vote final et la conclusion commune : Au-delà de l'opposition}

Le but pédagogique n'est pas de désigner un « gagnant », mais de faire émerger la synthèse culturelle : en Chine comme au Cameroun, la famille est le socle (\zh{根本} gēnběn), mais l'amitié (\zh{友谊} yǐyì) est le ciment de la vie sociale (\zh{社会生活} shèhuì shēnghuó).

\begin{experience}[Déroulé de la conclusion collective (5 min)]
\begin{enumerate}
    \item \textbf{Vote à main levée :} \zh{谁觉得家人最重要？举手。} (Shéi juéde jiārén zuì zhòngyào? Jǔshǒu.) / \zh{谁觉得朋友也很重要？} (Shéi juéde péngyou yě hěn zhòngyào?)
    \item \textbf{Formulation de la synthèse par l'enseignant (modèle) :} \\
    \zh{同学们，经过讨论，我们大家都同意：} (Tóngxuémen, jīngguò tǎolùn, wǒmen dàjiā dōu tóngyì:) \\
    \zh{家人给我们生命、爱和根基；朋友给我们陪伴、理解和支持。} \\
    (Jiārén gěi wǒmen shēngmìng, ài hé gēnjī; péngyou gěi wǒmen péibàn, lǐjiě hé zhīchí.) \\
    \zh{所以，家人和朋友都很重要，缺一不可。} (Suǒyǐ, jiārén hé péngyou dōu hěn zhòngyào, quē yī bù kě.)
    \item \textbf{Répétition chorale de la phrase finale :} \zh{家人和朋友都很重要，缺一不可。} (Jiārén hé péngyou dōu hěn zhòngyào, quē yī bù kě.) — « Famille et amis sont tous importants, aucun ne peut manquer. »
\end{enumerate}
\end{experience}

\begin{savaistu}[Le mot \texorpdfstring{\zh{支持}}{zhīchí} : votre allié pour le Bac]
Le verbe \zh{支持} (zhīchí) signifie « soutenir, appuyer, soutenir financièrement/moralement ». C'est un mot « à haut rendement » pour l'expression écrite et orale du Baccalauréat (niveau HSK 4).
\begin{itemize}
    \item \textbf{Structure :} \zh{Subject + 支持 + Object / Verbe Phrase}.
    \item \textbf{Exemples Bac :}
    \begin{itemize}
        \item \zh{父母支持我学中文。} (Fùmǔ zhīchí wǒ xué Zhōngwén.) — Mes parents soutiennent mon apprentissage du chinois.
        \item \zh{我支持你的决定。} (Wǒ zhīchí nǐ de juédìng.) — Je soutiens ta décision.
        \item \zh{国家支持教育事业。} (Guójiā zhīchí jiàoyù shìyè.) — L'État soutient la cause de l'éducation.
    \end{itemize}
    \item \textbf{Différence \zh{帮助} (bāngzhù) vs \zh{支持} (zhīchí) :} \zh{帮助} = aide concrète, action ponctuelle (aider à porter, aider à faire les devoirs). \zh{支持} = appui moral, financier, politique, ou accord sur un choix de vie (soutenir un projet, une carrière, un avis). \textbf{Au débat :} « Les amis nous \zh{支持} (soutiennent nos choix) » est plus précis que « Les amis nous \zh{帮助} ».
\end{itemize}
\end{savaistu}

\courssub{Entraînement final : Simulation d'exposé noté (Préparation au Bac)}

Pour clore cette section, chaque élève tire au sort une carte « Rôle » et doit présenter un exposé de \textbf{1 minute 30 secondes} devant la classe (ou en binôme enregistré sur téléphone pour auto-évaluation). La grille d'évaluation porte sur :
\begin{enumerate}
    \item \textbf{Prononciation des tons} (notamment \zh{因为 yīnwèi}, \zh{重要 zhòngyào}, \zh{朋友 péngyou}, \zh{支持 zhīchí}).
    \item \textbf{Structure} : Avis clair + \zh{因为} + Raison + \zh{比如} + Exemple.
    \item \textbf{Richesse lexicale} : Utilisation d'au moins 3 mots de la liste imposée (\zh{爱, 帮助, 团结, 朋友, 支持, 理解, 选择}).
    \item \textbf{Cohérence} : L'exemple illustre bien la raison.
\end{enumerate}

\begin{exercice}[Simulation d'exposé : Carte « Pour » / Carte « Contre »]
\textbf{Carte A (Pour) :} « Vous êtes l'aîné(e) d'une famille de 5 enfants à Yaoundé. Expliquez pourquoi la solidarité familiale (\zh{团结}) a permis à vos frères/sœurs de réussir leurs études. Utilisez \zh{因为} et \zh{比如}. »

\textbf{Carte B (Contre/Nuancé) :} « Vous êtes étudiant(e) à Douala, loin de votre village natal. Vos amis de fac (\zh{同学}) sont devenus votre "deuxième famille". Expliquez pourquoi ce soutien (\zh{支持}) est vital pour vous. Utilisez \zh{因为} et \zh{例如}. »

\textbf{Carte C (Synthèse - niveau excellence) :} « Comparez le rôle du \zh{叔叔阿姨} (oncles/tantes - au sens large camerounais) et celui des \zh{同窗} (camarades de promo / condisciples) dans la réussite sociale d'un jeune Camerounais aujourd'hui. Concluez par \zh{家人和朋友都很重要，缺一不可。} »
\end{exercice}

\begin{corrige}[Exercice : Simulation d'exposé - Pistes de correction pour l'enseignant]
\textbf{Carte A (Modèle) :} \zh{我觉得家人最重要，因为家人团结互助。比如，我大哥毕业后，给弟弟妹妹交学费，全家一起努力，现在我们都上大学了。所以，家人的团结是力量。} (Tons à surveiller : \zh{团结 tuánjié}, \zh{互助 hùzhù}, \zh{交学费 jiāo xuéfèi}, \zh{力量 lìliàng}).

\textbf{Carte B (Modèle) :} \zh{我认为朋友也很重要，因为朋友在异乡支持我们。例如，我刚到杜阿拉不认识路，同学带我吃饭、找房子，像哥哥姐姐一样。所以，朋友就是第二个家。} (Tons : \zh{异乡 yìxiāng}, \zh{带我 dài wǒ}, \zh{找房子 zhǎo fángzi}, \zh{哥哥姐姐 gēge jiějie}).

\textbf{Carte C (Modèle) :} \zh{喀麦隆的叔叔阿姨给我们传统教育和人脉关系，同窗给我们现代知识和职业机会。两者结合，才能成功。家人和朋友都很重要，缺一不可。} (Vocabulaire avancé : \zh{人脉关系 rénmài guānxì}, \zh{职业机会 zhíyè jīhuì}, \zh{结合 jiéhé}, \zh{成功 chénggōng}).
\end{corrige}

\courssub{Bilan de la section : Ce qu'il faut retenir pour l'exposé final}

\begin{aretenir}[Check-list « Débat \& Exposé » -- À mémoriser]
\begin{enumerate}
    \item \textbf{Le connecteur roi :} \zh{因为} (yīnwèi) introduit la cause. Structure : \zh{Opinion, 因为 Raison.}
    \item \textbf{La structure en 3 temps :} \zh{我觉得... 因为... 比如...} (Avis -- Raison -- Exemple).
    \item \textbf{Vocabulaire imposé actif :} \zh{爱 ài, 帮助 bāngzhù, 团结 tuánjié, 朋友 péngyou, 支持 zhīchí, 理解 lǐjiě, 选择 xuǎnzé}.
    \item \textbf{Formules de politesse débat :} \zh{我不同意，因为...} (Désaccord justifié) / \zh{我赞同，而且...} (Accord + ajout).
    \item \textbf{Phrase de synthèse culturelle :} \zh{家人和朋友都很重要，缺一不可。} (Jiārén hé péngyou dōu hěn zhòngyào, quē yī bù kě.)
    \item \textbf{Tons critiques à réviser :} \zh{重} (zhòng - 4ème ton) dans \zh{重要}, \zh{为} (wèi - 4ème ton) dans \zh{因为}, \zh{持} (chí - 2ème ton) dans \zh{支持}.
\end{enumerate}
\end{aretenir}

Cette section termine la partie « Expression orale » du Module 1. Les élèves disposent désormais des outils linguistiques (vocabulaire, grammaire \zh{因为}, structure argumentative) et socioculturels (valeurs famille/amis Cameroun-Chine) pour aborder sereinement l'\textbf{Expression écrite} (rédaction d'un petit texte d'opinion) et la \textbf{Révision générale} qui suivront.

\coursec{Bilan et entraînement à l'exposé final}

Après le débat guidé « La famille d'abord ? », nous consolidons maintenant tout ce qui a été appris dans cette leçon : vocabulaire de la famille, structures pour exprimer l'opinion et la solidarité, formules fonctionnelles pour structurer un exposé, et surtout la prononciation correcte des tons. Cette dernière section vous prépare à l'examen oral : vous devez être capable de présenter un exposé court, clair, bien prononcé et structuré sur votre famille et l'amour familial.

\courssub{Synthèse orale collective : les 5 étapes de l'exposé, les formules, les tons difficiles}

\begin{experience}[Synthèse collective — Rappel des acquis]
L'enseignant anime une mise en commun rapide (5 min) au tableau. Les élèves, par groupes de 4, rappellent :
\begin{enumerate}
\item Les \textbf{5 étapes} obligatoires d'un exposé oral réussi.
\item Les \textbf{formules fonctionnelles} pour chaque étape (salutation, introduction, développement, exemples, conclusion).
\item Les \textbf{tons difficiles} identifiés dans les dialogues et le débat (ex. : 3e ton + 3e ton, 4e ton suivi de 4e ton, ton neutre sur les particules).
\end{enumerate}
Un secrétaire par groupe note au tableau ; la classe valide ou corrige collectivement.
\end{experience}

\begin{popfigure}
\begin{tikzpicture}[
    node distance=1.2cm and 0.8cm,
    stepnode/.style={rectangle, rounded corners, draw=popblue, thick, fill=popblueL, minimum width=2.8cm, minimum height=1.2cm, align=center, font=\small},
    arrownode/.style={-Stealth, thick, popdark}
]
\node[stepnode, align=center] (e1){问好\\Wèn hǎo\\Saluer};
\node[stepnode, right=of e1, align=center] (e2){介绍主题\\Jièshào zhǔtí\\Annoncer le sujet};
\node[stepnode, right=of e2, align=center] (e3){说看法\\Shuō kànfǎ\\Exprimer son avis};
\node[stepnode, right=of e3, align=center] (e4){举例子\\Jǔ lìzi\\Donner des exemples};
\node[stepnode, right=of e4, align=center] (e5){谢谢\\Xièxie\\Remercier \& conclure};

\draw[arrownode] (e1.east) -- (e2.west);
\draw[arrownode] (e2.east) -- (e3.west);
\draw[arrownode] (e3.east) -- (e4.west);
\draw[arrownode] (e4.east) -- (e5.west);

\node[below=0.3cm of e1, font=\tiny, popmuted] {大家好！};
\node[below=0.3cm of e2, font=\tiny, popmuted] {今天我要谈一谈…};
\node[below=0.3cm of e3, font=\tiny, popmuted] {我觉得…因为…};
\node[below=0.3cm of e4, font=\tiny, popmuted] {比如…例如…};
\node[below=0.3cm of e5, font=\tiny, popmuted] {我的话说完了，谢谢大家！};
\end{tikzpicture}
\legende{Frise des 5 étapes obligatoires de l'exposé oral (HSK 2-3)}
\end{popfigure}

\begin{aretenir}[Les 5 étapes de l'exposé oral type]
\begin{enumerate}
\item \textbf{问好 (Wèn hǎo)} — Saluer l'auditoire : \zh{大家好！} (Dàjiā hǎo !) ou \zh{老师好，同学们好！} (Lǎoshī hǎo, tóngxuémen hǎo !).
\item \textbf{介绍主题 (Jièshào zhǔtí)} — Annoncer le sujet : \zh{今天我要谈一谈我的家人。} (Jīntiān wǒ yào tántan wǒ de jiārén.) / \zh{我的题目是：家庭的爱与团结。} (Wǒ de tímù shì: Jiātíng de ài yǔ tuánjié.)
\item \textbf{说看法 (Shuō kànfǎ)} — Exprimer son opinion avec justification : \zh{我觉得家人最重要，因为他们总是支持我。} (Wǒ juéde jiārén zuì zhòngyào, yīnwèi tāmen zǒngshì zhīchí wǒ.)
\item \textbf{举例子 (Jǔ lìzi)} — Illustrer par un exemple concret : \zh{比如，上个星期我生病了，妈妈照顾我。} (Bǐrú, shàng gè xīngqī wǒ shēngbìng le, māma zhàogù wǒ.)
\item \textbf{谢谢 (Xièxie)} — Conclure poliment : \zh{我的话说完了，谢谢大家！} (Wǒ de huà shuō wán le, xièxie dàjiā !)
\end{enumerate}
\end{aretenir}

\begin{propriete}[Règle de changement de ton : le 3e ton (Sandhi)]
En chinois standard, le 3e ton (ton bas descendant-remontant \emph{ˇ}) change de réalisation selon le ton qui suit :
\begin{itemize}
\item \textbf{3e ton + 3e ton} : le premier 3e ton devient un \textbf{2e ton} (montant \emph{´}). Ex. : \zh{大家好} → \emph{Dájiā hǎo} ; \zh{很好} → \emph{hén hǎo}.
\item \textbf{3e ton + 1er, 2e, 4e ton ou ton neutre} : le premier 3e ton se réalise en \textbf{demi-3e ton} (bas, sans remontée). Ex. : \zh{因为} (yīnwèi, 1+4 inchangé) ; \zh{比如} (bǐrú, 3+2 → \emph{bírú} car le 2e ton « tire » le 3e vers le haut).
\item \textbf{3e ton isolé ou en fin de phrase} : ton complet (descendant-remontant).
\end{itemize}
\emph{Attention :} Le pinyin écrit garde le trait du 3e ton (ǎ), mais l'oral suit la règle ci-dessus.
\end{propriete}

\begin{center}
\begin{tabularx}{\linewidth}{|X|X|X|X|}
\hline
\rowcolor{popblueL} \textbf{Étape} & \textbf{Formule clé (caractères + pinyin)} & \textbf{Traduction} & \textbf{Point de ton à surveiller} \\
\hline
1. Saluer & \zh{大家好！} Dàjiā hǎo ! & Bonjour à tous ! & 3e + 3e → \textbf{2e + 3e} : \emph{Dájiā hǎo} \\
\hline
2. Sujet & \zh{今天我要谈一谈…} Jīntiān wǒ yào tántan… & Aujourd'hui je vais parler de… & \zh{谈一谈} tántan (2e + ton neutre) léger \\
\hline
3. Avis & \zh{我觉得…因为…} Wǒ juéde… yīnwèi… & Je pense que… parce que… & \zh{觉得} juéde (2e + ton neutre) ; \zh{因为} yīnwèi (1e + 4e) \\
\hline
4. Exemple & \zh{比如… / 例如…} Bǐrú… / Lìrú… & Par exemple… / Par exemple… & \zh{比如} bǐrú (3e + 2e) → \emph{bírú} \\
\hline
5. Conclure & \zh{我的话说完了，谢谢大家！} Wǒ de huà shuō wán le, xièxie dàjiā ! & J'ai fini, merci à tous ! & \zh{说完了} shuō wán le (1e + 2e + neutre) ; \zh{谢谢} xièxie (4e + ton neutre) \\
\hline
\end{tabularx}
\end{center}

\courssub{Lexique clé de la leçon : famille, solidarité, connecteurs}

\begin{center}
\begin{tabularx}{\linewidth}{|X|X|X|X|}
\hline
\rowcolor{popblueL} \textbf{Caractères} & \textbf{Pinyin} & \textbf{Nature} & \textbf{Traduction} \\
\hline
家人 & jiārén & nom & famille, membres de la famille \\
\hline
口 & kǒu & classifieur & classifieur pour les personnes (famille) \\
\hline
团结 & tuánjié & adj./verbe & uni, solidaire / s'unir \\
\hline
支持 & zhīchí & verbe & soutenir \\
\hline
鼓励 & gǔlì & verbe & encourager \\
\hline
照顾 & zhàogù & verbe & prendre soin de, s'occuper de \\
\hline
陪 & péi & verbe & accompagner, tenir compagnie \\
\hline
亲戚 & qīnqi & nom & parents, proches (famille élargie) \\
\hline
长辈 & zhǎngbèi & nom & aînés, générations supérieures \\
\hline
晚辈 & wǎnbèi & nom & cadets, générations inférieures \\
\hline
爱 & ài & verbe & aimer \\
\hline
重要 & zhòngyào & adj. & important \\
\hline
因为 & yīnwèi & conjonction & parce que \\
\hline
所以 & suǒyǐ & conjonction & donc, c'est pourquoi \\
\hline
比如 & bǐrú & conjonction & par exemple (oral) \\
\hline
例如 & lìrú & conjonction & par exemple (écrit/formel) \\
\hline
但是 & dànshì & conjonction & mais \\
\hline
大家 & dàjiā & pronom & tout le monde \\
\hline
\end{tabularx}
\end{center}

\courssub{Plan type de l'exposé final}

\begin{textebox}[Modèle d'exposé complet — Version standard]
大家好！我叫玛丽。今天我要谈一谈我的家人。我家有五口人：爸爸、妈妈、哥哥、妹妹和我。我很爱我的家人。我觉得家人最重要，因为他们总是帮助我、支持我。比如，上个星期我考试没考好，爸爸鼓励我，妈妈给我做好吃的。我们一家人很团结。我的话说完了，谢谢大家！
\par\medskip
Dàjiā hǎo ! Wǒ jiào Mǎlì. Jīntiān wǒ yào tántan wǒ de jiārén. Wǒ jiā yǒu wǔ kǒu rén: bàba, māma, gēge, mèimei hé wǒ. Wǒ hěn ài wǒ de jiārén. Wǒ juéde jiārén zuì zhòngyào, yīnwèi tāmen zǒngshì bāngzhù wǒ, zhīchí wǒ. Bǐrú, shàng gè xīngqī wǒ kǎoshì méi kǎohǎo, bàba gǔlì wǒ, māma gěi wǒ zuò hǎochī de. Wǒmen yī jiā rén hěn tuánjié. Wǒ de huà shuō wán le, xièxie dàjiā !
\par\medskip
Bonjour à tous ! Je m'appelle Marie. Aujourd'hui je vais parler de ma famille. Ma famille compte cinq personnes : papa, maman, grand frère, petite sœur et moi. J'aime beaucoup ma famille. Je pense que la famille est la plus importante, car ils m'aident et me soutiennent toujours. Par exemple, la semaine dernière j'ai raté mon examen, papa m'a encouragée, maman m'a cuisiné de bons plats. Notre famille est très unie. J'ai terminé mon propos, merci à tous !
\end{textebox}

\begin{textebox}[Modèle d'exposé — Contexte camerounais]
大家好！我叫保罗，我住在雅温得。今天我要谈一谈家庭的爱与团结。我家有六口人：爸爸、妈妈、爷爷、奶奶、姐姐和我。我们三代同堂，很团结。我觉得长辈很辛苦，所以我们要孝顺他们。比如，上个星期天，全家人一起包饺子，奶奶教我，爸爸洗菜，妈妈做饭，非常开心。我的话说完了，谢谢大家！
\par\medskip
Dàjiā hǎo ! Wǒ jiào Bǎoluó, wǒ zhù zài Yǎwēndé. Jīntiān wǒ yào tántan jiātíng de ài yǔ tuánjié. Wǒ jiā yǒu liù kǒu rén: bàba, māma, yéye, nǎinai, jiějie hé wǒ. Wǒmen sān dài tóngtáng, hěn tuánjié. Wǒ juéde zhǎngbèi hěn xīnkǔ, suǒyǐ wǒmen yào xiàoshùn tāmen. Bǐrú, shàng gè xīngqītiān, quán jiā rén yìqǐ bāo jiǎozi, nǎinai jiāo wǒ, bàba xǐ cài, māma zuò fàn, fēicháng kāixīn. Wǒ de huà shuō wán le, xièxie dàjiā !
\par\medskip
Bonjour à tous ! Je m'appelle Paul, j'habite à Yaoundé. Aujourd'hui je vais parler de l'amour et de la solidarité familiale. Ma famille compte six personnes : papa, maman, grand-père, grand-mère, grande sœur et moi. Nous vivons trois générations sous le même toit, très unis. Je trouve que les aînés ont une vie dure, donc nous devons leur être pieux (filiaux). Par exemple, dimanche dernier, toute la famille a fait des raviolis ensemble, grand-mère m'a appris, papa a lavé les légumes, maman a cuisiné, c'était très joyeux. J'ai terminé, merci à tous !
\end{textebox}

\begin{methode}[Préparer un PLAN (mots-clés), pas un texte appris par cœur]
\begin{enumerate}
\item \textbf{Lisez les modèles} ci-dessus et repérez la structure en 5 étapes.
\item \textbf{Écrivez seulement des mots-clés} sur une fiche (format carte de visite) : 
\begin{itemize}
\item Salutation : \zh{大家好}
\item Sujet : \zh{我的家人} / \zh{家庭的爱}
\item Avis + raison : \zh{很重要} + \zh{因为支持我} / \zh{长辈辛苦} + \zh{所以孝顺}
\item Exemple concret : \zh{生病 / 考试 / 生日 / 周末包饺子} + \zh{谁} + \zh{做了什么}
\item Conclusion : \zh{我的话说完了，谢谢大家！}
\end{itemize}
\item \textbf{Entraînez-vous à voix haute} 3 fois en ne regardant que vos mots-clés. Vous parlerez plus naturellement, avec vos propres phrases, et vous gérerez mieux les tons.
\item \textbf{Chronométrez-vous} : visez 1 min 30 à 2 min maximum.
\end{enumerate}
\end{methode}

\courssub{Préparation individuelle de 5 minutes : rédiger un plan de 6 à 8 phrases}

\begin{experience}[Préparation individuelle — 5 minutes]
Chaque élève prend une feuille blanche et rédige \textbf{son plan personnel} (pas des phrases complètes, seulement des mots-clés et connecteurs) en suivant le tableau ci-dessous. L'enseignant circule, aide à corriger les tons sur les mots-clés difficiles.
\end{experience}

\begin{center}
\begin{tabularx}{\linewidth}{|X|X|X|}
\hline
\rowcolor{popblueL} \textbf{Étape} & \textbf{Mots-clés à écrire (caractères + pinyin)} & \textbf{Votre contenu personnel} \\
\hline
1. 问好 & \zh{大家好！} Dàjiā hǎo ! & (identique pour tous) \\
\hline
2. 介绍主题 & \zh{今天我要谈一谈…} Jīntiān wǒ yào tántan… & \zh{我的家人} / \zh{家庭的爱与团结} … \\
\hline
3. 说看法 & \zh{我觉得…因为… / 所以…} Wǒ juéde… yīnwèi… / suǒyǐ… & \zh{家人很重要} / \zh{我们很团结} + \zh{因为…} / \zh{所以…} \\
\hline
4. 举例子 & \zh{比如… / 例如…} Bǐrú… / Lìrú… & \zh{上个星期 / 昨天 / 周末} + \zh{什么事} + \zh{谁} + \zh{做了什么} \\
\hline
5. 谢谢 & \zh{我的话说完了，谢谢大家！} Wǒ de huà shuō wán le, xièxie dàjiā ! & (identique pour tous) \\
\hline
\end{tabularx}
\end{center}

\begin{attention}[Ne pas traduire mot à mot du français — Pièges fréquents]
\begin{itemize}
\item \textbf{« Ma famille me manque »} $\neq$ \zh{我的家人缺我} (incorrect, \zh{缺} « manquer à qqn » s'emploie autrement). \\
$\rightarrow$ Dire : \zh{我想我的家人。} (Wǒ xiǎng wǒ de jiārén.) \emph{« Je pense à ma famille / Ma famille me manque. »}
\item \textbf{« Je m'occupe de ma grand-mère »} $\neq$ \zh{我管我的奶奶} (\zh{管} « gérer, contrôler » → trop autoritaire). \\
$\rightarrow$ Dire : \zh{我照顾奶奶。} (Wǒ zhàogù nǎinai.) ou \zh{我陪奶奶。} (Wǒ péi nǎinai. « Je tiens compagnie à grand-mère. »)
\item \textbf{« Nous sommes unis »} $\neq$ \zh{我们是联合} (\zh{联合} nom/verbe « union/alliance »). \\
$\rightarrow$ Dire : \zh{我们很团结。} (Wǒmen hěn tuánjié.) adjectif d'état.
\item \textbf{Classifieur familial} : \zh{五口人} (wǔ kǒu rén) pour « 5 personnes (famille) », pas \zh{五个人} dans le contexte familial formel.
\item \textbf{Connecteur \zh{所以} (suǒyǐ)} : En oral, \zh{因为…所以…} structure la pensée. Ne l'oubliez pas si la phrase est longue : \zh{我觉得家人重要，因为他们支持我，所以我很爱他们。}
\end{itemize}
\end{attention}

\courssub{Atelier prononciation : tons pièges de l'exposé}

\begin{experience}[Entraînement ciblé sur les tons — 5 minutes]
En binôme, lisez à voix haute les séries ci-dessous. L'écouteur note les erreurs de sandhi (3+3, 3+2, ton neutre). Changez de rôle.
\begin{enumerate}
\item \zh{大家好、很好、马马虎虎、姐姐、妹妹、想一想} (Focus : 3e + 3e → 2e + 3e)
\item \zh{比如、例如、因为、所以、鼓励、支持、照顾、团结} (Focus : 3e + 2e → 2e + 2e ; tons neutres sur \zh{鼓励 gǔli, 支持 zhīchi, 照顾 zhàogu})
\item \zh{谢谢、对不起、没关系、好的、是的} (Focus : ton neutre sur 2e syllabe)
\item \zh{我的话说完了、考试没考好、生日快乐、一家人} (Focus : enchainement naturel, \zh{了} le neutre)
\end{enumerate}
\end{experience}

\courssub{Passage au tableau de volontaires avec grille d'évaluation}

\begin{experience}[Passage à l'oral — 20 à 25 minutes]
\begin{enumerate}
\item L'enseignant appelle 4 à 6 volontaires (selon le temps). Chaque passage dure \textbf{1 min 30 max}.
\item Les autres élèves ont la \textbf{grille d'évaluation} (distribuée ou projetée) et notent leur camarade \textbf{bienveillamment}.
\item Après chaque passage : 30 secondes de retour oral de l'enseignant (1 point fort + 1 axe de progression).
\end{enumerate}
\end{experience}

\begin{center}
\begin{tabularx}{\linewidth}{|X|X|X|}
\hline
\rowcolor{popblueL} \textbf{Critère} & \textbf{Points} & \textbf{Descripteurs (niveau HSK 2-3 / Première A)} \\
\hline
\textbf{Tons (prononciation)} & 3 pts & 3 = tons corrects sur \textgreater 90\% des syllabes, y compris changements de ton (3+3, 3+2, 4+4, neutre) ; 2 = quelques erreurs mais compréhensible ; 1 = erreurs fréquentes gênant la compréhension. \\
\hline
\textbf{Formules fonctionnelles} & 3 pts & 3 = 5 étapes présentes, connecteurs (\zh{因为/所以/比如/但是}) utilisés correctement ; 2 = 4 étapes, quelques connecteurs ; 1 = structure incomplète. \\
\hline
\textbf{Fluidité / Débit} & 2 pts & 2 = débit naturel, peu d'hésitations, auto-correction possible ; 1 = nombreuses pauses, lecture du plan. \\
\hline
\textbf{Contenu / Vocabulaire} & 2 pts & 2 = vocabulaire de la leçon employé (\zh{亲戚, 团结, 支持, 鼓励, 照顾, 陪, 长辈, 晚辈}…), exemple personnel pertinent ; 1 = vocabulaire minimal, exemple vague. \\
\hline
\textbf{TOTAL} & \textbf{/10} & \textbf{Seuil de réussite : 6/10} \\
\hline
\end{tabularx}
\end{center}

\begin{exoresolu}[Simulation d'évaluation — Exposé d'un élève]
\textbf{Énoncé :} Un élève présente l'exposé suivant (enregistré). Évaluez-le avec la grille.
\par\medskip
\zh{大家好！我叫保罗。今天我要谈一谈我的家人。我家有四口人：爸爸、妈妈、姐姐和我。我觉得家人很重要，因为他们爱我。比如，我的生日，妈妈做蛋糕，姐姐给我礼物。我们很团结。我的话说完了，谢谢大家！}
\par\medskip
Dàjiā hǎo ! Wǒ jiào Bǎoluó. Jīntiān wǒ yào tántan wǒ de jiārén. Wǒ jiā yǒu sì kǒu rén: bàba, māma, jiějie hé wǒ. Wǒ juéde jiārén hěn zhòngyào, yīnwèi tāmen ài wǒ. Bǐrú, wǒ de shēngrì, māma zuò dàngāo, jiějie gěi wǒ lǐwù. Wǒmen hěn tuánjié. Wǒ de huà shuō wán le, xièxie dàjiā !
\par\medskip
\textbf{Solution détaillée :}
\begin{itemize}
\item \textbf{Tons (3/3)} : \zh{保罗} Bǎoluó (3e+2e) correct ; \zh{四口人} sì kǒu rén (4e+3e+2e) correct ; \zh{因为} yīnwèi (1e+4e) correct ; \zh{比如} bǐrú → \emph{bírú} (changement 3+2) correct ; \zh{团结} tuánjié (2e+2e) correct ; \zh{蛋糕} dàngāo (4e+1e) correct. Pas d'erreur majeure.
\item \textbf{Formules (3/3)} : 5 étapes respectées ; connecteurs \zh{因为}, \zh{比如} bien placés.
\item \textbf{Fluidité (2/2)} : Débit régulier, pas de lecture, hésitation légère sur \zh{蛋糕} dàngāo (acceptable).
\item \textbf{Contenu (2/2)} : Vocabulaire ciblé (\zh{四口人, 重要, 爱, 生日, 做蛋糕, 礼物, 团结}) ; exemple personnel concret.
\end{itemize}
\textbf{Note finale : 10/10 — Excellent.}
\end{exoresolu}

\courssub{Auto-évaluation : chaque élève coche les objectifs atteints de la leçon}

\begin{experience}[Auto-évaluation individuelle — 3 minutes]
Chaque élève reçoit la fiche ci-dessous (ou la recopie dans son cahier) et coche honnêtement. L'enseignant récupère pour adapter la révision du module.
\end{experience}

\begin{center}
\begin{tabularx}{\linewidth}{|X|c|c|c|}
\hline
\rowcolor{popblueL} \textbf{Objectif de la leçon (ZHA03)} & \textbf{Acquis} & \textbf{En cours} & \textbf{Non acquis} \\
\hline
Je sais saluer et annoncer un sujet oralement (\zh{大家好，今天我要谈一谈…}) & $\square$ & $\square$ & $\square$ \\
\hline
Je sais exprimer mon avis avec \zh{我觉得…因为… / 所以…} & $\square$ & $\square$ & $\square$ \\
\hline
Je sais donner un exemple avec \zh{比如… / 例如…} & $\square$ & $\square$ & $\square$ \\
\hline
Je sais conclure avec \zh{我的话说完了，谢谢大家！} & $\square$ & $\square$ & $\square$ \\
\hline
Je prononce correctement les tons difficiles : \zh{大家好} (3+3→2+3), \zh{比如} (3+2→2+2), \zh{谢谢} (4e+neutre), \zh{鼓励/支持/照顾} (ton neutre) & $\square$ & $\square$ & $\square$ \\
\hline
J'utilise le vocabulaire de la solidarité : \zh{团结, 支持, 鼓励, 照顾, 陪, 亲戚, 长辈, 晚辈} & $\square$ & $\square$ & $\square$ \\
\hline
Je participe à un débat en donnant mon accord / désaccord : \zh{我同意 / 我不太同意 / 我觉得不一样} & $\square$ & $\square$ & $\square$ \\
\hline
Je peux présenter un exposé de 1 min 30 sur ma famille sans lire & $\square$ & $\square$ & $\square$ \\
\hline
\end{tabularx}
\end{center}

\courssub{Devoir : enregistrer son exposé à la maison et le rejouer en classe}

\begin{methode}[Devoir maison — Enregistrement audio]
\begin{enumerate}
\item \textbf{Enregistrez-vous} (smartphone, dictaphone, ou application WeChat/Voice Memo) en présentant votre exposé final (1 min 30 à 2 min) \textbf{sans lire}, uniquement avec votre plan mots-clés.
\item \textbf{Réécoutez-vous} : cochez sur votre grille d'auto-évaluation ce qui va / ce qui ne va pas (tons, fluidité, formules).
\item \textbf{Ré-enregistrez} une deuxième version corrigée.
\item \textbf{Apportez le fichier} (ou envoyez-le à l'enseignant via la plateforme de l'établissement) pour la prochaine séance. En classe, nous en écouterons 3-4 anonymement pour une analyse collective bienveillante.
\end{enumerate}
\end{methode}

\begin{savaistu}[À l'épreuve orale de chinois (Baccalauréat / HSKK)]
L'examinateur évalue \textbf{d'abord la prononciation des tons et l'aisance (fluency)}. Mieux vaut des phrases simples, courtes, bien prononcées, avec les bonnes formules de liaison (\zh{因为、所以、比如、但是}), que des phrases longues, complexes, mais hésitantes et truffées d'erreurs de tons. \textbf{La régularité du débit et la justesse tonale comptent pour 50\% de la note orale.} Entraînez-vous chaque jour 5 minutes à lire un court texte à voix haute en exagérant les tons : c'est le secret des meilleurs élèves.
\end{savaistu}

\begin{exercice}[Entraînement final — Complétez l'exposé]
Complétez les parties manquantes (entre crochets) avec les formules appropriées, puis lisez l'exposé complet à voix haute en respectant les tons.
\par\medskip
\zh{\_\_\_\_\_\_\_\_\_\_！我叫\_\_\_\_\_\_\_\_\_\_。\_\_\_\_\_\_\_\_\_\_我的家人。我家有\_\_\_\_\_\_\_\_\_\_口人：\_\_\_\_\_\_\_\_\_\_。我很\_\_\_\_\_\_\_\_\_\_我的家人。\_\_\_\_\_\_\_\_\_\_家人最重要，\_\_\_\_\_\_\_\_\_\_他们总是\_\_\_\_\_\_\_\_\_\_我、\_\_\_\_\_\_\_\_\_\_我。\_\_\_\_\_\_\_\_\_\_，\_\_\_\_\_\_\_\_\_\_。我们一家人很\_\_\_\_\_\_\_\_\_\_。\_\_\_\_\_\_\_\_\_\_，谢谢大家！}
\par\medskip
\textit{Indice : 5 étapes, 8-9 phrases. Utilisez : 大家好, 今天我要谈一谈, 六, 爸爸妈妈哥哥姐姐和我, 爱, 我觉得, 因为, 支持, 鼓励, 比如, 生日妈妈做面条, 团结, 我的话说完了.}
\end{exercice}

\begin{corrige}[Exercice — Entraînement final]
\zh{大家好！我叫保罗。今天我要谈一谈我的家人。我家有六口人：爸爸、妈妈、哥哥、姐姐和我。我很爱我的家人。我觉得家人最重要，因为他们总是支持我、鼓励我。比如，我的生日，妈妈做面条。我们一家人很团结。我的话说完了，谢谢大家！}
\par\medskip
Dàjiā hǎo ! Wǒ jiào Bǎoluó. Jīntiān wǒ yào tántan wǒ de jiārén. Wǒ jiā yǒu liù kǒu rén: bàba, māma, gēge, jiějie hé wǒ. Wǒ hěn ài wǒ de jiārén. Wǒ juéde jiārén zuì zhòngyào, yīnwèi tāmen zǒngshì zhīchí wǒ, gǔlì wǒ. Bǐrú, wǒ de shēngrì, māma zuò miàntiáo. Wǒmen yī jiā rén hěn tuánjié. Wǒ de huà shuō wán le, xièxie dàjiā !
\par\medskip
Traduction : Bonjour à tous ! Je m'appelle Paul. Aujourd'hui je vais parler de ma famille. Ma famille compte six personnes : papa, maman, grand frère, grande sœur et moi. J'aime beaucoup ma famille. Je pense que la famille est la plus importante, car ils me soutiennent et m'encouragent toujours. Par exemple, pour mon anniversaire, maman fait des nouilles. Notre famille est très unie. J'ai terminé, merci à tous !
\end{corrige}

\begin{aretenir}[Check-list finale avant l'oral]
\begin{itemize}
\item[$\square$] Mon plan tient sur une carte (mots-clés seulement).
\item[$\square$] Je connais par cœur la formule de conclusion : \zh{我的话说完了，谢谢大家！}
\item[$\square$] J'ai répété 3 fois à voix haute en me chronométrant (1 min 30 - 2 min).
\item[$\square$] J'ai vérifié les tons pièges : \zh{大家好} (2+3), \zh{比如} (2+2), \zh{因为} (1+4), \zh{谢谢} (4+neutre), \zh{鼓励/支持/照顾} (ton neutre 2e syllabe).
\item[$\square$] J'ai un exemple personnel concret (date, personne, action).
\item[$\square$] Je ne lis pas, je parle à l'auditoire (regard, sourire).
\end{itemize}
\end{aretenir}

\textbf{Félicitations !} Vous avez maintenant tous les outils pour réussir votre exposé oral sur l'amour familial et la solidarité. La clé : \textbf{structure claire + tons justes + exemple vivant = réussite}. Bon entraînement et à la prochaine leçon pour la révision générale du Module 1 !

\newpage

\coursec{Activité d'intégration}

\coursec{Leçon 3 — Expression orale : Amour de la famille et solidarité familiale}

\courssub{Objectifs d'apprentissage}
\begin{aretenir}[Compétences visées (Programme MINESEC 1ère A — Module 1)]
À l'issue de cette leçon (2 h), l'élève doit être capable de :
\begin{itemize}
    \item \textbf{Comprendre} des dialogues et courts exposés sur la famille et la solidarité (vocabulaire HSK 2-3).
    \item \textbf{Produire} un exposé oral structuré (1-2 min) selon un plan en 5 étapes : salutation, annonce du sujet, avis personnel (\zh{觉得}), exemple concret (\zh{例如}), conclusion (\zh{所以}) et remerciement.
    \item \textbf{Interagir} dans un débat guidé : poser des questions de relance (\zh{你呢？}, \zh{你觉得怎么样？}), répondre en justifiant (\zh{因为}...\zh{所以}...).
    \item \textbf{Prononcer} correctement les tons, en gérant le sandhi du 3e ton (3e + 3e $\to$ 2e + 3e) et l'intonation phrase (déclarative, interrogative, exclamative).
    \item \textbf{Mobiliser} le lexique (famille, sentiments, solidarité) et les formules fonctionnelles (donner son avis, relancer, conclure) dans une situation authentique (Forum de la jeunesse).
    \item \textbf{Comparer} les valeurs familiales Cameroun/Chine (respect des aînés, famille élargie, solidarité concrète).
\end{itemize}
\end{aretenir}

\courssub{Mise en route : Parler de sa famille}
\begin{experience}[Activité de démarrage — « Arbre généalogique éclair » (10 min)]
\begin{enumerate}
    \item \textbf{Individuel (2 min) :} Dessinez rapidement votre arbre familial sur 3 générations (grands-parents, parents, fratrie, vous). Notez les termes chinois à côté de chaque personne (\zh{爷爷、奶奶、外公、外婆、爸爸、妈妈、哥哥、姐姐、弟弟、妹妹、我}).
    \item \textbf{Binôme (5 min) :} Présentez votre famille à votre voisin en 5 phrases minimum. Utilisez : \zh{我家有...口人。} \zh{我有...} \zh{我爱我的...} \zh{我们很团结。}
    \item \textbf{Mise en commun (3 min) :} 3 élèves volontaires présentent la famille de leur camarade à la classe (\zh{他/她家有...}).
\end{enumerate}
\end{experience}

\begin{attention}[Rappel classificateurs famille]
\zh{口} (kǒu) pour le nombre de personnes dans le foyer (\zh{三口人}). \zh{个} (gè) pour les membres pris individuellement (\zh{两个哥哥}). \zh{位} (wèi) forme polie pour les aînés (\zh{两位爷爷}).
\end{attention}


\courssub{Dialogues modèles}
\begin{textebox}[Dialogue 1 : Discussion entre camarades — « Qui t'a le plus marqué ? »]
\textbf{Contexte :} Deux élèves de 1ère A, \zh{小林} (Xiǎo Lín) et \zh{阿明} (Ā Míng), parlent dans la cour du Lycée de la Réunification.
\zh{小林：阿明，这个周末你过得怎么样？} \\
Xiǎo Lín: Ā Míng, zhège zhōumò nǐ guò de zěnmeyàng? \\
(Xiao Lin : Ah Ming, comment s'est passé ton week-end ?) \\[4pt]
\zh{阿明：很好！我回家看望爷爷奶奶了。他们很高兴。} \\
Ā Míng: Hěn hǎo! Wǒ huíjiā kànwàng yéye nǎinai le. Tāmen hěn gāoxìng. \\
(Ah Ming : Très bien ! Je suis rentré voir grand-père et grand-mère. Ils étaient très contents.) \\[4pt]
\zh{小林：真好。我觉得家庭很重要，大家在一起最开心。} \\
Xiǎo Lín: Zhēn hǎo. Wǒ juéde jiātíng hěn zhòngyào, dàjiā zài yīqǐ zuì kāixīn. \\
(Xiao Lin : C'est bien. Je trouve que la famille est importante, être ensemble est le plus joyeux.) \\[4pt]
\zh{阿明：我也是。因为家人互相爱护，所以我们有力量。你呢？} \\
Ā Míng: Wǒ yě shì. Yīnwèi jiārén hùxiāng àihù, suǒyǐ wǒmen yǒu lìliang. Nǐ ne? \\
(Ah Ming : Moi aussi. Parce que la famille s'aime mutuellement, nous avons de la force. Et toi ?) \\[4pt]
\zh{小林：我最喜欢和姐姐在一起，因为她照顾我。谢谢你们！} \\
Xiǎo Lín: Wǒ zuì xǐhuan hé jiějie zài yīqǐ, yīnwèi tā zhàogù wǒ. Xièxie nǐmen! \\
(Xiao Lin : J'aime le plus être avec ma grande sœur, car elle prend soin de moi. Merci !)
\end{textebox}

\begin{textebox}[Dialogue 2 : Mini-exposé improvisé — « Mon exemple de solidarité »]
\textbf{Contexte :} L'enseignant demande un exemple concret de solidarité familiale.
\zh{老师：谁能举个例子？} \\
Lǎoshī: Shéi néng jǔ gè lìzi? \\
(Professeur : Qui peut donner un exemple ?) \\[4pt]
\zh{学生：我举个例子。上个星期我考试不好，心情不好。妈妈给我做了好吃的，爸爸陪我聊天，我就好多了。} \\
Xuésheng: Wǒ jǔ gè lìzi. Shàng gè xīngqī wǒ kǎoshì bù hǎo, xīnqíng bù hǎo. Māma gěi wǒ zuò le hǎochī de, bàba péi wǒ liáotiān, wǒ jiù hǎo duō le. \\
(Élève : Je donne un exemple. La semaine dernière, j'ai raté mon examen, j'étais de mauvaise humeur. Maman m'a cuisiné de bons plats, papa a discuté avec moi, et ça m'a beaucoup remonté le moral.) \\[4pt]
\zh{老师：很好。因为家人的照顾，所以你感觉好多了。这就是团结的力量。} \\
Lǎoshī: Hěn hǎo. Yīnwèi jiārén de zhàogù, suǒyǐ nǐ gǎnjué hǎo duō le. Zhè jiù shì tuánjié de lìliang. \\
(Professeur : Très bien. Grâce aux soins de la famille, tu te sens beaucoup mieux. C'est ça la force de la solidarité.)
\end{textebox}

\courssub{Formules fonctionnelles pour exposer et débattre}
\begin{propriete}[Boîte à outils orale — Structures clés (HSK 2-3)]
\begin{center}
\begin{tabularx}{\linewidth}{|X|X|X|}
\hline
\rowcolor{poporangeL} \textbf{Fonction} & \textbf{Formule (Caractères + Pinyin)} & \textbf{Traduction / Usage} \\ \hline
\textbf{Saluer / Ouvrir} & \zh{各位同学，大家好！} (Gèwèi tóngxué, dàjiā hǎo!) & Chers camarades, bonjour à tous ! (Formel) \\ \hline
& \zh{同学们好，老师好。} (Tóngxuémen hǎo, lǎoshī hǎo.) & Bonjour les camarades, bonjour prof. (Standard) \\ \hline
\textbf{Annoncer le sujet} & \zh{今天我的演讲题目是……} (Jīntiān wǒ de yǎnjiǎng tímù shì…) & Aujourd'hui mon sujet d'exposé est… \\ \hline
& \zh{我想谈谈……} (Wǒ xiǎng tántan…) & Je voudrais parler de… \\ \hline
\textbf{Exprimer son avis} & \zh{我觉得……} (Wǒ juéde…) & Je trouve que… / Je pense que… (Subjectif) \\ \hline
& \zh{我认为……} (Wǒ rènwéi…) & Je considère que… (Plus formel) \\ \hline
& \zh{在我看来……} (Zài wǒ kàn lái…) & À mon avis… / Pour moi… \\ \hline
\textbf{Justifier (Cause $\to$ Conséquence)} & \zh{因为……，所以……} (Yīnwèi…, suǒyǐ…) & Parce que…, donc… \textbf{Ordre fixe !} \\ \hline
\textbf{Illustrer (Exemple)} & \zh{例如……} (Lìrú…) & Par exemple… (Écrit / Oral soutenu) \\ \hline
& \zh{比较……} (Bǐjiào…) & Par exemple / Disons… (Oral courant) \\ \hline
& \zh{上个星期……} (Shàng gè xīngqī…) & La semaine dernière… (Marqueur temps passé) \\ \hline
\textbf{Relancer / Interagir} & \zh{你呢？} (Nǐ ne?) & Et toi ? (Rebond sur même thème) \\ \hline
& \zh{你觉得怎么样？} (Nǐ juéde zěnmeyàng?) & Qu'en penses-tu ? / Comment trouves-tu ça ? \\ \hline
& \zh{你同意吗？} (Nǐ tóngyì ma?) & Es-tu d'accord ? \\ \hline
\textbf{Conclure / Remercier} & \zh{所以，我们要……} (Suǒyǐ, wǒmen yào…) & Donc, nous devons… (Engagement) \\ \hline
& \zh{谢谢大家！} (Xièxie dàjiā!) & Merci à tous ! \\ \hline
& \zh{我的演讲完了。} (Wǒ de yǎnjiǎng wán le.) & Mon exposé est fini. \\ \hline
\end{tabularx}
\end{center}
\end{propriete}

\courssub{Les tons à travailler : Le 3e ton en chaîne (Sandhi)}
\begin{methode}[Règle d'or du 3e ton + 3e ton $\to$ 2e ton + 3e ton]
\begin{enumerate}
    \item \textbf{Identifiez} les mots en 3e ton (ton bas-descendant-remontant : \zh{ǎ}, \zh{ě}, \zh{ǐ}, \zh{ǒ}, \zh{ǔ}, \zh{ǚ}).
    \item \textbf{Repérez} les séquences de \textbf{deux 3e tons consécutifs} (même mot ou mots collés).
    \item \textbf{Transformez} le \textbf{premier} 3e ton en \textbf{2e ton} (ton montant : \zh{á}, \zh{é}, \zh{í}, \zh{ó}, \zh{ú}, \zh{ǘ}). Le second reste 3e ton.
    \item \textbf{Entraînez-vous} par paires : exagérer le mouvement montant du 1er, puis bas du 2e.
\end{enumerate}
\end{methode}

\begin{attention}[Pièges de la leçon — Mots du vocabulaire concernés]
\begin{itemize}
    \item \zh{互相} : hǔ + xiāng $\to$ \textbf{hú xiāng} (2 + 1)
    \item \zh{演讲} : yǎn + jiǎng $\to$ \textbf{yán jiǎng} (2 + 3)
    \item \zh{所以} : suǒ + yǐ $\to$ \textbf{suó yǐ} (2 + 3)
    \item \zh{奶奶} : nǎi + nai $\to$ \textbf{nái nai} (2 + neutre) — \textit{2e 3e ton devient neutre souvent}.
    \item \zh{好吃} : hǎo + chī $\to$ \textbf{háo chī} (2 + 1)
    \item \zh{感动} : gǎn + dòng $\to$ \textbf{gǎn dòng} (3 + 4) \textit{\textbf{Pas de changement} car le 2e est 4e ton ! (Demi-3e ton seulement).}
    \item \zh{爱护} : \textbf{ài hù} (4 + 4) \textit{\textbf{Attention !} « 爱 » est 4e ton, pas 3e. Pas de sandhi.}
    \item \zh{照顾} : \textbf{zhào gù} (4 + 4) \textit{\textbf{Attention !} « 照 » est 4e ton. Pas de sandhi.}
    \item \zh{生病} : \textbf{shēng bìng} (1 + 4) \textit{\textbf{Attention !} « 生 » est 1er ton. Pas de sandhi.}
\end{itemize}
\end{attention}


\begin{popfigure}
\begin{tikzpicture}[
    node distance=1.2cm and 1.5cm,
    box/.style={draw=poporange, thick, rounded corners, fill=poporangeL, minimum width=2.5cm, minimum height=1cm, align=center},
    arrow/.style={-Stealth, thick, poporange}
]
\node[box] (t1) {3e ton + 3e ton};
\node[box, right=of t1] (t2) {1er devient 2e ton};
\node[box, right=of t2] (t3) {2e ton + 3e ton};
\draw[arrow] (t1) -- (t2);
\draw[arrow] (t2) -- (t3);
\node[below=0.3cm of t1, text=popdark, font=\small] {\zh{互相} hǔ+xiāng};
\node[below=0.3cm of t3, text=popdark, font=\small] {\zh{互相} \textbf{hú}+xiāng};
\node[below=0.3cm of t2, text=popgreen, font=\small\bfseries] {Règle};
\end{tikzpicture}
\legende{Sandhi du 3e ton : application sur \zh{互相} (hùxiāng $\to$ húxiāng).}
\end{popfigure}

\courssub{Jeux de rôle guidés : Préparation au forum}
\begin{experience}[Atelier « Tribune libre » — Entraînement par paires (20 min)]
\begin{enumerate}
    \item \textbf{Rôles :} A = Orateur (prépare un exposé 5 étapes, 6-8 phrases). B = Journaliste (prépare 2 questions de relance : \zh{你呢？} + \zh{你觉得怎么样？}).
    \item \textbf{Contraintes :} A doit utiliser \textbf{au moins} : 1x \zh{觉得}, 1x \zh{因为}...\zh{所以}..., 1x \zh{例如}, 3 mots du lexique « Famille \& Solidarité ». B doit utiliser \zh{你呢？} et \zh{你觉得怎么样？}.
    \item \textbf{Déroulement :} A parle (notes autorisées). B écoute, note 1 point fort / 1 point à améliorer (tons, structure). B pose ses 2 questions. A répond \textbf{obligatoirement} avec \zh{因为}...\zh{所以}... Inversion des rôles.
    \item \textbf{Retour classe :} 2 binômes volontaires passent au tableau. Classe évalue avec la grille (voir Ressources).
\end{enumerate}
\end{experience}

\courssub{Activité d'intégration : Forum de la jeunesse — « La famille, notre force »}
\begin{textebox}[Contexte : Forum de la jeunesse de Yaoundé — \zh{家庭与团结} (Jiātíng yǔ tuánjié)]
Le Lycée de la Réunification à Yaoundé organise son premier \textbf{Forum de la jeunesse} sur le thème \zh{家庭与团结}. Vous êtes délégué(e) de votre classe. Chaque élève monte à la tribune pour livrer un \textbf{mini-exposé} (6 à 8 phrases, 1-2 min) puis répondre à deux questions de ses camarades. Un jury d'élèves note la prestation sur une grille (tons, formules, fluidité, contenu). La séance se clôture par la récitation collective de la devise : \zh{家人互相帮助，团结就是力量！}
\end{textebox}

\courssub{Consignes de l'activité finale}
\begin{methode}[Déroulement pas à pas (45 min)]
\begin{enumerate}
    \item \textbf{Préparation individuelle (10 min) :} Rédigez votre exposé en suivant \textbf{strictement} le plan en 5 étapes ci-dessous. Utilisez au moins 3 formules fonctionnelles et 5 mots du lexique. Entraînez-vous à voix haute : surlignez les syllabes en 3e ton qui changent (sandhi).
    \begin{itemize}
        \item \textbf{Étape 1 — Salutation :} \zh{各位同学，大家好。} / \zh{同学们好，老师好。}
        \item \textbf{Étape 2 — Annonce du sujet :} \zh{今天我的演讲题目是……} / \zh{我想谈谈……}
        \item \textbf{Étape 3 — Avis personnel :} \zh{我觉得……因为……} / \zh{我认为……}
        \item \textbf{Étape 4 — Exemple personnel (Cameroun/Chine) :} \zh{例如，……} / \zh{上个星期，……}
        \item \textbf{Étape 5 — Conclusion \& Remerciement :} \zh{所以，我们要……谢谢大家！}
    \end{itemize}
    \item \textbf{Passage à la tribune (1-2 min/élève) :} Exposé sans lecture (carte bristol autorisée). Classe en écoute active.
    \item \textbf{Relance interactive (1 min) :} Deux camarades désignés posent :
    \begin{itemize}
        \item Camarade A : \zh{你呢？你最喜欢和谁在一起？}
        \item Camarade B : \zh{你觉得怎么样？家庭重要吗？}
    \end{itemize}
    L'orateur répond \textbf{obligatoirement} avec \zh{因为} + raison + \zh{所以} + conséquence.
    \item \textbf{Évaluation par les pairs :} Grille d'observation (voir Ressources). Titre \zh{最佳演讲者} décerné.
    \item \textbf{Clôture collective :} Toute la classe se lève, récite 3 fois : \zh{家人互相帮助，团结就是力量！}
\end{enumerate}
\end{methode}

\courssub{Ressources à mobiliser}
\begin{definition}[Lexique clé — Famille, sentiments, solidarité (HSK 2-3)]
\begin{center}
\begin{tabularx}{\linewidth}{|X|X|X|X|}
\hline
\rowcolor{popblueL} \textbf{Caractères} & \textbf{Pinyin} & \textbf{Nature} & \textbf{Traduction} \\ \hline
家庭 & jiātíng & n. & famille, foyer \\ \hline
家人 & jiārén & n. & membres de la famille, proches \\ \hline
团结 & tuánjié & v./adj. & s'unir, solidarité / uni, solidaire \\ \hline
力量 & lìliang & n. & force, puissance \\ \hline
互相 & hùxiāng & adv. & mutuellement, l'un l'autre \\ \hline
爱护 & àihù & v. & aimer et protéger, chérir \\ \hline
照顾 & zhàogù & v. & prendre soin de, s'occuper de \\ \hline
感动 & gǎndòng & v./adj. & être ému / touchant, émouvant \\ \hline
演讲 & yǎnjiǎng & v./n. & faire un discours / discours, exposé \\ \hline
题目 & tímù & n. & sujet, titre, thème \\ \hline
重要 & zhòngyào & adj. & important \\ \hline
所以 & suǒyǐ & conj. & donc, par conséquent \\ \hline
大家 & dàjiā & pron. & tout le monde, chacun \\ \hline
孤单 & gūdān & adj. & seul, solitaire \\ \hline
幸福 & xìngfú & adj. & heureux, bonheur \\ \hline
\end{tabularx}
\end{center}
\end{definition}

\begin{propriete}[Structures grammaticales indispensables pour l'exposé et le débat]
\begin{center}
\begin{tabularx}{\linewidth}{|X|X|X|}
\hline
\rowcolor{poporangeL} \textbf{Structure} & \textbf{Exemple (Caractères + Pinyin)} & \textbf{Traduction} \\ \hline
Suj. + \cle{觉得} + [Phrase d'opinion] & \zh{我觉得家庭很重要。} (Wǒ juéde jiātíng hěn zhòngyào.) & Je trouve que la famille est très importante. \\ \hline
\cle{因为} + Raison , \cle{所以} + Conséquence & \zh{因为家人互相爱护，所以我很幸福。} (Yīnwèi jiārén hùxiāng àihù, suǒyǐ wǒ hěn xìngfú.) & Comme la famille s'aime mutuellement, je suis très heureux. \\ \hline
\cle{例如} / \cle{比较} + Exemple concret & \zh{例如，妈妈给我做好吃的。} (Lìrú, māma gěi wǒ zuò hǎochī de.) & Par exemple, maman me cuisine de bons plats. \\ \hline
\cle{你呢？} (Relance elliptique) & \zh{我喜欢和奶奶在一起。你呢？} (Wǒ xǐhuan hé nǎinai zài yīqǐ. Nǐ ne?) & J'aime être avec grand-mère. Et toi ? \\ \hline
\cle{你觉得怎么样？} (Demander avis) & \zh{这个题目很好。你觉得怎么样？} (Zhège tímù hěn hǎo. Nǐ juéde zěnmeyàng?) & Ce sujet est bien. Qu'en penses-tu ? \\ \hline
Suj. + \cle{给} + Bénéficiaire + V & \zh{哥哥给我讲笑话。} (Gēge gěi wǒ jiǎng xiàohuà.) & Grand frère me raconte des blagues. (\zh{给} = pour/moi) \\ \hline
Temps + Suj. + V + \cle{了} (Changement d'état) & \zh{上个星期我生病了。} (Shàng gè xīngqī wǒ shēngbìng le.) & La semaine dernière, je suis tombé malade. \\ \hline
\end{tabularx}
\end{center}
\end{propriete}

\begin{attention}[Pièges fréquents pour francophones — \textbf{À afficher au tableau}]
\begin{itemize}
    \item \textbf{Ordre \zh{因为}...\zh{所以}... :} La cause \textbf{précède toujours} la conséquence. Interdit : *\zh{所以我很高兴，因为...}*.
    \item \textbf{\zh{互相} vs \zh{一起} :} \zh{互相} = action réciproque (l'un l'autre) \textit{Ex: \zh{我们互相帮助}} ; \zh{一起} = action commune (ensemble) \textit{Ex: \zh{我们一起吃饭}}.
    \item \textbf{Sandhi 3e ton :} \zh{互相} (hǔxiāng $\to$ \textbf{húxiāng}), \zh{演讲} (yǎnjiǎng $\to$ \textbf{yánjiǎng}), \zh{所以} (suǒyǐ $\to$ \textbf{suóyǐ}), \zh{好吃} (hǎochī $\to$ \textbf{háochī}). \textbf{Pas de sandhi} sur \zh{爱护} (ài hù, 4+4), \zh{照顾} (zhào gù, 4+4), \zh{生病} (shēng bìng, 1+4), \zh{感动} (gǎn dòng, 3+4).
    \item \textbf{\zh{题目} vs \zh{问题} :} \zh{题目} = sujet/titre (exposé, devoir) ; \zh{问题} = question/problème.
    \item \textbf{\zh{给} position :} Suj + \zh{给} + Bénéf + V (+ Obj). Ne pas mettre \zh{给} après le verbe (*\zh{讲给我笑话} $\times$ $\to$ \zh{给我讲笑话} $\checkmark$).
\end{itemize}
\end{attention}


\begin{methode}[Grille d'évaluation par les pairs (à photocopier / projeter)]
\begin{center}
\begin{tabularx}{\linewidth}{|X|c|c|c|c|}
\hline
\rowcolor{popgreenL} \textbf{Critère} & \textbf{1 - Insuffisant} & \textbf{2 - Passable} & \textbf{3 - Bon} & \textbf{4 - Excellent} \\ \hline
\textbf{Tons \& Prononciation} & Beaucoup d'erreurs, incompréhensible & Erreurs sur tons 2/3, compréhensible & Tons globalement justes, sandhi 3e géré & Tons parfaits, intonation naturelle, fluide \\ \hline
\textbf{Formules fonctionnelles} & Aucune formule apprise & 1 formule correcte & 2-3 formules (\zh{我觉得}, \zh{因为}...) & 3+ formules, transitions naturelles \\ \hline
\textbf{Fluidité / Débit} & Hésitations longues, lecture mot à mot & Hésitations, regards notes fréquents & Débit soutenu, regards public, notes discrètes & Débit naturel, regard soutenu, sans notes \\ \hline
\textbf{Contenu / Structure} & Hors sujet, pas de structure & Structure incomplète (3/5 étapes) & Plan 5 étapes respecté, exemple personnel & Plan respecté, exemple vivant, vocabulaire riche \\ \hline
\end{tabularx}
\end{center}
\end{methode}

\courssub{Production type et corrigé commenté}
\begin{exoresolu}[Modèle d'exposé + Interaction + Analyse linguistique]
\textbf{Énoncé :} Produire un exposé modèle respectant le plan en 5 étapes, suivi de l'interaction type (2 questions / 2 réponses avec \zh{因为}...\zh{所以}...).

\tcblower

\textbf{1. Exposé modèle (Durée : env. 1 min 15 s)}

\begin{textebox}[Exposé de l'élève : « Ma famille, mon refuge » — \zh{我的家，我的港湾} (Wǒ de jiā, wǒ de gǎngwān)]
\textbf{Étape 1 — Salutation :} \\
\zh{各位同学，大家好！} \\
Gèwèi tóngxué, dàjiā hǎo ! \\
(Chers camarades, bonjour à tous !) \\[6pt]

\textbf{Étape 2 — Annonce du sujet :} \\
\zh{今天我的演讲题目是：「家庭，我们的力量」。} \\
Jīntiān wǒ de yǎnjiǎng tímù shì: « Jiātíng, wǒmen de lìliang ». \\
(Aujourd'hui, mon sujet d'exposé est : « La famille, notre force ».) \\[6pt]

\textbf{Étape 3 — Avis personnel :} \\
\zh{我觉得家庭非常重要，因为家人互相爱护，互相帮助。} \\
Wǒ juéde jiātíng fēicháng zhòngyào, yīnwèi jiārén hùxiāng àihù, hùxiāng bāngzhù. \\
(Je trouve que la famille est extrêmement importante, car les proches s'aiment et s'aident mutuellement.) \\[6pt]

\textbf{Étape 4 — Exemple personnel (Contexte camerounais) :} \\
\zh{例如，上个星期我生病了，妈妈照顾我，爸爸买药，哥哥给我讲笑话，我非常感动。} \\
Lìrú, shàng gè xīngqī wǒ shēngbìng le, māma zhàogù wǒ, bàba mǎi yào, gēge gěi wǒ jiǎng xiàohuà, wǒ fēicháng gǎndòng. \\
(Par exemple, la semaine dernière j'ai été malade, maman a pris soin de moi, papa a acheté des médicaments, grand frère m'a raconté des blagues, j'ai été très ému.) \\[6pt]

\textbf{Étape 5 — Conclusion \& Remerciement :} \\
\zh{所以，我们要爱护家人，多回家看看。谢谢大家！} \\
Suǒyǐ, wǒmen yào àihù jiārén, duō huíjiā kànkan. Xièxie dàjiā ! \\
(Donc, nous devons chérir nos proches, rentrer souvent à la maison. Merci à tous !)
\end{textebox}

\textbf{2. Interaction type (Relance \& Réponse avec \zh{因为}...\zh{所以}...)}

\begin{center}
\begin{tabularx}{\linewidth}{|X|X|}
\hline
\rowcolor{poppurpleL} \textbf{Rôle} & \textbf{Échange (Caractères + Pinyin + Traduction)} \\ \hline
\textbf{Camarade A} & \zh{你呢？你最喜欢和谁在一起？} \\
& Nǐ ne? Nǐ zuì xǐhuan hé shéi zài yīqǐ? \\
& (Et toi ? Avec qui aimes-tu être le plus ?) \\ \hline
\textbf{Orateur (Réponse)} & \zh{我最喜欢和奶奶在一起，因为奶奶给我做好吃的，所以我很开心。} \\
& Wǒ zuì xǐhuan hé nǎinai zài yīqǐ, yīnwèi nǎinai gěi wǒ zuò hǎochī de, suǒyǐ wǒ hěn kāixīn. \\
& (J'aime le plus être avec grand-mère, car elle me cuisine de bons plats, donc je suis très content.) \\ \hline
\textbf{Camarade B} & \zh{你觉得怎么样？家庭重要吗？} \\
& Nǐ juéde zěnmeyàng? Jiātíng zhòngyào ma? \\
& (Qu'en penses-tu ? La famille est-elle importante ?) \\ \hline
\textbf{Orateur (Réponse)} & \zh{我觉得非常重要。因为没有家庭，人会孤单，所以家庭是力量。} \\
& Wǒ juéde fēicháng zhòngyào. Yīnwèi méiyǒu jiātíng, rén huì gūdān, suǒyǐ jiātíng shì lìliang. \\
& (Je trouve que c'est extrêmement important. Sans famille, on est seul, donc la famille est la force.) \\ \hline
\end{tabularx}
\end{center}

\textbf{3. Commentaire linguistique détaillé (Pour l'enseignant / Auto-correction élève)}

\begin{itemize}
    \item \textbf{Tons \& Sandhi (3e ton) — \textbf{Points critiques à marquer au crayon} :}
    \begin{itemize}
        \item \zh{互相} (hǔxiāng $\to$ \textbf{húxiāng}) — 3+1 $\to$ 2+1.
        \item \zh{演讲} (yǎnjiǎng $\to$ \textbf{yánjiǎng}) — 3+3 $\to$ 2+3.
        \item \zh{所以} (suǒyǐ $\to$ \textbf{suóyǐ}) — 3+3 $\to$ 2+3.
        \item \zh{好吃} (hǎochī $\to$ \textbf{háochī}) — 3+1 $\to$ 2+1.
        \item \zh{奶奶} (nǎinai $\to$ \textbf{nái nai}) — 3+neutre $\to$ 2+neutre.
        \item \zh{爱护} (\textbf{ài hù}, 4+4) — \textbf{Pas de sandhi} (1er caractère 4e ton).
        \item \zh{照顾} (\textbf{zhào gù}, 4+4) — \textbf{Pas de sandhi}.
        \item \zh{生病} (\textbf{shēng bìng}, 1+4) — \textbf{Pas de sandhi}.
        \item \zh{感动} (\textbf{gǎn dòng}, 3+4) — Demi-3e ton sur \zh{感}, pas de changement en 2e ton.
    \end{itemize}
    \item \textbf{Structure \zh{因为}...\zh{所以}... :} Utilisée 3 fois dans l'exposé + 2 fois dans les réponses. Vérifier : \zh{因为} en tête de proposition (pas après sujet). Ex: \zh{因为家人互相爱护，所以...} $\checkmark$ ; *\zh{家庭因为...所以...}* $\times$ (oral standard).
    \item \textbf{Marqueurs de temps :} \zh{上个星期} placé \textbf{avant} le sujet \zh{我} (Temps + Sujet + V) ou après (Sujet + Temps + V). Les deux corrects HSK 2.
    \item \textbf{Complément \zh{给} (Bénéficiaire) :} \zh{哥哥给我讲笑话} / \zh{奶奶给我做好吃的}. Structure : Suj + \zh{给} + Bénéf + V + Obj. Point HSK 3 clé.
    \item \textbf{Particule \zh{了} (Accompli / Changement d'état) :} \zh{生病了} (tomber malade / être malade maintenant), \zh{做了} (action passée accomplie). Marque l'actualisation.
    \item \textbf{Relance \zh{你呢？} :} Particule modale \zh{ね} (ne) pour rebondir sur le prédicat précédent (\zh{和谁在一起}). Économie de langue.
    \item \textbf{Vocabulaire bonus (niveau +) :} \zh{港湾} (gǎngwān - port d'attache/refuge, métaphore famille), \zh{避风港} (bìfēnggǎng - abri), \zh{孝顺} (xiàoshùn - piété filiale, valeur clé Chine).
\end{itemize}
\end{exoresolu}

\courssub{Débat guidé : La famille d'abord ?}
\begin{experience}[Débat structuré — « Travail / Études loin de la famille : comment garder le lien ? » (15 min)]
\begin{enumerate}
    \item \textbf{Lancer le thème :} \zh{现在很多年轻人去大城市（杜阿拉、雅温得、北京、上海）工作或上学，远离家庭。怎么保持团结？} (Aujourd'hui beaucoup de jeunes partent en ville... Comment garder la solidarité ?)
    \item \textbf{Groupes de 4 :} Chaque groupe a 3 cartes « Argument » (ex: \zh{视频通话}, \zh{寄钱/红包}, \zh{回家过节}, \zh{每周打电话}). Ils préparent 3 arguments avec \zh{因为}...\zh{所以}...
    \item \textbf{Tour de table :} Un rapporteur par groupe présente : \zh{我们组觉得……因为……所以……} (Notre groupe pense que... car... donc...).
    \item \textbf{Synthèse prof :} Relever les structures, corriger les tons, valider le vocabulaire (\zh{视频通话} shìpín tōnghuà, \zh{红包} hóngbāo, \zh{过节} guòjié).
    \item \textbf{Ouverture culturelle :} \zh{中国：春节必回家 (Chūnjié bì huíjiā) — La fête du Printemps, retour obligatoire.} vs \zh{喀麦隆：节日全家聚餐 (Jiérì quánjiā jùcān) — Fêtes = repas familial élargi.}
\end{enumerate}
\end{experience}

\courssub{Bilan et entraînement à l'exposé final}
\begin{aretenir}[Bilan — Forum de la jeunesse : \zh{家庭与团结}]
\begin{itemize}
    \item \textbf{Compétence orale validée :} Discours continu structuré (monologue planifié 5 étapes) + Échange interactif (dialogue non préparé, relance \zh{你呢？} / \zh{你觉得怎么样？}).
    \item \textbf{Ancrage culturel :} Exemples personnels (maladie, soins, grand frère, grand-mère) ancrés dans la réalité camerounaise (famille élargie, respect aînés, solidarité concrète). Comparaison Chine : \zh{孝顺} (xiàoshùn - piété filiale), \zh{春节回家} (Chūnjié huíjiā).
    \item \textbf{Outils linguistiques fixés (Schéma reproductible HSK 3/4) :}
    \begin{center}
    \zh{Salutation $\to$ 题目 $\to$ 我觉得...因为... $\to$ 例如... $\to$ 所以... 谢谢大家!}
    \end{center}
    \item \textbf{Point de vigilance phonétique :} Chaînes de 3e tons (\zh{互相爱护} húxiāng àihù, \zh{演讲题目} yánjiǎng tímù). \textbf{Exercice ciblé « 3e+3e » en échauffement prochaine séance.}
    \item \textbf{Trace écrite obligatoire :} Fiche de préparation (exposé + 2 réponses types) collée dans cahier « Production orale » + Grille d'auto-évaluation complétée.
\end{itemize}
\begin{center}
\textbf{\zh{家人互相帮助，团结就是力量！}} \\
Jiārén hùxiāng bāngzhù, tuánjié jiù shì lìliang ! \\
(Les membres de la famille s'entraident, l'union fait la force !)
\end{center}
\end{aretenir}

\begin{savaistu}[Regard croisé Cameroun — Chine : La famille, socle commun]
\begin{itemize}
    \item \textbf{Cameroun :} Famille élargie (\zh{大家庭} dàjiātíng) centrale. Respect des aînés (\zh{尊老} zūnlǎo). Solidarité financière (\zh{cotisation}, \zh{tontine}) et présence physique (deuils, mariages, fêtes). « \emph{La famille, c'est la banque, l'assurance, l'hôpital} ».
    \item \textbf{Chine :} Famille nucléaire + devoir filial (\zh{孝 xiào} / \zh{孝顺 xiàoshùn}). Fête du Printemps (\zh{春节 Chūnjié}) = plus grande migration humaine annuelle (\zh{春运 chūnyùn}) pour retour au village. Soins aux parents âgés = obligation légale et morale (\zh{养老 yǎnglǎo}).
    \item \textbf{Point commun :} \zh{家和万事兴} (Jiā hé wàn shì xīng) — « Quand la famille est en harmonie, tout prospère ». Devise partagée : \zh{团结就是力量} (Tuánjié jiù shì lìliang).
\end{itemize}
\end{savaistu}

\newpage

\coursec{Méthodes et exercices résolus}

\coursec{Leçon 3 — Expression orale : amour de la famille et la solidarité familiale}

\courssub{Mise en route : parler de sa famille}

\begin{experience}[Échauffement : « Ma famille en 3 mots »]
\textbf{Consigne :} En binôme, chaque élève pioche 3 cartes « adjectifs » (ex: \zh{勤劳}, \zh{温柔}, \zh{严厉}, \zh{乐观}, \zh{体贴}, \zh{健康}) et doit présenter un membre de sa famille en justifiant son choix en une phrase.
\textbf{Exemple :} \zh{我选「体贴」，因为妈妈每天给我做早餐。} (\emph{Wǒ xuǎn « tǐtiē », yīnwèi māma měitiān gěi wǒ zuò zǎocān.})
\textbf{Objectif :} Réactiver le lexique HSK 1-2 (membres, adjectifs, \cle{因为}) et travailler la prononciation des tons (notamment neutres sur la reduplication).
\end{experience}

\begin{textebox}[Dialogue modèle : Présentation rapide]
\textbf{A :} 你家有几口人？\\
\textbf{B :} 我家有五口人：爸爸、妈妈、哥哥、妹妹和我。\\
\textbf{A :} 你们关系怎么样？\\
\textbf{B :} 很好。大家都很\cle{孝顺}，\cle{团结}。哥哥常\cle{帮助}妹妹\cle{复习}功课。\\
\textbf{A :} 真\cle{幸福}！\\
\textbf{Pinyin :} Nǐ jiā yǒu jǐ kǒu rén? / Wǒ jiā yǒu wǔ kǒu rén: bàba, māma, gēge, mèimei hé wǒ. / Nǐmen guānxi zěnmeyàng? / Hěn hǎo. Dàjiā dōu hěn \textbf{xiàoshùn}, \textbf{tuánjié}. Gēge cháng \textbf{bāngzhù} mèimei \textbf{fùxí} gōngkè. / Zhēn \textbf{xìngfú}!
\textbf{Traduction :} Combien de personnes chez toi ? / Nous sommes cinq : papa, maman, grand frère, petite sœur et moi. / Comment s'entend votre famille ? / Très bien. Tous sont très respectueux/filiaux, unis. Grand frère aide souvent petite sœur à réviser. / Vraiment heureux !
\end{textebox}

\courssub{Dialogues modèles et structures clés}

\begin{propriete}[Structure \cle{因为...所以...} (Cause $\rightarrow$ Conséquence)]
\textbf{Schéma :} \zh{因为} + Cause \textbf{，} \zh{所以} + Conséquence.
\textbf{Emploi :} Lien logique fort, réponse à « pourquoi ». \cle{所以} souvent omis à l'oral, mais requis à l'écrit/Bac.
\begin{center}
\begin{tabularx}{\linewidth}{|X|X|X|}
\hline \rowcolor{popblueL} \textbf{Structure} & \textbf{Exemple (Caractères + Pinyin)} & \textbf{Traduction} \\ \hline
\zh{因为...所以...} & \zh{因为父母养育我们不容易，所以我们要孝顺。} (\emph{Yīnwèi fùmǔ yǎngyù wǒmen bù róngyì, suǒyǐ wǒmen yào xiàoshùn.}) & Comme les parents nous élèvent difficilement, nous devons être filiaux. \\ \hline
\zh{因为...所以...} & \zh{因为我生病了，所以全家人来照顾我。} (\emph{Yīnwèi wǒ shēngbìng le, suǒyǐ quánjiā rén lái zhàogù wǒ.}) & Parce que j'étais malade, toute la famille est venue s'occuper de moi. \\ \hline
\end{tabularx}
\end{center}
\textbf{Attention :} Ne pas confondre avec \cle{所以} (donc, conséquence) et \cle{但是} (mais, concession).
\end{propriete}

\begin{propriete}[Structure \cle{虽然...但是/可是...} (Concession)]
\textbf{Schéma :} \zh{虽然} + Fait réel \textbf{，} \zh{但是/可是} + Fait contraire à l'attente.
\textbf{Emploi :} Exprimer une nuance, un contraste. \cle{虽然} ne s'utilise \textbf{jamais} avec \cle{所以}.
\begin{center}
\begin{tabularx}{\linewidth}{|X|X|X|}
\hline \rowcolor{popblueL} \textbf{Structure} & \textbf{Exemple} & \textbf{Traduction} \\ \hline
\zh{虽然...但是...} & \zh{虽然爷爷奶奶年纪大了，但是身体很健康。} (\emph{Suīrán yéye nǎinai niánjì dà le, dànshì shēntǐ hěn jiànkāng.}) & Bien que papy/mamy soient âgés, ils sont en bonne santé. \\ \hline
\zh{虽然...可是...} & \zh{虽然我很忙，可是我每周给家里打电话。} (\emph{Suīrán wǒ hěn máng, kěshì wǒ měi zhōu gěi jiālǐ dǎ diànhuà.}) & Bien que occupé, j'appelle la famille chaque semaine. \\ \hline
\end{tabularx}
\end{center}
\end{propriete}

\begin{propriete}[Structure \cle{一方面...另一方面...} (Énumération d'arguments contraires)]
\textbf{Schéma :} \zh{一方面} + Argument A \textbf{，} \zh{另一方面} + Argument B.
\textbf{Niveau :} HSK 4 / Terminale. Indispensable pour le débat nuancé.
\begin{center}
\begin{tabularx}{\linewidth}{|X|X|X|}
\hline \rowcolor{popblueL} \textbf{Structure} & \textbf{Exemple} & \textbf{Traduction} \\ \hline
\zh{一方面...另一方面...} & \zh{一方面诚实是美德，另一方面要保护父母感受。} (\emph{Yīfāngmiàn chéngshí shì měidé, lìngyīfāngmiàn yào bǎohù fùmǔ gǎnshòu.}) & D'un côté l'honnêteté est une vertu, de l'autre il faut protéger les sentiments des parents. \\ \hline
\end{tabularx}
\end{center}
\end{propriete}

\begin{propriete}[Classificateur \cle{口} (kǒu) : nombre de personnes du foyer)]
\textbf{Règle :} Pour exprimer la taille du ménage (\zh{我家有...口人}), on utilise \textbf{uniquement} \cle{口}.
\textbf{Erreur fréquente :} \zh{*我家有四个 (gè) 人} (trop générique), \zh{*三位 (wèi) 人} (honorifique, écrit).
\textbf{Exemples :} \zh{三口人 sān kǒu rén} (3 personnes), \zh{六口人 liù kǒu rén} (6 personnes).
\end{propriete}

\begin{propriete}[Particules aspectuelles : \cle{了} (le) vs \cle{过} (guò) vs \cle{着} (zhe)]
\begin{center}
\begin{tabularx}{\linewidth}{|X|X|X|X|}
\hline \rowcolor{popblueL} \textbf{Particule} & \textbf{Fonction} & \textbf{Exemple} & \textbf{Traduction} \\ \hline
\cle{了} (le) & Action terminée / Changement d'état \textbf{maintenant} (souvent + temps précis) & \zh{我昨天回家了。} (\emph{Wǒ zuótiān huí jiā le.}) & Je suis rentré hier (action close). \\ \hline
\cle{过} (guò) & Expérience vécue \textbf{au moins une fois} dans le passé (jamais de temps précis) & \zh{我去过北京。} (\emph{Wǒ qù guò Běijīng.}) & Je suis déjà allé à Pékin (expérience). \\ \hline
\cle{着} (zhe) & État durable / Action en cours (souvent verbe 1 + \cle{着} + verbe 2) & \zh{他手里拿着书。} (\emph{Tā shǒuli ná zhe shū.}) & Il tient un livre en main (état). \\ \hline
\end{tabularx}
\end{center}
\textbf{Interdiction :} \zh{*我去过北京了} (double marqueur). \zh{我结婚了} (je me suis marié \textbf{et je le suis encore} / changement d'état) vs \zh{我结过婚} (j'ai été marié \textbf{une fois dans ma vie}, peut-être divorcé).
\end{propriete}

\courssub{Formules fonctionnelles pour exposer et débattre}

\begin{methode}[Présenter un exposé oral simple sur sa famille (1-2 min)]
\textbf{Objectif :} Produire un monologue structuré : Description statique $\rightarrow$ Narration passée $\rightarrow$ Expression sentimentale.
\textbf{Étapes :}
\begin{enumerate}
    \item \textbf{Accroche (开场白) :} \zh{各位同学，大家好。今天我想谈谈我的家庭。} (\emph{Gèwèi tóngxué, dàjiā hǎo. Jīntiān wǒ xiǎng tántan wǒ de jiātíng.})
    \item \textbf{Composition (家庭成员) :} \zh{我家有六口人：爷爷、奶奶、爸爸、妈妈、姐姐和我。} (\emph{Wǒ jiā yǒu liù kǒu rén: yéye, nǎinai, bàba, māma, jiějie hé wǒ.}) \textit{Rappel : classificateur \cle{口}.}
    \item \textbf{Portrait \& Qualité (描述与品质) :} Structure \zh{Sujet + 很/真/特别 + Adjectif}. \zh{爸爸很勤劳，妈妈很体贴，姐姐很聪明。} (\emph{Bàba hěn qínláo, māma hěn tǐtiē, jiějie hěn cōngming.})
    \item \textbf{Exemple de solidarité (团结事例) :} Narration passée (\cle{了}, \cle{过}). Mots-clés : \cle{一起} (ensemble), \cle{帮助} (aider), \cle{鼓励} (encourager), \cle{照顾} (soigner). \zh{去年高考前，我压力很大。全家人都来鼓励我，姐姐帮我复习，妈妈做好吃的，爸爸陪我散步。} (\emph{Qùnián gāokǎo qián, wǒ yālì hěn dà. Quánjiā rén dōu lái gǔlì wǒ, jiějie bāng wǒ fùxí, māma zuò hǎochī de, bàba péi wǒ sànbù.})
    \item \textbf{Conclusion sentimentale (结语) :} \zh{我觉得家是最温暖的港湾，家人永远是我最坚强的后盾。我爱我的家，谢谢大家！} (\emph{Wǒ juéde jiā shì zuì wēnnuǎn de gǎngwān, jiārén yǒngyuǎn shì wǒ zuì jiānqiáng de hòudùn. Wǒ ài wǒ de jiā, xièxie dàjiā!})
    \item \textbf{Entraînement tonal ciblé :} Repérer sandhi \cle{不} / \cle{一} et tons neutres (voir section Tons).
\end{enumerate}
\textbf{Structures à retenir :} \zh{我家有...} | \zh{...很/真/特别+形容词} | \zh{大家一起+动词} | \zh{我觉得...因为...} | \zh{虽然...但是...}
\end{methode}

\begin{exoresolu}[Exposé type : « Ma famille unie face aux épreuves »]
\textbf{Énoncé :} Préparez un exposé oral d'une minute sur « La solidarité dans ma famille ». Contraintes : $\ge$ 5 membres, $\ge$ 3 adjectifs de qualité, 1 exemple concret d'entraide, conclusion type. Structure logique/chronologique.
\tcblower
\textbf{Solution détaillée :}
\textbf{1. Analyse :} Plan narratif : Présentation $\rightarrow$ Description $\rightarrow$ Narration (passé) $\rightarrow$ Sentiment (présent).
\textbf{2. Vocabulaire clé (HSK 2-3) :} \zh{爷爷 yéye}, \zh{奶奶 nǎinai}, \zh{姐姐 jiějie}, \zh{弟弟 dìdi} | \zh{健康 jiànkāng}, \zh{严厉 yánlì}, \zh{体贴 tǐtiē}, \zh{乐观 lèguān} | \zh{照顾 zhàogù}, \zh{鼓励 gǔlì}, \zh{做饭 zuòfàn}, \zh{复习 fùxí}.
\textbf{3. Texte modèle (Caractères + Pinyin + Traduction) :}
\begin{textebox}[Exposé modèle]
各位同学，大家好。\\
今天我想谈谈我的家庭。\\
我家有六口人：爷爷、奶奶、爸爸、妈妈、姐姐和我。\\
爷爷奶奶虽然年纪大了，但身体很\cle{健康}，精神很好。爸爸很\cle{勤劳}，妈妈很\cle{体贴}，姐姐很\cle{聪明}。\\
去年高考前，我压力很大，心情不好。全家人都来\cle{鼓励}我。姐姐帮我\cle{复习}功课，妈妈每天给我\cle{做}好吃的饭，爸爸陪我\cle{散步}放松。\\
最后我考上了大学。我觉得\cle{家}是最\cle{温暖}的港湾，\cle{家人}永远是我最坚强的后盾。\\
我爱我的家，谢谢大家！
\end{textebox}
\textbf{Pinyin :} Gèwèi tóngxué, dàjiā hǎo. / Jīntiān wǒ xiǎng tántan wǒ de jiātíng. / Wǒ jiā yǒu liù kǒu rén: yéye, nǎinai, bàba, māma, jiějie hé wǒ. / Yéye nǎinai suīrán niánjì dà le, dàn shēntǐ hěn \textbf{jiànkāng}, jīngshén hěn hǎo. Bàba hěn \textbf{qínláo}, māma hěn \textbf{tǐtiē}, jiějie hěn \textbf{cōngming}. / Qùnián gāokǎo qián, wǒ yālì hěn dà, xīnqíng bù hǎo. Quánjiā rén dōu lái \textbf{gǔlì} wǒ. Jiějie bāng wǒ \textbf{fùxí} gōngkè, māma měitiān gěi wǒ \textbf{zuò} hǎochī de fàn, bàba péi wǒ \textbf{sànbù} fàngsōng. / Zuìhòu wǒ kǎo shàng le dàxué. Wǒ juéde \textbf{jiā} shì zuì \textbf{wēnnuǎn} de gǎngwān, \textbf{jiārén} yǒngyuǎn shì wǒ zuì jiānqiáng de hòudùn. / Wǒ ài wǒ de jiā, xièxie dàjiā!
\textbf{Traduction :} Chers camarades, bonjour. Aujourd'hui je parle de ma famille. Nous sommes six : papy, mamy, papa, maman, grande sœur et moi. Papy et mamy, bien que âgés, sont en bonne santé, bon moral. Papa travailleur, maman attentionnée, sœur intelligente. L'an dernier avant le Bac, grosse pression, mauvais moral. Toute la famille m'a encouragé. Sœur a aidé réviser, maman cuisinait bon chaque jour, papa m'accompagnait marcher détendre. Finalement j'ai réussi entrée université. Je trouve que famille est le port le plus chaleureux, les proches sont à jamais mon plus solide soutien. J'aime ma famille, merci !
\textbf{4. Points de vigilance oral :}
\begin{itemize}
    \item \zh{六口人 liù kǒu rén} (4-3-2) -- \cle{口} obligatoire, pas \cle{个}.
    \item \zh{虽然...但是...} : structure concession attendue au Bac.
    \item \zh{港湾 gǎngwān} (port/abri) : image métaphorique valorisée.
    \item \zh{考上了 kǎo shàng le} (V + Complément de résultat + \cle{led}).
\end{itemize}
\end{exoresolu}

\begin{methode}[Participer à un débat guidé : « La famille d'abord ? »]
\textbf{Objectif :} Interagir pour exprimer accord/désaccord nuancé, relancer, argumenter avec \cle{因为...所以...}, \cle{虽然...但是...}, \cle{一方面...另一方面...}.
\textbf{Étapes :}
\begin{enumerate}
    \item \textbf{Prendre la parole :} \zh{我觉得...}, \zh{在我看来...}, \zh{我想补充一点...}
    \item \textbf{Accord :} \zh{我同意你的看法。}, \zh{你说得很对。}, \zh{完全正确。}
    \item \textbf{Désaccord poli :} \zh{我有不同的看法。}, \zh{这倒不一定。}, \zh{一方面...，另一方面...}
    \item \textbf{Argumenter :} Structure rigide \zh{因为...，所以...}. Ex: \zh{因为父母养育我们不容易，所以我们要孝顺。}
    \item \textbf{Relancer :} \zh{你怎么看？}, \zh{如果...会怎么样？}, \zh{举个例子吧。}
    \item \textbf{Conclure :} \zh{总之，家庭很重要。}
\end{enumerate}
\textbf{Formules clés (à mémoriser) :}
\begin{center}
\begin{tabularx}{\linewidth}{|X|X|X|}
\hline \rowcolor{popblueL} \textbf{Fonction} & \textbf{Chinois (Pinyin)} & \textbf{Français} \\ \hline
Accord fort & \zh{我完全赞成 (Wǒ wánquán zànchéng)} & Je suis tout à fait d'accord \\ \hline
Accord partiel & \zh{有道理，但是... (Yǒu dàoli, dànshì...)} & C'est vrai, mais... \\ \hline
Désaccord poli & \zh{恐怕不太对吧 (Kǒngpà bù tài duì ba)} & Je crains que ce ne soit pas tout à fait exact \\ \hline
Demander avis & \zh{你怎么看这个问题？ (Nǐ zěnme kàn zhège wèntí?)} & Comment vois-tu ce problème ? \\ \hline
Citer exemple & \zh{比如说... (Bǐrú shuō...)} & Par exemple... \\ \hline
Nuancer & \zh{视情况而定 (Shì qíngkuàng ér dìng)} & Selon les circonstances / Ça dépend \\ \hline
\end{tabularx}
\end{center}
\end{methode}

\begin{exoresolu}[Débat simulé : « Faut-il tout dire à ses parents ? »]
\textbf{Énoncé :} Complétez le dialogue entre \zh{小明 (Xiǎomíng)} (transparence) et \zh{小红 (Xiǎohóng)} (protection). Utilisez \cle{因为...所以...}, \cle{虽然...但是...}, \cle{一方面...另一方面...}. Registre oral (particules \cle{吧}, \cle{呢}, \cle{啊}).
\tcblower
\textbf{Dialogue complété :}
\begin{textebox}[Débat : 家庭透明度]
\textbf{小明}：小红，你觉得我们应该对父母\cle{隐瞒} (yǐnmán) certaines choses吗？
\textbf{小红}：\cle{这倒不一定} (Zhè dào bù yídìng)。\cle{一方面} (Yīfāngmiàn)，诚实是美德；\cle{另一方面} (lìngyīfāngmiàn)，父母年纪大了，\cle{如果} (rúguǒ) 全告诉他们，他们会\cle{担心} (dānxīn)。\cle{所以} (Suǒyǐ)，有时候\cle{善意的谎言} (shànyì de huǎngyán) 是必要的。
\textbf{小明}：\cle{我有不同看法} (Wǒ yǒu bùtóng kànfǎ)。\cle{因为} (Yīnwèi) 信任是家庭的基础，\cle{所以} (suǒyǐ) 一旦撒谎，\cle{就} (jiù) 破坏了信任。\cle{虽然} (Suīrán) 说真话可能让他们\cle{生气} (shēngqì) 或担心，\cle{但是} (dànshì) 从长远看，坦诚更好。
\textbf{小红}：\cle{有道理} (Yǒu dàoli)。\cle{不过} (Bùguò)，\cle{视情况而定} (shì qíngkuàng ér dìng) 吧。比如考砸了，先不说，等心情好再谈。
\textbf{小明}：\cle{也行} (Yě xíng)。\cle{总之} (Zǒngzhī)，目的是让家庭\cle{和睦} (hémù)。
\end{textebox}
\textbf{Analyse grammaticale :}
\begin{itemize}
    \item \zh{隐瞒 yǐnmán} (v. cacher) -- HSK 4.
    \item \zh{这倒不一定} : \cle{倒} renforce la nuance « contrairement à l'attente ».
    \item \zh{一方面...另一方面...} : pondération.
    \item \zh{一旦...就...} : conditionnelle hypothétique (niveau Terminale/HSK 4).
    \item \cle{视情况而定} : locution figée.
    \item Particules finales : \cle{吧} (suggestion), \cle{呢} (relance), \cle{啊} (constat).
\end{itemize}
\textbf{Entraînement tonal :} \zh{隐瞒 yǐn mán} (3-2), \zh{担心 dān xīn} (1-1), \zh{善意的谎言 shàn yì de huǎng yán} (4-4 - 3-2), \zh{和睦 hé mù} (2-4).
\end{exoresolu}

\courssub{Les tons à travailler : pièges et sandhi}

\begin{methode}[Maîtriser les tons critiques du vocabulaire familial (HSK 1-3)]
\textbf{Objectif :} Éliminer les erreurs fossilisées. La mauvaise tonalité change le sens ou rend incompréhensible.
\textbf{Démarche :} 1. Identifier paires minimales/pièges. 2. S'entraîner syllabes $\rightarrow$ mots $\rightarrow$ phrases (exagérer courbe mélodique). 3. Appliquer en contexte (dialogues, sandhi \cle{不}, \cle{一}).
\textbf{Tableau des 10 « Tons Pièges » :}
\begin{center}
\begin{tabularx}{\linewidth}{|X|X|X|X|}
\hline \rowcolor{popblueL} \textbf{Mot (Caractères)} & \textbf{Pinyin Correct} & \textbf{Erreur fréquente (Francophone)} & \textbf{Sens / Contexte} \\ \hline
\zh{父母} & \textbf{fù mǔ} (4-3) & fù mù (4-4) ou fǔ mǔ (3-3) & Parents. \cle{母} mǔ (3) = mère. \\ \hline
\zh{爷爷 / 奶奶} & \textbf{yé ye / nǎi nai} (2-neutre / 3-neutre) & yé yé (2-2) / nǎi nǎi (3-3) & 2e syllabe \textbf{ton neutre} (轻声) obligatoire ! \\ \hline
\zh{哥哥 / 姐姐 / 弟弟 / 妹妹} & \textbf{gē ge / jiě jie / dì di / mèi mei} (1-neutre / 3-neutre / 4-neutre / 4-neutre) & Tons pleins sur 2e syllabe & Frères/Sœurs. \textbf{Règle :} Reduplication $\rightarrow$ 2e = ton neutre. \\ \hline
\zh{勤劳} & \textbf{qín láo} (2-2) & qín lǎo (2-3) & Travaillador. \cle{劳} 2e ton ! (pas \cle{老} lǎo = vieux). \\ \hline
\zh{体贴} & \textbf{tǐ tiē} (3-1) & tì tiē (4-1) & Attentionné. \cle{体} 3e ton (bas). \\ \hline
\zh{孝顺} & \textbf{xiào shùn} (4-4) & xiǎo shùn (3-4) & Filial. \cle{孝} 4e ton (chute). Pas \cle{小} xiǎo (petit). \\ \hline
\zh{团圆} & \textbf{tuán yuán} (2-2) & tuán yuǎn (2-3) & Réunion (fête). \cle{圆} 2e ton. \\ \hline
\zh{和睦} & \textbf{hé mù} (2-4) & hè mù (4-4) / hú mù (2-3) & Harmonieux. \cle{和} = hé (2) ici (pas hè/hú/huo). \\ \hline
\zh{照顾} & \zh{zhào gù} (4-4) & zhǎo gù (3-4) & S'occuper. \cle{照} 4e ton. \\ \hline
\zh{鼓励} & \zh{gǔ lì} (3-4) & gǔ lǐ (3-3) & Encourager. \cle{励} 4e ton. \\ \hline
\end{tabularx}
\end{center}

\end{methode}

\begin{propriete}[Règles d'or du Sandhi (变调) : \cle{不} et \cle{一}]
\textbf{1. \cle{不} (bù, ton 4 d'origine) :}
\begin{itemize}
    \item Devant ton \textbf{4} $\rightarrow$ devient \textbf{ton 2 (bú)}. Ex: \zh{不客气 bú kèqi} (2-4-4), \zh{不是 bú shì} (2-4).
    \item Devant tons \textbf{1, 2, 3} $\rightarrow$ reste \textbf{ton 4 (bù)}. Ex: \zh{不忙 bù máng} (4-2), \zh{不来 bù lái} (4-2), \zh{不想 bù xiǎng} (4-3).
\end{itemize}
\textbf{2. \cle{一} (yī, ton 1 d'origine) :}
\begin{itemize}
    \item Seul, comptage, ordinal, fin de phrase $\rightarrow$ \textbf{ton 1 (yī)}. Ex: \zh{一 yī}, \zh{第一 dì yī}, \zh{看一看 kàn yi kàn} (neutre).
    \item Devant ton \textbf{4} $\rightarrow$ \textbf{ton 2 (yí)}. Ex: \zh{一定 yí dìng}, \zh{一样 yí yàng}.
    \item Devant tons \textbf{1, 2, 3} $\rightarrow$ \textbf{ton 4 (yì)}. Ex: \zh{一起 yì qǐ} (4-3), \zh{一家人 yì jiā rén} (4-1-2), \zh{一起 yì qǐ}.
    \item Au milieu de reduplication (V-\cle{一}-V) $\rightarrow$ \textbf{ton neutre (yi)}. Ex: \zh{看一看 kàn yi kàn}, \zh{想一想 xiǎng yi xiǎng}.
\end{itemize}
\end{propriete}

\begin{exoresolu}[Dictée tonale \& Correction de phrases types]
\textbf{Énoncé :} Écrivez le pinyin avec tons corrects. Corrigez les erreurs soulignées.
\tcblower
\textbf{1. Dictée tonale (Pinyin cible) :}
\begin{enumerate}
    \item \zh{我的爷爷奶奶很健康。} $\rightarrow$ \textbf{Wǒ de yéye nǎinai hěn jiànkāng.} (yéye, nǎinai : ton neutre).
    \item \zh{爸爸很勤劳，妈妈很体贴。} $\rightarrow$ \textbf{Bàba hěn qínláo, māma hěn tǐtiē.}
    \item \zh{我们一起吃饭，很团圆。} $\rightarrow$ \textbf{Wǒmen yìqǐ chīfàn, hěn tuányuán.} (Sandhi \cle{一} $\rightarrow$ yì devant ton 3 \cle{起}).
    \item \zh{因为孝顺，所以家庭和睦。} $\rightarrow$ \textbf{Yīnwèi xiàoshùn, suǒyǐ jiātíng hémù.}
    \item \zh{我不累，我想照顾弟弟。} $\rightarrow$ \textbf{Wǒ bú lèi, wǒ xiǎng zhàogù dìdi.} (Sandhi \cle{不} $\rightarrow$ bú devant ton 4 \cle{累}).
\end{enumerate}
\textbf{2. Correction d'erreurs (Élève $\rightarrow$ Correct) :}
\begin{center}
\begin{tabularx}{\linewidth}{|X|X|X|}
\hline \rowcolor{popblueL} \textbf{Phrase élève (Erreurs)} & \textbf{Correction} & \textbf{Explication} \\ \hline
Wǒ de yéye nǎinai hěn jiànkang. & \textbf{Wǒ de yéye nǎinai hěn jiànkāng.} & \cle{康} kāng (1er ton haut), pas ton neutre. \\ \hline
Bàba hěn qínlǎo. & \textbf{Bàba hěn qínláo.} & \cle{劳} láo (2e ton montant), pas lǎo (3e ton = vieux). \\ \hline
Māma hěn tìtiē. & \textbf{Māma hěn tǐtiē.} & \cle{体} tǐ (3e ton bas), pas tì (4e ton = substitut). \\ \hline
Wǒmen yīqǐ chīfàn. & \textbf{Wǒmen yìqǐ chīfàn.} & Sandhi \cle{一} : devant ton 3 (\cle{起 qǐ}) $\rightarrow$ \cle{yì} (4e ton). \\ \hline
Wǒ bù xiǎng qù. & \textbf{Wǒ bù xiǎng qù.} & \cle{想} xiǎng (3e ton). Devant ton 3, \cle{不} reste \cle{bù} (4). \textbf{Pas de changement !} \\ \hline
Wǒ bù kèqi. & \textbf{Wǒ bú kèqi.} & Devant ton 4 (\cle{客 kè}) $\rightarrow$ \cle{bú} (2). \\ \hline
Jiātíng hémù hěn hǎo. & \textbf{Jiātíng hémù hěn hǎo.} & \cle{和} hé (2), \cle{睦} mù (4). Correct. \\ \hline
Tā xiǎoshun. & \textbf{Tā xiàoshùn.} & \cle{孝} xiào (4), \cle{顺} shùn (4). Pas xiǎo (3 = petit). \\ \hline
\end{tabularx}
\end{center}
\textbf{Bilan :} Erreurs majeures = 1. Tons neutres (famille), 2. Confusion \cle{ao/ou} (\cle{láo} vs \cle{lǎo}), 3. Sandhi \cle{不/一} (conditionnel au ton \textbf{suivant}), 4. Homophones \cle{xiào/xiǎo}.
\end{exoresolu}

\courssub{Jeux de rôle guidés}

\begin{experience}[Jeu de rôle : « Le dîner du dimanche »]
\textbf{Scénario :} 3-4 élèves (Père, Mère, Grand-frère/Grande-sœur, Cadet). Contexte : Dîner familial, on parle de la semaine, des notes, de la fatigue, on se passe les plats.
\textbf{Rôles \& Contraintes linguistiques :}
\begin{itemize}
    \item \textbf{Père/Mère :} Utiliser \cle{叮嘱} (rappeler/conseiller), \cle{关心} (se soucier), structure \cle{多+V} / \cle{少+V} (ex: \zh{多吃菜，少喝可乐}).
    \item \textbf{Aîné(e) :} Utiliser \cle{鼓励} (encourager), \cle{帮助} (aider), \cle{分享} (partager).
    \item \textbf{Cadet(te) :} Utiliser \cle{撒娇} (faire le capricieux/mignon), \cle{答应} (promettre), \cle{谢谢} / \cle{我知道了}.
    \item \textbf{Tous :} Politesse \cle{您} (nín) pour aînés, classificateur \cle{口}, tons neutres kinship.
\end{itemize}
\textbf{Exemple de réplique :} \zh{妈妈：儿子，多吃点排骨，长身体。别光吃饭啊！} (\emph{Māma: Érzi, duō chī diǎn páigǔ, zhǎng shēntǐ. Bié guāng chīfàn a!})
\textbf{Débrief :} Enregistrement $\rightarrow$ auto-correction tons (surtout neutres) + ordre mots (Temps/Lieu avant V).
\end{experience}

\courssub{Débat guidé : la famille d'abord ?}
(Voir \textbf{Méthode « Participer à un débat guidé »} et \textbf{Exoresolu « Débat simulé »} ci-dessus pour les structures et le modèle complet.)
\begin{exercice}[Sujets de débat pour l'entraînement (niveau Première)]
\begin{enumerate}
    \item \zh{「成绩不好，要不要告诉父母？」} (Mauvaises notes : le dire ou pas ?)
    \item \zh{「大学要不要离家乡远一点？」} (Étudier loin ou près de la famille ?)
    \item \zh{「做家务是父母的事，还是孩子的事？」} (Tâches ménagères : parents ou enfants ?)
    \item \zh{「网络时代，家人聚在一起却都在看手机，怎么办？」} (Écrans lors des réunions familiales).
\end{enumerate}
\textbf{Consigne :} Par groupes de 4. 5 min préparation (fiche mots-clés), 5 min débat. Un observateur note : structures utilisées, tons, politesse. Rotation des rôles.
\end{exercice}

\courssub{Expression écrite : messages et vœux familiaux}

\begin{methode}[Rédiger un message court / carte de vœux (SMS, WeChat, Carte)]
\textbf{Objectif :} Produit écrit court pour fête (Nouvel An, Fête Mères/Pères, Anniversaire) selon codes socioculturels chinois.
\textbf{Étapes :}
\begin{enumerate}
    \item \textbf{Identification :} Destinataire (\zh{给妈妈}, \zh{给爷爷奶奶}) + Occasion (\zh{春节}, \zh{母亲节}, \zh{生日}).
    \item \textbf{Structure type :}
    \begin{itemize}
        \item \textbf{Appel (称呼) :} \zh{亲爱的妈妈：} (proches) | \zh{尊敬的爷爷奶奶：} (respect aînés).
        \item \textbf{Corps (正文) :} \zh{祝您...} + \textbf{Chengyu (成语) 4 caractères}.
        \item \textbf{Signature (落款) :} \zh{您的女儿/儿子：xxx} | \zh{孙子/孙女：xxx} + \zh{日期}.
    \end{itemize}
    \item \textbf{Chengyu indispensables (吉祥话) :}
    \begin{center}
    \begin{tabularx}{\linewidth}{|X|X|X|}
    \hline \rowcolor{popblueL} \textbf{Chinois} & \textbf{Pinyin} & \textbf{Sens / Usage} \\ \hline
    \zh{身体健康} & shēntǐ jiànkāng & Santé (Priorité \#1 aînés) \\ \hline
    \zh{万事如意} & wànshì rúyì & Tout selon vos vœux \\ \hline
    \zh{阖家幸福} & hèjiā xìngfú & Bonheur famille entière \\ \hline
    \zh{福寿安康} & fú shòu ānkāng & Bonheur, longévité, santé (Grands-parents) \\ \hline
    \zh{笑口常开} & xiàokǒu chángkāi & Sourire permanent (Joie) \\ \hline
    \end{tabularx}
    \end{center}
    \item \textbf{Écriture caractères (ordre traits) :} \zh{福 fú} : 示 (esprit) + 畐 (phonétique) -- haut $\rightarrow$ bas. \zh{寿 shòu} : compact, haut $\rightarrow$ bas. \zh{家 jiā} : 宀 (toit) + 豕 (cochon) -- toit d'abord.
    \item \textbf{Traduction Thème (Fr $\rightarrow$ Ch) :} Pas de calque. « Joyeuse fête Maman » $\neq$ \zh{快乐母亲节}. Correct : \zh{妈妈，母亲节快乐！祝您身体健康，万事如意！}
\end{enumerate}
\end{methode}

\begin{exoresolu}[Rédaction \& Traduction : Carte WeChat pour la Fête du Printemps]
\textbf{Énoncé :} Vous êtes \zh{李明 (Lǐ Míng)}, étudiant à Yaoundé. Message WeChat aux grands-parents pour \zh{春节 Chūnjié}. $\ge$ 3 chengyu 4 caractères. Traduction FR.
\tcblower
\textbf{1. Rédaction (Caractères + Pinyin) :}
\begin{textebox}[Message WeChat - 微信消息]
尊敬的爷爷、奶奶：\\
\indent 春节快乐！\\
\indent 孙儿李明在喀麦隆雅温得向您二老拜年了！\\
\indent 祝您二老：\cle{福寿安康}，\cle{身体健康}，\cle{万事如意}，\cle{阖家幸福}！\\
\indent 孙儿虽然不在身边，但心一直在家。等放假一定回家陪您吃\cle{饺子}，包\cle{春卷}。\\
\indent 请保重身体！\\
您的孙儿：李明\\
二〇二五年二月一日
\end{textebox}
\textbf{Pinyin :} Zūnjìng de yéye, nǎinai: / Chūnjié kuàilè! / Sūnér Lǐ Míng zài Kāmàilóng Yǎwēndé xiàng nín èr lǎo bàinián le! / Zhù nín èr lǎo: \textbf{fú shòu ānkāng}, \textbf{shēntǐ jiànkāng}, \textbf{wànshì rúyì}, \textbf{hèjiā xìngfú}! / Sūnér suīrán bù zài shēnbiān, dàn xīn yīzhí zài jiā. Děng fàngjià yídìng huíjiā péi nín chī \textbf{jiǎozi}, bāo \textbf{chūnjuǎn}. / Qǐng bǎozhòng shēntǐ! / Nín de sūnér: Lǐ Míng. / Èrlíng'èrwǔ nián èryuè yīrì.
\textbf{2. Traduction française soignée :}
« Chers grand-père, grand-mère,
Joyeux Nouvel An (Fête du Printemps) !
Votre petit-fils Li Ming, depuis Yaoundé au Cameroun, vous souhaite la bonne année à tous les deux !
Je vous souhaite : bonheur-longévité-santé, santé physique, que tout se passe comme vous le désirez, bonheur pour toute la famille !
Bien que votre petit-fils ne soit pas à vos côtés, mon cœur est toujours à la maison. Quand j'aurai les vacances, je rentrerai sans faute pour vous accompagner manger des raviolis et faire des rouleaux de printemps.
Prenez bien soin de vous !
Votre petit-fils : Li Ming
1er février 2025 »
\textbf{3. Analyse linguistique \& culturelle :}
\begin{itemize}
    \item \zh{尊敬的 zūnjìng de} : Appel formel aînés (vs \zh{亲爱的} parents/proches).
    \item \zh{二老 èr lǎo} : Respectueux « vous deux les anciens ».
    \item \zh{拜年 bài nián} : Verbe spécifique « souhaiter bonne année (en s'inclinant) ».
    \item \zh{饺子 jiǎozi} / \zh{春卷 chūnjuǎn} : Symboles Nouvel An (richesse/printemps).
    \item \zh{保重 bǎozhòng} : « Prendre soin de soi » (clôture aînés/malades).
    \item \textbf{Ordre vœux :} Santé (\cle{身体健康}, \cle{福寿安康}) $\rightarrow$ Réussite (\cle{万事如意}) $\rightarrow$ Famille (\cle{阖家幸福}). Hiérarchie confucéenne : Individu $\rightarrow$ Monde $\rightarrow$ Groupe.
\end{itemize}
\end{exoresolu}

\courssub{Compréhension de l'oral : témoignage}

\begin{methode}[Analyser un texte oral authentique (Méthode « 3 Écoutes »)]
\textbf{Objectif :} Extraire infos clés (Qui ? Quoi ? Quand ? Sentiment ?) et répondre en chinois.
\textbf{Démarche :}
\begin{enumerate}
    \item \textbf{1ère écoute (Global) :} Thème, locuteurs, ton (joyeux/triste/sérieux). Pas d'écriture.
    \item \textbf{2ème écoute (Détail) :} Noter \textbf{mots-clés} (noms, verbes, adjectifs, chiffres, temps) + \textbf{connecteurs} : \zh{因为/所以}, \zh{虽然/但是}, \zh{而且}, \zh{不过}.
    \item \textbf{3ème écoute (Vérif) :} Vérifier réponses précises. Reconstituer structure (Intro $\rightarrow$ Exemple $\rightarrow$ Conclusion).
\end{enumerate}
\textbf{Types questions (Bac) :} QCM/Vrai-Faux (\zh{判断对错}), Ouvertes courtes (\zh{为什么...？}, \zh{谁...？}), Complétion/Résumé (\zh{补全句子}, \zh{用中文概括大意 50 caractères max}).
\textbf{Astuce :} Notes en \textbf{pinyin abrégé} ou \textbf{mots-clés FR} pendant l'écoute (pas caractères complets = trop lent).
\end{methode}

\begin{exoresolu}[Compréhension orale : Témoignage de Wang Wei (Douala)]
\textbf{Énoncé :} Écoutez \zh{王伟 (Wáng Wěi)}, étudiant chinois à Douala. Répondez en chinois (caractères + pinyin).
\textit{(Script audio pour enseignant/auto-correction)}
\tcblower
\textbf{Script audio :}
\zh{王伟：大家好。我叫王伟，今年二十二岁，在杜阿拉大学学习法语。我家在北京，有爸爸、妈妈和妹妹。妹妹叫王芳，今年十八岁，上高中。去年圣诞节，我没回家，很想念家人。妈妈微信视频给我，看见妹妹长高了，也懂事了。妈妈说：「伟伟，在外面要照顾好自己，少吃外卖，多吃蔬菜。」我听了很感动，也很内疚。今年春节，我决定请假两周回北京团圆。给爸爸买双皮鞋，给妈妈买件围巾，给妹妹买个书包。希望全家人健康、快乐、团圆。}
\textbf{Questions \& Réponses modèles :}
\begin{enumerate}
    \item \zh{王伟在哪里学习？} $\rightarrow$ \zh{在杜阿拉大学。} (\emph{Zài Dù'ālā Dàxué.})
    \item \zh{王伟家有几口人？谁？} $\rightarrow$ \zh{四口人：爸爸、妈妈、妹妹和他。} (\emph{Sì kǒu rén: bàba, māma, mèimei hé tā.})
    \item \zh{去年圣诞节王伟怎么了？} $\rightarrow$ \zh{他没回家，很想念家人。} (\emph{Tā méi huí jiā, hěn xiǎngniàn jiārén.})
    \item \zh{妈妈在视频里对他说了什么？} $\rightarrow$ \zh{在外面要照顾好自己，少吃外卖，多吃蔬菜。} (\emph{Zài wàimiàn yào zhàogù hǎo zìjǐ, shǎo chī wàimài, duō chī shūcài.})
    \item \zh{王伟今年春节打算做什么？} $\rightarrow$ \zh{请假两周回北京团圆，给家人买礼物。} (\emph{Qǐngjià liǎng zhōu huí Běijīng tuányuán, gěi jiārén mǎi lǐwù.})
    \item \zh{用一个成语概括王伟的心情。} $\rightarrow$ \zh{归心似箭} (\emph{guīxīn sìjiàn} -- le cœur du retour vole comme une flèche / impatient de rentrer).
\end{enumerate}
\textbf{Grammaire repérée :}
\begin{itemize}
    \item \zh{没回家 méi huí jiā} : Négation passé (\cle{没}+V).
    \item \zh{想念 xiǎngniàn} : Verbe transitif direct.
    \item \zh{长高了 zhǎng gāo le} / \zh{懂事了 dǒngshì le} : Changement d'état (\cle{led} final).
    \item \zh{少吃...多吃...} : Structure conseil/impératif.
    \item \zh{决定请假... juédìng qǐngjià...} : Chaîne verbale (V1+V2).
\end{itemize}
\end{exoresolu}

\courssub{Éléments socioculturels : Cameroun – Chine}

\begin{savaistu}[Famille : Valeurs partagées et différences]
\textbf{Points communs (Valeurs universelles) :}
\begin{itemize}
    \item \textbf{Respect des aînés / Piété filiale :} Cameroun : respect absolu aux parents/grands-parents, on ne contredit pas l'aîné en public. Chine : \cle{孝 xiào} (piété filiale) est la vertu cardinale confucéenne (\zh{百善孝为先} -- « Parmi cent vertus, la piété filiale vient en premier »).
    \item \textbf{Solidarité élargie :} Cameroun : famille élargie (oncles, tantes, cousins = « frères/sœurs »), entraide financière/éducative forte. Chine : historiquement clan/famille élargie ; aujourd'hui famille nucléaire (politique enfant unique 1979-2015, puis 2 puis 3 enfants), mais liens forts maintenus pour fêtes (\zh{春节}, \zh{中秋节}) et soin aînés.
    \item \textbf{Éducation :} Sacrifice parental pour la scolarité des enfants (Cameroun : « l'enfant est la retraite des parents » ; Chine : \cle{望子成龙} wàng zǐ chéng lóng -- « espérer que le fils devienne un dragon » = réussite sociale).
\end{itemize}
\textbf{Différences notables :}
\begin{itemize}
    \item \textbf{Nommage / Appellation :} Cameroun : souvent prénom + titre (Tonton Paul, Maman Marie) ou termes génériques « Papa », « Maman » pour oncles/tantes. Chine : \textbf{système très codifié} selon lignée (paternelle/maternelle), âge relatif, sexe. Ex: \zh{伯父 bófu} (oncle paternel aîné), \zh{叔叔 shūshu} (oncle paternel cadet), \zh{舅舅 jiùjiu} (oncle maternel), \zh{姑姑 gūgu} (tante paternelle), \zh{姨妈 yímā} (tante maternelle). \textit{Erreur fréquente :} Appeler un oncle maternel \cle{叔叔}.
    \item \textbf{Charge du soin aux aînés :} Chine : traditionnellement le fils aîné (\zh{长子 zhǎngzǐ}) et sa femme (\zh{儿媳 érxí}) cohabitent avec parents. Loi 2013 : obligation légale de visite/soin (\zh{养老 yǎnglǎo}). Cameroun : charge partagée entre tous les enfants, souvent rotation ou contribution financière ; cohabitation moins systématique en ville.
    \item \textbf{Fêtes :} Chine : \zh{春节} (réunion impérative, \cle{团圆}), \zh{清明节} (nettoyage tombes ancêtres). Cameroun : Fêtes chrétiennes (Noël, Pâques), fêtes traditionnelles (funérailles = moment majeur de réunion clanique, \cle{deuil} = solidarité active).
\end{itemize}
\textbf{Pour le Bac oral :} Savoir citer \zh{百善孝为先} et expliquer \cle{团圆} vs réunion familiale camerounaise (funérailles, Noël).
\end{savaistu}

\courssub{Bilan et entraînement à l'exposé final}

\begin{attention}[Les 5 erreurs qui coûtent des points au Bac (Leçon 3)]
\begin{enumerate}
    \item \textbf{Tons : Famille = Ton Neutre obligatoire.}
    \textbf{Erreur :} \zh{爸爸 bàba (4-4), 妈妈 māma (1-1), 哥哥 gēge (1-1), 妹妹 mèimei (4-4)}.
    \textbf{Correct :} \zh{爸爸 bàba, 妈妈 māma, 哥哥 gēge, 妹妹 mèimei, 弟弟 dìdi, 姐姐 jiějie, 爷爷 yéye, 奶奶 nǎinai}. \textit{Règle : Reduplication noms de parenté $\rightarrow$ 2e syllabe ton neutre (轻声 qīngshēng).} Perte 1 pt/mot à l'oral.
    \item \textbf{Caractères confondus : \cle{孝} vs \cle{校} vs \cle{效} | \cle{养} vs \cle{样} | \cle{团} vs \cle{抟}.}
    \textbf{Erreur :} \zh{校顺 xiào shùn} (école+obéir) au lieu de \zh{孝顺 xiàoshùn}. \zh{抟圆 tuán yuán} au lieu de \zh{团圆 tuányuán}.
    \textbf{Astuce :} \cle{孝} = \cle{老} (vieux) sur \cle{子} (enfant) = l'enfant porte le vieux. \cle{团} = \cle{囗} (enceinte) + \cle{专} (spécial) = groupe uni.
    \item \textbf{Ordre des mots : Complément Temps/Lieu AVANT le verbe.}
    \textbf{Erreur (Calque FR) :} \zh{*我回家了昨天} / \zh{*我在北京住}.
    \textbf{Correct :} \zh{我昨天回家了} (S + Temps + Lieu + V + \cle{led}). \zh{我住在北京} (S + V + \cle{在} + Lieu). \textit{Règle :} \cle{在} (zài) introduit lieu \textbf{avant} V d'action/état ; \cle{led} suit V principal.
    \item \textbf{Mots-outils : \cle{了} (le) vs \cle{过} (guò) vs \cle{着} (zhe).}
    \textbf{Erreur :} \zh{*我去过北京了} (Double marqueur). \zh{我结婚了} (changement d'état actuel) vs \zh{我结过婚} (expérience vie).
    \textbf{Règle :} \cle{led} = action close/changement \textbf{maintenant} (souvent + temps précis). \cle{guò} = expérience vécue \textbf{au moins une fois} (jamais temps précis). \cle{zhe} = durée/état en cours.
    \item \textbf{Classificateurs : \cle{口} (kǒu) pour taille du foyer.}
    \textbf{Erreur :} \zh{*我家有四个 (gè) 人} / \zh{*三位 (wèi) 人}.
    \textbf{Correct :} \zh{我家有四口人}. \textit{Note :} \cle{位} (wèi) honorifiques (\zh{三位老师}), \cle{个} (gè) générique mais moins naturel pour « taille foyer ».
\end{enumerate}
\textbf{Conseil final :} Pour l'exposé final, préparez \textbf{fiches bristol} : 1. Plan (mots-clés), 2. 3 phrases complexes types (\cle{因为...所以...}, \cle{虽然...但是...}, \cle{一旦...就...}), 3. 2 chengyu (\zh{归心似箭}, \zh{阖家幸福}), 4. Vocabulaire « pièges tons » surlignés. Entraînez-vous devant miroir / enregistrement.
\end{attention}

\begin{exercice}[Entraînement final : Préparation de l'exposé noté]
\textbf{Consigne :} Rédigez votre exposé personnel (1 min 30) sur le thème : \zh{「我的家庭：爱与团结」}. Contraintes obligatoires :
\begin{enumerate}
    \item Accroche + Composition famille (\cle{口}).
    \item 4 membres + 4 adjectifs qualités (dont 1 \cle{孝顺} ou \cle{体贴} ou \cle{勤劳}).
    \item 1 anecdote solidarité passée (\cle{了} ou \cle{过}) avec \cle{一起}, \cle{帮助}, \cle{鼓励} ou \cle{照顾}.
    \item 1 phrase concession (\cle{虽然...但是...}).
    \item Conclusion avec 1 chengyu (\zh{港湾}, \zh{后盾}, \zh{归心似箭}, \zh{阖家幸福}).
    \item Signature orale polie.
\end{enumerate}
\textbf{Remise :} Version caractères + pinyin (tons marqués) + traduction FR. Entraînement oral en binôme (chronomètre + grille évaluation : prononciation, structures, vocabulaire, fluidité, interaction).
\end{exercice}

\begin{corrige}[Exercice : Exposé final]
\textbf{Grille d'auto-évaluation / Évaluation par les pairs (sur 20) :}
\begin{itemize}
    \item \textbf{Prononciation \& Tons (5 pts) :} Tons neutres kinship respectés ? Sandhi \cle{不/一} corrects ? Tons pièges (\cle{xiàoshùn}, \cle{qínláo}, \cle{tǐtiē}, \cle{hémù}) justes ? Intonation phrase naturelle.
    \item \textbf{Structures grammaticales (5 pts) :} \cle{因为...所以...} / \cle{虽然...但是...} / \cle{一旦...就...} utilisés correctement ? Ordre mots (Temps/Lieu avant V) respecté ? \cle{led/guò/zhe} appropriés ? Classificateur \cle{口} utilisé ?
    \item \textbf{Vocabulaire \& Chengyu (4 pts) :} Richesse (adjectifs variés, verbes d'action précis) ? 1 chengyu minimum bien placé ? Registre soutenu (\cle{尊敬的}, \cle{二老}, \cle{拜年}) si contexte formel ?
    \item \textbf{Cohérence \& Contenu (3 pts) :} Structure logique (Présent $\rightarrow$ Passé $\rightarrow$ Sentiment) ? Anecdote concrète et touchante ? Respect temps imparti (1'30 $\pm$ 15'') ?
    \item \textbf{Présentation / Interaction (3 pts) :} Regard public, pas de lecture intégrale, voix claire, capacité à répondre aux 2 questions du jury (ex: \zh{你最敬佩谁？为什么？} / \zh{如果你有矛盾怎么解决？}).
\end{itemize}
\textbf{Modèle de fiche bristol (recto) :}
\begin{center}
\begin{tabularx}{\linewidth}{|X|}
\hline \rowcolor{popblueL} \textbf{MON EXPOSÉ : « 我的家庭：爱与团结 »} \\ \hline
\textbf{1. Accroche :} 各位同学，大家好。今天我想谈谈我的家庭。 \\ \hline
\textbf{2. Composition :} 我家有 \_\_\_ 口人：\_\_\_、\_\_\_、\_\_\_、\_\_\_ 和我。 \\ \hline
\textbf{3. Qualités :} \_\_\_ 很 \_\_\_；\_\_\_ 很 \_\_\_；\_\_\_ 很 \_\_\_；\_\_\_ 很 \_\_\_。 \\ \hline
\textbf{4. Anecdote (Passé) :} \_\_\_年\_\_\_，我\_\_\_。全家人\_\_\_。\_\_\_帮我\_\_\_，\_\_\_给我\_\_\_。 \\ \hline
\textbf{5. Concession :} 虽然\_\_\_，但是\_\_\_。 \\ \hline
\textbf{6. Conclusion + Chengyu :} 我觉得家是\_\_\_。\_\_\_（成语）。我爱我的家！ \\ \hline
\textbf{7. Mots pièges TONS :} \zh{xiàoshùn (4-4)} | \zh{qínláo (2-2)} | \zh{tǐtiē (3-1)} | \zh{yéye (2-0)} | \zh{nǎinai (3-0)} | \zh{bú kèqi (2-4-4)} | \zh{yìqǐ (4-3)} \\ \hline
\end{tabularx}
\end{center}
\textbf{Verso :} Traduction française de secours + 2 questions anticipées pour le jury.
\end{corrige}

\newpage

\coursec{Exercices}

\courssub{Je vérifie mes connaissances}

\begin{exercice}[QCM : Vocabulaire de la famille et de la solidarité]
Choisis la bonne réponse parmi les propositions A, B, C.
\begin{enumerate}
    \item Comment dit-on « grand-père paternel » en chinois ?
    \begin{itemize}
        \item A. 外公 (wàigōng)
        \item B. 爷爷 (yéye)
        \item C. 爸爸 (bàba)
    \end{itemize}
    \item Quel mot signifie « prendre soin de / s'occuper de » ?
    \begin{itemize}
        \item A. 帮助 (bāngzhù)
        \item B. 照顾 (zhàogù)
        \item C. 关心 (guānxīn)
    \end{itemize}
    \item Quelle structure exprime « Je pense que la famille est très importante » ?
    \begin{itemize}
        \item A. 我觉得家庭很重要。 (Wǒ juéde jiātíng hěn zhòngyào.)
        \item B. 我喜欢家庭很重要。 (Wǒ xǐhuan jiātíng hěn zhòngyào.)
        \item C. 我要家庭很重要。 (Wǒ yào jiātíng hěn zhòngyào.)
    \end{itemize}
    \item Pour dire « Parce que je aime ma famille, je veux les aider », on utilise :
    \begin{itemize}
        \item A. 所以 (suǒyǐ)
        \item B. 因为 (yīnwèi)
        \item C. 但是 (dànshì)
    \end{itemize}
\end{enumerate}
\end{exercice}

\begin{exercice}[Vrai ou Faux : Compréhension d'un dialogue]
Écoute (ou lis) le dialogue suivant et coche V (Vrai) ou F (Faux). Justifie ta réponse en français en te basant sur le texte.

\begin{textebox}[Dialogue : Qui fait la cuisine ce soir ?]
A: 妹妹，今天谁做饭？ (Mèimei, jīntiān shéi zuòfàn ?)
B: 哥哥做饭。我洗碗。 (Gēge zuòfàn. Wǒ xǐwǎn.)
A: 爸爸呢？ (Bàba ne ?)
B: 爸爸很忙，他在工作。 (Bàba hěn máng, tā zài gōngzuò.)
A: 我们一起帮助妈妈。 (Wǒmen yīqǐ bāngzhù māma.)
B: 好主意！ (Hǎo zhǔyì !)
\end{textebox}

\begin{enumerate}
    \item Le grand-père fait la cuisine ce soir. \hfill [ ] V \quad [ ] F
    \item La sœur cadette lave la vaisselle. \hfill [ ] V \quad [ ] F
    \item Le père est occupé au travail. \hfill [ ] V \quad [ ] F
    \item Le frère aîné propose d'aider la mère. \hfill [ ] V \quad [ ] F
\end{enumerate}
\end{exercice}

\begin{exercice}[Vocabulaire : Associer les paires]
Relie chaque mot chinois à sa définition française et à son pinyin correct (attention aux tons).

\begin{center}
\begin{tabularx}{\linewidth}{|X|X|X|X|}
\hline
\rowcolor{popblueL} \textbf{Caractères} & \textbf{Pinyin} & \textbf{Nature} & \textbf{Traduction} \\
\hline
团结 & tuánjié & verbe & s'unir, se serrer les coudes \\
\hline
孝顺 & xiàoshùn & adj. & filial, respectueux envers les parents \\
\hline
支持 & zhīchí & verbe & soutenir, appuyer \\
\hline
争吵 & zhēngchǎo & verbe & se disputer, se quereller \\
\hline
和睦 & hémù & adj. & harmonieux, en bons termes \\
\hline
\end{tabularx}
\end{center}

\textbf{Consigne :} Écris le numéro de la définition correspondante devant chaque mot.
1. Se disputer \quad 2. Soutenir \quad 3. Filial \quad 4. Harmonieux \quad 5. S'unir
\end{exercice}

\begin{exercice}[Pinyin et Tons : Discrimination auditive]
Associe chaque syllabe au bon caractère chinois en entendant la différence de ton (ou en lisant le pinyin). Entoure le bon caractère.

\begin{enumerate}
    \item \zh{jiā} (1er ton) \quad A. 假 (faux) \quad B. 家 (famille/maison) \quad C. 价 (prix)
    \item \zh{ài} (4e ton) \quad A. 爱 (aimer) \quad B. 艾 (armoise) \quad C. 唉 (soupirer)
    \item \zh{zhàogù} (4e + 4e ton) \quad A. 照顾 (s'occuper de) \quad B. 找古 (chercher l'ancien) \quad C. 罩鼓 (couvrir le tambour)
    \item \zh{xiōngdì} (1er + 4e ton) \quad A. 凶敌 (ennemi féroce) \quad B. 兄弟 (frères) \quad C. 胸弟 (poitrine jeune frère)
    \item \zh{jiějie} (3e + ton neutre) \quad A. 解解 (détacher) \quad B. 姐姐 (sœur aînée) \quad C. 借借 (emprunter)
\end{enumerate}
\end{exercice}

\courssub{Je m'entraîne}

\begin{exercice}[Traduction : Phrases à traduire en chinois]
Traduis les phrases suivantes en caractères chinois (avec pinyin). Utilise le vocabulaire de la leçon (famille, solidarité, formules d'opinion).

\begin{enumerate}
    \item Ma famille est très unie, nous nous aimons beaucoup.
    \item Je pense que prendre soin de ses parents, c'est important.
    \item Mon grand-père est malade, ma mère s'occupe de lui à l'hôpital.
    \item Bien que nous nous disputions parfois, nous nous soutenons toujours.
    \item À ton avis, qui doit faire la cuisine ce soir ?
\end{enumerate}
\end{exercice}

\begin{exercice}[Transformation : Restructuration de phrases]
Transforme les phrases suivantes selon le modèle indiqué. Écris le résultat en caractères + pinyin.

\begin{enumerate}
    \item \textbf{Modèle « 因为…所以… » :} Combine les deux phrases.
    \begin{itemize}
        \item Phrase 1 : 我爱我的家庭。 (Wǒ ài wǒ de jiātíng.)
        \item Phrase 2 : 我每天帮助父母。 (Wǒ měitiān bāngzhù fùmǔ.)
        \item $\rightarrow$ \underline{\hspace{8cm}}
    \end{itemize}
    \item \textbf{Modèle « 虽然…但是… » :} Exprime la concession.
    \begin{itemize}
        \item Idée : Nous sommes pauvres / Nous sommes heureux.
        \item $\rightarrow$ \underline{\hspace{8cm}}
    \end{itemize}
    \item \textbf{Modèle « 不但…而且… » :} Ajoute une information positive.
    \begin{itemize}
        \item Base : 妹妹很孝顺。 (Mèimei hěn xiàoshùn.)
        \item Ajout : Elle est très polie. (很有礼貌 / hěn yǒu lǐmào)
        \item $\rightarrow$ \underline{\hspace{8cm}}
    \end{itemize}
    \item \textbf{Question ouverte avec « 怎么 » ou « 为什么 » :}
    \begin{itemize}
        \item Réponse : 因为我想照顾生病的奶奶。 (Yīnwèi wǒ xiǎng zhàogù shēngbìng de nǎinai.)
        \item Question : \underline{\hspace{8cm}}
    \end{itemize}
\end{enumerate}
\end{exercice}

\begin{exercice}[Texte à trous : Dialogue à compléter]
Complète le dialogue suivant avec les formules fonctionnelles ou le vocabulaire approprié (caractères chinois). Le pinyin est donné comme indice.

\begin{textebox}[Dialogue : Préparer la fête du Printemps]
A: 春节快到了，\underline{\hspace{3cm}} (wǒmen yīqǐ zuò shénme) ?
B: \underline{\hspace{3cm}} (wǒ juéde) 我们应该打扫卫生，\underline{\hspace{3cm}} (yěyào) 贴春联。
A: \underline{\hspace{3cm}} (hǎo zhǔyì)！爸爸做饭，妈妈包饺子，\underline{\hspace{3cm}} (wǒmen) 做什么？
B: 我们洗碗和扫地。\underline{\hspace{3cm}} (nǐ ne) ?
A: 我\underline{\hspace{3cm}} (tóngyì)。全家\underline{\hspace{3cm}} (tuánjié) 一心，其利断金！
B: 说得好！\underline{\hspace{3cm}} (xièxie dàjiā)！
\end{textebox}

\textbf{Mots-outils à placer :} 你呢？(nǐ ne) / 我觉得 (wǒ juéde) / 好主意 (hǎo zhǔyì) / 我也要 (wǒ yě yào) / 团结 (tuánjié) / 我们一起做什么 (wǒmen yīqǐ zuò shénme) / 同意 (tóngyì) / 谢谢大家 (xièxie dàjiā)
\end{exercice}

\begin{exercice}[Remise en ordre : Étapes d'un exposé oral]
Remets les 5 étapes suivantes dans l'ordre logique pour présenter un exposé oral structuré sur « Ma famille et moi ». Numérote-les de 1 à 5. Puis, pour chaque étape, écris une phrase-type en chinois (caractères + pinyin) que tu pourrais dire.

\begin{enumerate}
    \item \_\_\_ Conclure et remercier l'auditoire.
    \item \_\_\_ Présenter les membres de sa famille (qui, âge, métier).
    \item \_\_\_ Saluer et annoncer le sujet.
    \item \_\_\_ Développer un exemple de solidarité (raconter une petite histoire).
    \item \_\_\_ Exprimer son opinion / sentiment sur l'importance de la famille.
\end{enumerate}
\end{exercice}

\courssub{Je me prépare au Bac}

\begin{exercice}[Compréhension de l'écrit : Texte original — « Le dimanche à Yaoundé »]
Lis le texte suivant et réponds aux questions en français (sauf indication contraire).

\begin{textebox}[Le dimanche à Yaoundé]
我叫大卫，住在喀麦隆的雅温得。每个星期天，我们全家都很忙，但也很快乐。
早上八点，爸爸去市场买菜，妈妈和奶奶在厨房做饭。我和妹妹洗蔬菜，哥哥扫地。中午十二点，我们围坐在大桌子旁吃饭。桌上有恩多莱、香蕉和米饭，很好吃。
吃饭后，爸爸说：“今天我有点累，谁帮我按摩一下肩膀？”我说：“我来！”妹妹说：“我给爸爸倒水。”奶奶笑着说：“我们这个家真团结，真和睦。”
下午，我们一起看足球比赛，为喀麦隆队加油。晚上，爸爸给我们讲中国的故事，说中国人也很重视家庭团聚。我觉得无论在雅温得还是在北京，家人的爱和互相照顾都是最宝贵的财富。
(Dàwèi, Kāmàilóng, Yāwēndé, měige xīngqītiān, wéizuò, ēnduǒlái, ānmó, tuánjié, hémù, bǎoguì de cáifù)
\end{textebox}

\textbf{Questions de compréhension :}
\begin{enumerate}
    \item Où habite David ? Quel jour la famille est-elle décrite comme occupée ?
    \item Qui fait quoi le matin ? Cite au moins 3 activités différentes réparties entre les membres de la famille.
    \item Qu'est-ce qu'on mange le midi ? Donne les noms des plats en chinois (caractères) et en français.
    \item Comment le père exprime-t-il sa fatigue ? Quelle est la réaction immédiate de David et de sa sœur ?
    \item Quelle comparaison culturelle fait le père le soir ? Quel est le sentiment final de David (dernière phrase) ?
\end{enumerate}

\textbf{Questions de langue :}
\begin{enumerate}
    \setcounter{enumi}{5}
    \item Repère dans le texte deux structures « 因为…所以… » ou « 虽然…但是… » (ou équivalent logique) et recopie-les.
    \item Trouve 3 verbes d'action quotidiens (ex: 买, 做, 洗) et 2 adjectifs d'état/d'humeur.
    \item Explique la fonction de « 都 » (dōu) dans « 我们全家都很忙 » et de « 一起 » (yīqǐ) dans « 我们一起看 ».
\end{enumerate}
\end{exercice}

\begin{exercice}[Thème et Version : Traduction dirigée]
\textbf{A. Thème (Français $\rightarrow$ Chinois) :} Traduis le paragraphe suivant en chinois (caractères + pinyin). Environ 60-80 caractères.
\begin{quote}
« Chez moi, la solidarité familiale est une règle d'or. Quand ma grand-mère est tombée malade le mois dernier, mon père a pris congé pour l'accompagner à l'hôpital. Ma mère cuisinait des repas nutritifs tous les jours. Moi, j'aidais ma grand-mère à prendre ses médicaments et je lui lisais des histoires. Mes frères et sœurs nettoyaient la maison. Bien que nous soyons fatigués, personne ne s'est plaint. Je pense que c'est ça, l'amour familial : se soutenir dans les moments difficiles. »
\end{quote}

\textbf{B. Version (Chinois $\rightarrow$ Français) :} Traduis le texte suivant en français naturel.
\begin{zhbox}
现在的年轻人很忙，工作压力大，常常没时间陪父母。但是，孝顺不一定要天天在身边。一个电话、一个视频、一句“爸妈，你们吃饭了吗？”，都能让父母感到温暖。我也决定，以后每个周末给家里打电话，多关心关心他们的身体。家永远是我们最温暖的港湾。
\end{zhbox}
\end{exercice}

\begin{exercice}[Expression écrite / Orale : Sujet type Baccalauréat]
\textbf{Sujet :} « \textbf{家人最重要} » (La famille, c'est le plus important).
Tu dois préparer un exposé oral de 2 à 3 minutes (équivalent à 120-150 caractères chinois) pour l'épreuve orale du Baccalauréat.

\textbf{Consignes :}
\begin{enumerate}
    \item Rédige d'abord le texte écrit de ton exposé (brouillon structuré) en respectant la structure apprise : Salutation/Sujet $\rightarrow$ Présentation famille $\rightarrow$ Exemple concret de solidarité (vocabulaire : 照顾, 帮助, 团结, 孝顺) $\rightarrow$ Opinion personnelle avec \zh{因为}...\zh{所以}... / \zh{我觉得}... $\rightarrow$ Conclusion/Remerciements.
    \item Utilise au moins 5 formules fonctionnelles vues en classe (ex: \zh{首先}, \zh{其次}, \zh{比如说}, \zh{总之}, \zh{谢谢大家}).
    \item Soigne les tons (marque-les sur le pinyin final).
    \item Prépare 2 arguments POUR et 2 arguments CONTRE l'idée « La famille passe avant tout » (travail, études, indépendance) pour anticiper le débat avec le jury.
\end{enumerate}

\begin{center}
\begin{tabularx}{\linewidth}{|X|X|}
\hline
\rowcolor{popblueL} \textbf{Arguments POUR (家人最重要)} & \textbf{Arguments CONTRE (个人发展也重要)} \\
\hline
1. \zh{因为家人永远支持我……} & 1. \zh{虽然家人重要，但是……} \\
\hline
2. & 2. \\
\hline
\end{tabularx}
\end{center}

\textbf{Espace de rédaction (brouillon) :}
\vspace{4cm}
\end{exercice}

\newpage

\coursec{Corrigés des exercices}

\begin{corrige}[Exercice 1 — QCM : Vocabulaire de la famille et de la solidarité]
\begin{enumerate}
    \item \textbf{B.} \zh{爷爷} \emph{(yéye)} = grand-père paternel. \zh{外公} \emph{(wàigōng)} désigne le grand-père \emph{maternel} ; \zh{爸爸} \emph{(bàba)} = papa.
    \item \textbf{B.} \zh{照顾} \emph{(zhàogù)} = prendre soin de, s'occuper de (soins concrets : nourrir, accompagner). \zh{帮助} \emph{(bāngzhù)} = aider (action ponctuelle) ; \zh{关心} \emph{(guānxīn)} = se soucier de (sentiment).
    \item \textbf{A.} \zh{我觉得家庭很重要。} \emph{(Wǒ juéde jiātíng hěn zhòngyào.)} « Je pense que la famille est très importante. » \zh{觉得} est la formule d'opinion ; B et C sont grammaticalement fautives (\zh{喜欢} et \zh{要} ne peuvent pas introduire ici une proposition complète).
    \item \textbf{B.} \zh{因为} \emph{(yīnwèi)} = parce que (cause). Structure : \zh{因为我爱我的家人，所以我想帮助他们。} \emph{(Yīnwèi wǒ ài wǒ de jiārén, suǒyǐ wǒ xiǎng bāngzhù tāmen.)} \zh{所以} \emph{(suǒyǐ)} = donc (conséquence) ; \zh{但是} \emph{(dànshì)} = mais (opposition).
\end{enumerate}
\textbf{Point de vigilance :} ne pas confondre la branche paternelle (\zh{爷爷} / \zh{奶奶}) et la branche maternelle (\zh{外公} / \zh{外婆}), distinction obligatoire en chinois.
\end{corrige}

\begin{corrige}[Exercice 2 — Vrai ou Faux : Qui fait la cuisine ce soir ?]
\begin{enumerate}
    \item \textbf{F.} Le grand-père n'apparaît pas dans le dialogue ; c'est le frère aîné qui cuisine : \zh{哥哥做饭} \emph{(Gēge zuòfàn)}.
    \item \textbf{V.} \zh{我洗碗} \emph{(Wǒ xǐwǎn)} : c'est la sœur cadette (\zh{妹妹}, interlocuteur B) qui dit « je lave la vaisselle ».
    \item \textbf{V.} \zh{爸爸很忙，他在工作。} \emph{(Bàba hěn máng, tā zài gōngzuò.)} = « Papa est très occupé, il est en train de travailler. »
    \item \textbf{V.} \zh{我们一起帮助妈妈。} \emph{(Wǒmen yīqǐ bāngzhù māma.)} : c'est A (qui s'adresse à \zh{妹妹}, donc l'aîné) qui propose d'aider la mère ensemble.
\end{enumerate}
\textbf{Point de vigilance :} repérer qui parle grâce aux termes d'adresse (\zh{妹妹} = on parle \emph{à} la petite sœur) ; \zh{呢} \emph{(ne)} sert à relancer la question : « et papa ? ».
\end{corrige}

\begin{corrige}[Exercice 3 — Vocabulaire : Associer les paires]
\begin{center}
\begin{tabularx}{\linewidth}{|X|X|X|}
\hline
\rowcolor{popblueL} \textbf{Caractères} & \textbf{Pinyin} & \textbf{N° de définition} \\
\hline
团结 & \emph{tuánjié} & 5 (s'unir) \\
\hline
孝顺 & \emph{xiàoshùn} & 3 (filial) \\
\hline
支持 & \emph{zhīchí} & 2 (soutenir) \\
\hline
争吵 & \emph{zhēngchǎo} & 1 (se disputer) \\
\hline
和睦 & \emph{hémù} & 4 (harmonieux) \\
\hline
\end{tabularx}
\end{center}
\textbf{Point de vigilance :} attention aux tons de \zh{孝顺} \emph{(xiàoshùn, 4e + 4e)} et à ne pas confondre \zh{支持} \emph{(zhīchí, soutenir moralement)} avec \zh{帮助} \emph{(bāngzhù, aider concrètement)}.
\end{corrige}

\begin{corrige}[Exercice 4 — Pinyin et Tons : Discrimination auditive]
\begin{enumerate}
    \item \textbf{B.} \zh{家} \emph{(jiā, 1er ton)} = famille/maison. \zh{假} \emph{(jiǎ, 3e ton)} = faux ; \zh{价} \emph{(jià, 4e ton)} = prix.
    \item \textbf{A.} \zh{爱} \emph{(ài, 4e ton)} = aimer. \zh{唉} \emph{(āi, 1er ton)} = soupirer.
    \item \textbf{A.} \zh{照顾} \emph{(zhàogù)} = s'occuper de. B et C sont des suites de syllabes sans sens courant, pièges homophones.
    \item \textbf{B.} \zh{兄弟} \emph{(xiōngdì)} = frères. \zh{凶} \emph{(xiōng)} = féroce est un homophone du 1er ton : seul le contexte (et l'écrit) permet de trancher.
    \item \textbf{B.} \zh{姐姐} \emph{(jiějie)} = sœur aînée, avec deuxième syllabe au ton neutre, typique des redoublements de parenté (\zh{妈妈} \emph{(māma)}, \zh{哥哥} \emph{(gēge)}, \zh{妹妹} \emph{(mèimei)}).
\end{enumerate}
\textbf{Point de vigilance :} dans les mots de famille redoublés, la deuxième syllabe perd son ton (ton neutre) : on dit \emph{jiějie}, jamais \emph{jiějiě}.
\end{corrige}

\begin{corrige}[Exercice 5 — Traduction : Phrases à traduire en chinois]
\begin{enumerate}
    \item \zh{我的家庭很团结，我们非常相爱。} \emph{(Wǒ de jiātíng hěn tuánjié, wǒmen fēicháng xiāng'ài.)}
    \item \zh{我觉得照顾父母很重要。} \emph{(Wǒ juéde zhàogù fùmǔ hěn zhòngyào.)}
    \item \zh{我爷爷生病了，我妈妈在医院照顾他。} \emph{(Wǒ yéye shēngbìng le, wǒ māma zài yīyuàn zhàogù tā.)}
    \item \zh{虽然我们有时候争吵，但是我们总是互相支持。} \emph{(Suīrán wǒmen yǒu shíhou zhēngchǎo, dànshì wǒmen zǒngshì hùxiāng zhīchí.)}
    \item \zh{你觉得今天谁做饭？} \emph{(Nǐ juéde jīntiān shéi zuòfàn?)} ou \zh{在你看来，今天晚上谁应该做饭？} \emph{(Zài nǐ kànlái, jīntiān wǎnshang shéi yīnggāi zuòfàn?)}
\end{enumerate}
\textbf{Points de vigilance :} (a) pas de \zh{是} devant un adjectif : \zh{很重要} et non *\zh{是很重要} ; (b) \zh{虽然} et \zh{但是} peuvent coexister dans la même phrase en chinois, contrairement au français ; (c) \zh{互相} \emph{(hùxiāng, mutuellement)} se place avant le verbe.
\end{corrige}

\begin{corrige}[Exercice 6 — Transformation : Restructuration de phrases]
\begin{enumerate}
    \item \zh{因为我爱我的家庭，所以我每天帮助父母。} \emph{(Yīnwèi wǒ ài wǒ de jiātíng, suǒyǐ wǒ měitiān bāngzhù fùmǔ.)} — Cause puis conséquence ; virgule chinoise entre les deux propositions.
    \item \zh{虽然我们很穷，但是我们很幸福。} \emph{(Suīrán wǒmen hěn qióng, dànshì wǒmen hěn xìngfú.)} — Concession : \zh{虽然} en tête, \zh{但是} dans la seconde proposition.
    \item \zh{妹妹不但很孝顺，而且很有礼貌。} \emph{(Mèimei búdàn hěn xiàoshùn, érqiě hěn yǒu lǐmào.)} — Addition : \zh{不但…而且…} ; attention, \zh{不} passe au 2e ton devant un 4e ton : \emph{búdàn}.
    \item \zh{你为什么想留在家里？} \emph{(Nǐ wèishénme xiǎng liú zài jiā lǐ?)} ou \zh{你为什么不去玩？} \emph{(Nǐ wèishénme bú qù wán?)} — La réponse commence par \zh{因为}, donc la question doit contenir \zh{为什么}.
\end{enumerate}
\textbf{Point de vigilance :} \zh{因为…所以…} et \zh{虽然…但是…} s'emploient en paires en chinois ; ne pas traduire le français « parce que » seul par \zh{因为} seul quand la conséquence suit.
\end{corrige}

\begin{corrige}[Exercice 7 — Texte à trous : Préparer la fête du Printemps]
Dialogue complété :

A : \zh{春节快到了，我们一起做什么？} \emph{(Chūnjié kuài dào le, wǒmen yīqǐ zuò shénme?)}

B : \zh{我觉得我们应该打扫卫生，也要贴春联。} \emph{(Wǒ juéde wǒmen yīnggāi dǎsǎo wèishēng, yě yào tiē chūnlián.)}

A : \zh{好主意！爸爸做饭，妈妈包饺子，我们做什么？} \emph{(Hǎo zhǔyì! Bàba zuòfàn, māma bāo jiǎozi, wǒmen zuò shénme?)}

B : \zh{我们洗碗和扫地。你呢？} \emph{(Wǒmen xǐwǎn hé sǎodì. Nǐ ne?)}

A : \zh{我同意。全家团结一心，其利断金！} \emph{(Wǒ tóngyì. Quánjiā tuánjié yīxīn, qí lì duàn jīn!)}

B : \zh{说得好！谢谢大家！} \emph{(Shuō de hǎo! Xièxie dàjiā!)}

\textbf{Point de vigilance :} \zh{也要} \emph{(yě yào)} = « il faut aussi » ; le proverbe \zh{团结一心，其利断金} signifie « unis d'un même cœur, on peut briser le métal » — l'équivalent chinois de « l'union fait la force ».
\end{corrige}

\begin{corrige}[Exercice 8 — Remise en ordre : Étapes d'un exposé oral]
Ordre logique : 1. Saluer et annoncer le sujet ; 2. Présenter les membres de sa famille ; 3. Développer un exemple de solidarité ; 4. Exprimer son opinion ; 5. Conclure et remercier.

Phrases-types :
\begin{enumerate}
    \item \zh{大家好！今天我说一说我的家庭。} \emph{(Dàjiā hǎo! Jīntiān wǒ shuō yi shuō wǒ de jiātíng.)}
    \item \zh{我家有四口人：爸爸、妈妈、妹妹和我。} \emph{(Wǒ jiā yǒu sì kǒu rén: bàba, māma, mèimei hé wǒ.)}
    \item \zh{比如说，上个月奶奶生病了，我们一起照顾她。} \emph{(Bǐrú shuō, shàng ge yuè nǎinai shēngbìng le, wǒmen yīqǐ zhàogù tā.)}
    \item \zh{我觉得家人最重要，因为他们永远支持我。} \emph{(Wǒ juéde jiārén zuì zhòngyào, yīnwèi tāmen yǒngyuǎn zhīchí wǒ.)}
    \item \zh{我的介绍完了，谢谢大家！} \emph{(Wǒ de jièshào wán le, xièxie dàjiā!)}
\end{enumerate}
\textbf{Point de vigilance :} le classificateur des membres de la famille est \zh{口} \emph{(kǒu)} : \zh{四口人}, jamais *\zh{四个人} pour compter les bouches d'un foyer.
\end{corrige}

\begin{corrige}[Exercice 9 — Compréhension de l'écrit : « Le dimanche à Yaoundé »]
\textbf{Questions de compréhension :}
\begin{enumerate}
    \item David habite à Yaoundé, au Cameroun (\zh{住在喀麦隆的雅温得} \emph{(zhù zài Kāmàilóng de Yǎwēndé)}). La famille est décrite comme occupée le dimanche (\zh{每个星期天} \emph{(měi ge xīngqítiān)}).
    \item Le matin : le père va au marché acheter des légumes (\zh{爸爸去市场买菜} \emph{(bàba qù shìchǎng mǎi cài)}) ; la mère et la grand-mère cuisinent dans la cuisine (\zh{妈妈和奶奶在厨房做饭} \emph{(māma hé nǎinai zài chúfáng zuòfàn)}) ; David et sa sœur lavent les légumes (\zh{我和妹妹洗蔬菜} \emph{(wǒ hé mèimei xǐ shūcài)}) ; le frère aîné balaie (\zh{哥哥扫地} \emph{(gēge sǎodì)}).
    \item On mange du ndolé (\zh{恩多莱} \emph{(Ēnduǒlái)}), des bananes plantains (\zh{香蕉} \emph{(xiāngjiāo)}) et du riz (\zh{米饭} \emph{(mǐfàn)}).
    \item Le père dit : \zh{今天我有点累，谁帮我按摩一下肩膀？} \emph{(Jīntiān wǒ yǒudiǎn lèi, shéi bāng wǒ ànmó yíxià jiānbǎng?)} (« Aujourd'hui je suis un peu fatigué, qui peut me masser les épaules ? »). David répond aussitôt \zh{我来！} \emph{(Wǒ lái!)} (« Je m'en charge ! ») et sa sœur \zh{我给爸爸倒水} \emph{(Wǒ gěi bàba dào shuǐ)} (« Je verse de l'eau à papa »).
    \item Le soir, le père raconte des histoires de Chine et explique que les Chinois attachent aussi beaucoup d'importance à la réunion familiale (\zh{中国人也很重视家庭团聚} \emph{(Zhōngguórén yě hěn zhòngshì jiātíng tuánjù)}). Sentiment final de David : que ce soit à Yaoundé ou à Pékin, l'amour familial et l'entraide sont le trésor le plus précieux (\zh{最宝贵的财富} \emph{(zuì bǎoguì de cáifù)}).
\end{enumerate}
\textbf{Questions de langue :}
\begin{enumerate}
    \setcounter{enumi}{5}
    \item Le texte ne contient pas la paire explicite, mais deux équivalents logiques : \zh{我们全家都很忙，但也很快乐} \emph{(wǒmen quánjiā dōu hěn máng, dàn yě hěn kuàilè)} (concession avec \zh{但} \emph{(dàn)}) et \zh{无论在雅温得还是在北京，家人的爱…都是最宝贵的财富} \emph{(wúlùn zài Yǎwēndé háishì zài Běijīng, jiārén de ài… dōu shì zuì bǎoguì de cáifù)} (concession avec \zh{无论…都…} \emph{(wúlùn… dōu…)}).
    \item Verbes d'action : \zh{买} \emph{(mǎi, acheter)}, \zh{做} \emph{(zuò, faire)}, \zh{洗} \emph{(xǐ, laver)}, \zh{扫} \emph{(sǎo, balayer)}, \zh{讲} \emph{(jiǎng, raconter)}. Adjectifs : \zh{忙} \emph{(máng, occupé)}, \zh{快乐} \emph{(kuàilè, joyeux)}, \zh{累} \emph{(lèi, fatigué)}.
    \item \zh{都} \emph{(dōu)} = « tous » : il résume le sujet pluriel \zh{我们全家} et se place avant l'adjectif/verbe. \zh{一起} \emph{(yīqǐ)} = « ensemble » : adverbe placé avant le verbe, il insiste sur l'action collective, cœur du thème de la solidarité.
\end{enumerate}
\textbf{Barème indicatif (20 pts) :} compréhension 10 pts (2 pts par question), questions de langue 6 pts (2 pts chacune), qualité du français 4 pts.
\end{corrige}

\begin{corrige}[Exercice 10 — Thème et Version : Traduction dirigée]
\textbf{A. Thème — proposition de traduction :}

\zh{在我家，家庭团结是一条黄金规则。上个月奶奶生病了，爸爸请假陪她去医院。妈妈每天做有营养的饭。我帮助奶奶吃药，给她讲故事。我的哥哥姐姐打扫房子。虽然我们很累，但是没有人抱怨。我觉得这就是家人的爱：在困难的时候互相支持。}

\emph{(Wǒ jiā, jiātíng tuánjié shì yì tiáo huángjīn guīzé. Shàng ge yuè nǎinai shēngbìng le, bàba qǐngjià péi tā qù yīyuàn. Māma měitiān zuò yǒu yíngyǎng de fàn. Wǒ bāngzhù nǎinai chīyào, gěi tā jiǎng gùshi. Wǒ de gēge jiějie dǎsǎo fángzi. Suīrán wǒmen hěn lèi, dànshì méiyǒu rén bàoyuàn. Wǒ juéde zhè jiùshì jiārén de ài: zài kùnnan de shíhou hùxiāng zhīchí.)}

\textbf{Points de vigilance :} « tomber malade » = \zh{生病} + \zh{了} \emph{(le)} (changement d'état) ; « prendre congé » = \zh{请假} \emph{(qǐngjià)} ; « se plaindre » = \zh{抱怨} \emph{(bàoyuàn)} ; « personne ne » = \zh{没有人} \emph{(méiyǒu rén)} + verbe.

\textbf{B. Version — proposition de traduction :}

« Les jeunes d'aujourd'hui sont très occupés, la pression du travail est forte, et ils n'ont souvent pas le temps d'accompagner leurs parents. Mais être filial ne signifie pas forcément être à leurs côtés tous les jours. Un coup de téléphone, une vidéo, une simple phrase "Papa, maman, vous avez mangé ?", tout cela peut leur apporter de la chaleur. Moi aussi, j'ai décidé d'appeler ma famille chaque week-end et de mieux prendre de leurs nouvelles, de leur santé. La famille restera toujours notre port le plus chaleureux. »

\textbf{Barème indicatif (20 pts) :} thème 10 pts (exactitude grammaticale 4, vocabulaire 4, caractères 2) ; version 10 pts (fidélité 6, qualité du français 4).
\end{corrige}

\begin{corrige}[Exercice 11 — Expression écrite / Orale : « 家人最重要 »]
\textbf{Proposition d'exposé modèle (environ 140 caractères) :}

\zh{大家好！今天我说一说"家人最重要"。首先，我介绍一下我的家：我家有五口人，爸爸、妈妈、爷爷、妹妹和我。其次，我想讲一个小故事。比如说，去年爷爷生病了，爸爸每天下班以后去医院照顾他，妈妈做饭，我和妹妹打扫房间。虽然我们都很累，但是没有人抱怨。我觉得家人最重要，因为无论我遇到什么困难，他们总是支持我、鼓励我。总之，家是我们最温暖的港湾。我的介绍完了，谢谢大家！}

\emph{(Dàjiā hǎo! Jīntiān wǒ shuō yi shuō "jiārén zuì zhòngyào". Shǒuxiān, wǒ jièshào yíxià wǒ de jiā: wǒ jiā yǒu wǔ kǒu rén, bàba, māma, yéye, mèimei hé wǒ. Qícì, wǒ xiǎng jiǎng yí ge xiǎo gùshi. Bǐrú shuō, qùnián yéye shēngbìng le, bàba měitiān xiàbān yǐhòu qù yīyuàn zhàogù tā, māma zuòfàn, wǒ hé mèimei dǎsǎo fángjiān. Suīrán wǒmen dōu hěn lèi, dànshì méiyǒu rén bàoyuàn. Wǒ juéde jiārén zuì zhòngyào, yīnwèi wúlùn wǒ yùdào shénme kùnnan, tāmen zǒngshì zhīchí wǒ, gǔlì wǒ. Zǒngzhī, jiā shì wǒmen zuì wēnnuǎn de gǎngwān. Wǒ de jièshào wán le, xièxie dàjiā!)}

\textbf{Arguments pour le débat :}
\begin{center}
\begin{tabularx}{\linewidth}{|X|X|}
\hline
\rowcolor{popblueL} \textbf{POUR (\zh{家人最重要})} & \textbf{CONTRE (\zh{个人发展也重要})} \\
\hline
1. \zh{因为家人永远支持我，所以我有勇气面对困难。} \emph{(Yīnwèi jiārén yǒngyuǎn zhīchí wǒ, suǒyǐ wǒ yǒu yǒngqì miànduì kùnnan.)} & 1. \zh{虽然家人重要，但是年轻人也需要独立。} \emph{(Suīrán jiārén zhòngyào, dànshì niánqīngrén yě xūyào dúlì.)} \\
\hline
2. \zh{家给我们温暖和爱，钱买不到这些。} \emph{(Jiā gěi wǒmen wēnnuǎn hé ài, qián mǎi bu dào zhèxiē.)} & 2. \zh{如果只想家人，我们可能放弃好的工作机会。} \emph{(Rúguǒ zhǐ xiǎng jiārén, wǒmen kěnéng fàngqì hǎo de gōngzuò jīhuì.)} \\
\hline
\end{tabularx}
\end{center}

\textbf{Barème indicatif (20 pts) :} structure respectée (salutation $\rightarrow$ conclusion) 4 pts ; 5 formules fonctionnelles employées 4 pts ; grammaire et connecteurs (\zh{因为…所以}, \zh{虽然…但是}) 5 pts ; vocabulaire de la leçon (\zh{照顾}, \zh{团结}, \zh{孝顺}, \zh{支持}) 4 pts ; prononciation et tons à l'oral 3 pts.

\textbf{Points de vigilance :} (a) \zh{首先 / 其次 / 总之} \emph{(shǒuxiān / qícì / zǒngzhī)} rythment l'exposé, ne pas les oublier ; (b) \zh{无论…都…} \emph{(wúlùn… dōu…)} exige \zh{都} dans la seconde proposition ; (c) compter les personnes avec \zh{口} \emph{(kǒu)} ; (d) à l'oral, soigner les 4e tons de \zh{最重要} \emph{(zuì zhòngyào)} et le ton neutre de \zh{妹妹} \emph{(mèimei)} ; (e) terminer toujours par \zh{谢谢大家} \emph{(xièxie dàjiā)} avec un sourire et une légère inclinaison de tête.
\end{corrige}

\newpage

\coursec{Fiche bilan}

\begin{popfigure}
\begin{center}\tcbox[colback=popblue,colframe=popblue,coltext=white,arc=9pt,boxsep=4pt,left=6pt,right=6pt]{\bfseries\small Amour de la famille \& Solidarité}\end{center}
\vspace{-2pt}
\begin{tcbraster}[raster columns=3,raster equal height=rows,raster column skip=6pt,raster row skip=6pt,enhanced,blankest]
\begin{tcolorbox}[enhanced,colback=poporange,colframe=poporange,coltext=white,boxrule=0.9pt,arc=3pt,left=4pt,right=4pt,top=3pt,bottom=3pt,boxsep=1pt,halign=left]\bfseries\scriptsize Vocabulaire clé 孝顺 团聚 家务 支持 传统 争吵\end{tcolorbox}
\begin{tcolorbox}[enhanced,colback=popgreen,colframe=popgreen,coltext=white,boxrule=0.9pt,arc=3pt,left=4pt,right=4pt,top=3pt,bottom=3pt,boxsep=1pt,halign=left]\bfseries\scriptsize Structures 比较句 把字句 连词 语气词\end{tcolorbox}
\begin{tcolorbox}[enhanced,colback=poppurple,colframe=poppurple,coltext=white,boxrule=0.9pt,arc=3pt,left=4pt,right=4pt,top=3pt,bottom=3pt,boxsep=1pt,halign=left]\bfseries\scriptsize Formules 我认为 一方面 我赞成 请问\end{tcolorbox}
\begin{tcolorbox}[enhanced,colback=poppink,colframe=poppink,coltext=white,boxrule=0.9pt,arc=3pt,left=4pt,right=4pt,top=3pt,bottom=3pt,boxsep=1pt,halign=left]\bfseries\scriptsize Tons \& Prononciation 3e ton sandhi 轻声 儿化音\end{tcolorbox}
\begin{tcolorbox}[enhanced,colback=popgold,colframe=popgold,coltext=white,boxrule=0.9pt,arc=3pt,left=4pt,right=4pt,top=3pt,bottom=3pt,boxsep=1pt,halign=left]\bfseries\scriptsize Jeux de rôle Aider aux devoirs Préparer fête\end{tcolorbox}
\begin{tcolorbox}[enhanced,colback=popdark,colframe=popdark,coltext=white,boxrule=0.9pt,arc=3pt,left=4pt,right=4pt,top=3pt,bottom=3pt,boxsep=1pt,halign=left]\bfseries\scriptsize Débat guidé Famille d'abord ? Aînés vs Jeunes\end{tcolorbox}
\begin{tcolorbox}[enhanced,colback=popblue!60,colframe=popblue!60,coltext=white,boxrule=0.9pt,arc=3pt,left=4pt,right=4pt,top=3pt,bottom=3pt,boxsep=1pt,halign=left]\bfseries\scriptsize Culture CMR/CHN Fête du Printemps Ngondo Respect aînés\end{tcolorbox}
\end{tcbraster}

\legende{Carte mentale récapitulative — Leçon 3 : Expression orale (Famille \& Solidarité)}
\end{popfigure}

\begin{bilan}

\end{bilan}

\courssub{Lexique clé (HSK 2--3) — Famille \& Solidarité}

\begin{center}
\begin{tabularx}{\linewidth}{|X|X|X|X|}
\hline
\rowcolor{popblueL} \textbf{Caractères} & \textbf{Pinyin} & \textbf{Nature} & \textbf{Traduction} \\ \hline
父母 & fùmǔ & n. & parents \\ \hline
兄弟姐妹 & xiōngdì jiěmèi & n. & frères et sœurs \\ \hline
孝顺 & xiàoshùn & adj./v. & filial, respectueux (envers parents) ; piété filiale \\ \hline
团聚 & tuánjù & v. & se réunir (en famille) \\ \hline
关心 & guānxīn & v. & se soucier de, s'inquiéter pour \\ \hline
支持 & zhīchí & v./n. & soutenir ; soutien \\ \hline
家务 & jiāwù & n. & tâches ménagères \\ \hline
争吵 & zhēngchǎo & v./n. & se disputer ; dispute \\ \hline
和解 & huéjiě & v. & se réconcilier \\ \hline
传统 & chuántǒng & n. & tradition \\ \hline
长辈 & zhǎngbèi & n. & aîné, aînés \\ \hline
晚辈 & wǎnbèi & n. & cadet, cadets \\ \hline
\end{tabularx}
\end{center}

\courssub{Structures grammaticales à connaître par cœur}

\begin{center}
\begin{tabularx}{\linewidth}{|X|X|X|}
\hline
\rowcolor{popblueL} \textbf{Structure / Règle} & \textbf{Exemple (Caractères + Pinyin)} & \textbf{Traduction} \\ \hline
比较句 (比) : A 比 B + Adj. & 哥哥 \textbf{比} 我 \textbf{高}。 (Gēge \textbf{bǐ} wǒ \textbf{gāo}.) & Mon frère est plus grand que moi. \\ \hline
把字句 (disposition) : S + 把 + O + V + 结果/方向 & 我 \textbf{把} 家务 \textbf{做完了}。 (Wǒ \textbf{bǎ} jiāwù \textbf{zuò wán le}.) & J'ai fini de faire les tâches ménagères. \\ \hline
被字句 (passif) : S + 被 + Agent + V & 碗 \textbf{被} 弟弟 \textbf{打破了}。 (Wǎn \textbf{bèi} dìdi \textbf{dǎ pò le}.) & Le bol a été cassé par mon petit frère. \\ \hline
连词 : 不但…而且… & 他 \textbf{不但} 孝顺 \textbf{而且} 很能干。 (Tā \textbf{búdàn} xiàoshùn \textbf{érqiě} hěn nénggàn.) & Il est non seulement filial mais aussi très capable. \\ \hline
连词 : 因为…所以… & \textbf{因为} 爱家 \textbf{所以} 我回家。 (\textbf{Yīnwèi} ài jiā \textbf{suǒyǐ} wǒ huí jiā.) & Parce que j'aime la famille, je rentre à la maison. \\ \hline
语气词 : 吧 (suggestion) / 呢 (relance) & 我们一起做饭 \textbf{吧}？ (Wǒmen yìqǐ zuòfàn \textbf{ba}?) & On cuisine ensemble, d'accord ? \\ \hline
\end{tabularx}
\end{center}

\courssub{Méthodes express — Exposé \& Débat oral}

\begin{itemize}
    \item \textbf{Préparer un exposé (3 min) :} 1) Accroche (\zh{大家好，今天我谈谈……}) 2) 2--3 arguments avec \zh{首先、其次、最后} 3) Exemple personnel (Cameroun/Chine) 4) Conclusion (\zh{总之，家庭很重要}).
    \item \textbf{Participer au débat :} Écouter $\rightarrow$ Noter mots-clés $\rightarrow$ Réagir : \zh{我赞同……因为……} / \zh{我不太同意，我觉得……} $\rightarrow$ Relancer : \zh{你怎么看？} / \zh{还有什么想法？}.
    \item \textbf{Tons prioritaires :} Sandhi 3e ton (\zh{xiǎojiě} $\rightarrow$ \zh{xiáo jiě}), \zh{一} (yī/yí/yì), \zh{不} (bú/bù), neutralisation (\zh{ma, ba, ne, de}).
\end{itemize}



\begin{aretenir}[Checklist avant le Bac — Leçon 3]
$\square$ Je connais le vocabulaire de la famille (父母, 兄弟姐妹, 长辈, 晚辈) et de la solidarité (孝顺, 支持, 团聚, 家务).\\
$\square$ Je sais former une phrase de comparaison avec \textbf{比} (A 比 B Adj) sans ajouter \textbf{很}.\\
$\square$ Je sais utiliser \textbf{把} pour insister sur le résultat d'une action sur un objet connu (做家务 $\rightarrow$ 把家务做完).\\
$\square$ Je reconnais la structure passive \textbf{被} et je sais l'éviter à l'oral si l'agent est inconnu.\\
$\square$ J'utilise correctement les connecteurs \textbf{不但…而且…} et \textbf{因为…所以…} pour structurer mon exposé.\\
$\square$ Je prononce correctement les changements de ton du 3e ton (deux 3e tons $\rightarrow$ 2e + 3e) et les tons de \textbf{一} / \textbf{不}.\\
$\square$ Je maîtrise 5 formules pour donner mon avis (\zh{我认为}, \zh{我觉得}, \zh{依我看来}) et 3 pour relancer (\zh{对吧？}, \zh{你呢？}, \zh{怎么样？}).\\
$\square$ Je suis capable de jouer un jeu de rôle : « Aider un frère/sœur aux devoirs » ou « Préparer la fête du Printemps / Ngondo ».\\
$\square$ Je peux citer une similitude et une différence entre la famille camerounaise et chinoise (ex: respect aînés, repas commun, fête).\\
$\square$ J'écris correctement les caractères clés : \zh{孝} (ordre traits: 孝), \zh{顺}, \zh{团}, \zh{聚}, \zh{家}, \zh{务}.\\
$\square$ Je traduis sans calque : « Je m'occupe de mes parents » $\rightarrow$ \zh{我照顾父母} (pas \zh{我管父母}).\\
$\square$ Je gère mon temps de parole (exposé 2-3 min, débat : interventions 30-45s).
\end{aretenir}

\findecours

\end{document}
